
\Data\ID{1003206001}
\Updated{2021-02-15 07:02:54}
\Level{A}
\SubArea{Quadratic Functions}
\LANGen{
\Question{
We are given three quadratic functions: 
\[ \begin{aligned}
f_1(x)&=-x^2-2, \\
f_2(x)&=-x^2-2x-4, \\
f_3(x)&=x^2+2.
\end{aligned} \]
Which of the given functions are increasing on the interval \( (-2;0) \)?}
{only \( f_1 \)}{only \( f_2 \)}{\( f_1 \) and \( f_2 \)}{all three given functions}{}{}}
\LANGpl{
\Question{
Dane są trzy funkcje kwadratowe:
\[ \begin{aligned}
f_1(x)&=-x^2-2, \\
f_2(x)&=-x^2-2x-4, \\
f_3(x)&=x^2+2.
\end{aligned} \]
Które z nich są rosnące w przedziale \( (-2;0) \)?}
{tylko \( f_1 \)}{tylko \( f_2 \)}{\( f_1 \) i \( f_2 \)}{wszystkie trzy funkcje}{}{}}
\LANGcs{
\Question{
Jsou dány tři kvadratické funkce: 
\[ \begin{aligned}
f_1(x)&=-x^2-2, \\
f_2(x)&=-x^2-2x-4, \\
f_3(x)&=x^2+2.
\end{aligned} \]
Která, nebo které z daných funkcí jsou v intervalu \( (-2;0) \) rostoucí?}
{pouze \( f_1 \)}{pouze \( f_2 \)}{\( f_1 \) a \( f_2 \)}{všechny tři zadané funkce}{}{}}
\LANGsk{
\Question{
Dané sú tri kvadratické funkcie: 
\[ \begin{aligned}
f_1(x)&=-x^2-2, \\
f_2(x)&=-x^2-2x-4, \\
f_3(x)&=x^2+2 .
\end{aligned} \]
Ktoré z daných funkcií sú rastúce na intervale \( (-2;0) \)?}
{len funkcia \( f_1 \)}{len funkcia \( f_2 \)}{funkcie \( f_1 \) a \( f_2 \)}{všetky tri dané funkcie}{}{}}
\LANGes{
\Question{
Tenemos tres funciones cuadráticas:
\[ \begin{aligned}
f_1(x)&=-x^2-2, \\
f_2(x)&=-x^2-2x-4, \\
f_3(x)&=x^2+2.
\end{aligned} \]
¿Cuál de las funciones dadas es creciente en el intervalo \( (-2;0) \)?}
{solamente \( f_1 \)}{solamente \( f_2 \)}{\( f_1 \) y \( f_2 \)}{todas las tres funciones}{}{}}
\End

\Data\ID{1003206201}
\Updated{2021-02-15 07:02:54}
\Level{A}
\SubArea{Quadratic Functions}
\LANGen{
\Question{
Given \( f(x)=2x^2-6x+8 \), find all input value(s) of \( f \) such that the output value of \( f \) is \( 5.5 \).}
{\( x_1=\frac52 \), \( x_2=\frac12 \)}{\( x=35.5 \)}{\( x_1=13 \), \( x_2=11 \)}{\( x_1=-\frac52 \), \( x_2=-\frac12 \)}{}{}}
\LANGpl{
\Question{
Dana jest funkcja \( f(x)=2x^2-6x+8 \), wyznacz wszystkie wartości wejściowe \( f \) takie, dla których wartości wyjściowe  \( f \) wynoszą \( 5{,}5 \).}
{\( x_1=\frac52 \), \( x_2=\frac12 \)}{\( x=35{,}5 \)}{\( x_1=13 \), \( x_2=11 \)}{\( x_1=-\frac52 \), \( x_2=-\frac12 \)}{}{}}
\LANGcs{
\Question{
Je dána kvadratická funkce \( f(x)=2x^2-6x+8 \). Pro která \( x \) bude hodnota funkce \( 5{,}5 \)?}
{\( x_1=\frac52 \), \( x_2=\frac12 \)}{\( x=35{,}5 \)}{\( x_1=13 \), \( x_2=11 \)}{\( x_1=-\frac52 \), \( x_2=-\frac12 \)}{}{}}
\LANGsk{
\Question{
Daná je kvadratická funkcia \( f(x)=2x^2-6x+8 \). Pre aké \( x \) bude funkcia \( f \) nadobúdať hodnotu \( 5{,}5 \)?}
{\( x_1=\frac52 \), \( x_2=\frac12 \)}{\( x=35{,}5 \)}{\( x_1=13 \), \( x_2=11 \)}{\( x_1=-\frac52 \), \( x_2=-\frac12 \)}{}{}}
\LANGes{
\Question{
Dada \( f(x)=2x^2-6x+8 \), encuentra todos los valores dela variable independiente \(x \) para los cuales la variable independiente  \( y \) es \( 5.5 \).}
{\( x_1=\frac52 \), \( x_2=\frac12 \)}{\( x=35.5 \)}{\( x_1=13 \), \( x_2=11 \)}{\( x_1=-\frac52 \), \( x_2=-\frac12 \)}{}{}}
\End

\Data\ID{1003206202}
\Updated{2021-02-15 07:02:36}
\Level{A}
\SubArea{Quadratic Functions}
\LANGen{
\Question{
Given \( f(x)=-\frac12x^2+x+\frac32 \), find all input values of \( f \) such that the output values of \( f \) are  positive.}
{\( x\in(-1;3) \)}{\( x\in(-\infty;-1)\cup(3;+\infty) \)}{\(  x\in(-3;1) \)}{\( x\in(-\infty;-3)\cup(1;+\infty) \)}{}{}}
\LANGpl{
\Question{
Dane jest \( f(x)=-\frac12x^2+x+\frac32 \), wyznacz wszystkie wejściowe wartości  \( f \) dla których wartości wyjściowe  \( f \) są dodatnie.}
{\( x\in(-1;3) \)}{\( x\in(-\infty;-1)\cup(3;+\infty) \)}{\(  x\in(-3;1) \)}{\( x\in(-\infty;-3)\cup(1;+\infty) \)}{}{}}
\LANGcs{
\Question{
Je dána kvadratická funkce \( f(x)=-\frac12x^2+x+\frac32 \). Pro která \( x \) nabývá funkce \( f \) kladných hodnot?}
{\( x\in(-1;3) \)}{\( x\in(-\infty;-1)\cup(3;+\infty) \)}{\(  x\in(-3;1) \)}{\( x\in(-\infty;-3)\cup(1;+\infty) \)}{}{}}
\LANGsk{
\Question{
Daná je kvadratická funkcia \( f(x)=-\frac12x^2+x+\frac32 \). Pre ktoré \( x \)  bude funkcia \( f \) nadobúdať kladné hodnoty?}
{\( x\in(-1;3) \)}{\( x\in(-\infty;-1)\cup(3;+\infty) \)}{\(  x\in(-3;1) \)}{\( x\in(-\infty;-3)\cup(1;+\infty) \)}{}{}}
\LANGes{
\Question{
Dada \( f(x)=-\frac12x^2+x+\frac32 \), encuentra todos los valores dela variable independiente \(x \) para los cuales la variable independiente  \( y \) es espositiva.}
{\( x\in(-1;3) \)}{\( x\in(-\infty;-1)\cup(3;+\infty) \)}{\(  x\in(-3;1) \)}{\( x\in(-\infty;-3)\cup(1;+\infty) \)}{}{}}
\End

\Data\ID{1103083101}
\Updated{2021-04-07 06:04:18}
\Level{A}
\SubArea{Quadratic Functions}
\IMG{1103083101.pdf}
\LANGen{
\Question{
The quadratic  function \( f \) is graphed in the picture.  Which of the following is its equation?}
{\( f(x)=-(x-1)^2+2 \)}{\( f(x)=(x-1)^2+2 \)}{\( f(x)=(x+1)^2+2 \)}{\( f(x)=-(x+1)^2+2 \)}{}{}}
\LANGpl{
\Question{
Funkcję kwadratową \( f \) przedstawiono za pomocą wykresu poniżej. Które z podanych równań jest jej wzorem?}
{\( f(x)=-(x-1)^2+2 \)}{\( f(x)=(x-1)^2+2 \)}{\( f(x)=(x+1)^2+2 \)}{\( f(x)=-(x+1)^2+2 \)}{}{}}
\LANGcs{
\Question{
Na obrázku je graf kvadratické funkce  \( f \). Z nabídnutých možností vyberte její předpis.}
{\( f(x)=-(x-1)^2+2 \)}{\( f(x)=(x-1)^2+2 \)}{\( f(x)=(x+1)^2+2 \)}{\( f(x)=-(x+1)^2+2 \)}{}{}}
\LANGsk{
\Question{
Na obrázku je graf kvadratickej funkcie \( f \). Aký predpis zodpovedá funkcii \( f \)?}
{\( f(x)=-(x-1)^2+2 \)}{\( f(x)=(x-1)^2+2 \)}{\( f(x)=(x+1)^2+2 \)}{\( f(x)=-(x+1)^2+2 \)}{}{}}
\LANGes{
\Question{
En el dibujo aparece la gráfica de una función cuadrática \( f \). ¿Cuál es su expression analítica?}
{\( f(x)=-(x-1)^2+2 \)}{\( f(x)=(x-1)^2+2 \)}{\( f(x)=(x+1)^2+2 \)}{\( f(x)=-(x+1)^2+2 \)}{}{}}
\End

\Data\ID{1103083102}
\Updated{2021-04-07 06:04:51}
\Level{A}
\SubArea{Quadratic Functions}
\IMG{1103083102.pdf}
\LANGen{
\Question{
The quadratic  function \( f \) is graphed in the picture. Which of the following is its equation?}
{\( f(x)=\frac13(x-3)^2+1 \)}{\( f(x)=4\cdot(x-3)^2+1 \)}{\( f(x)=(x+3)^2+4 \)}{\( f(x)=4\cdot(x+1)^2+3 \)}{}{}}
\LANGpl{
\Question{
Funkcję kwadratową \( f \) przedstawiono za pomocą wykresu poniżej. Które z podanych równań jest jej wzorem?}
{\( f(x)=\frac13(x-3)^2+1 \)}{\( f(x)=4\cdot(x-3)^2+1 \)}{\( f(x)=(x+3)^2+4 \)}{\( f(x)=4\cdot(x+1)^2+3 \)}{}{}}
\LANGcs{
\Question{
Na obrázku je graf kvadratické funkce  \( f \). Z nabídnutých možností vyberte její předpis.}
{\( f(x)=\frac13(x-3)^2+1 \)}{\( f(x)=4\cdot(x-3)^2+1 \)}{\( f(x)=(x+3)^2+4 \)}{\( f(x)=4\cdot(x+1)^2+3 \)}{}{}}
\LANGsk{
\Question{
Na obrázku je graf kvadratickej funkcie \( f \).  Aký predpis zodpovedá funkcii \( f \)?}
{\( f(x)=\frac13(x-3)^2+1 \)}{\( f(x)=4\cdot(x-3)^2+1 \)}{\( f(x)=(x+3)^2+4 \)}{\( f(x)=4\cdot(x+1)^2+3 \)}{}{}}
\LANGes{
\Question{
La gráfica de la función \( f \) está representada en el siguiente dibujo. ¿Cuál es su expresión analítica?}
{\( f(x)=\frac13(x-3)^2+1 \)}{\( f(x)=4\cdot(x-3)^2+1 \)}{\( f(x)=(x+3)^2+4 \)}{\( f(x)=4\cdot(x+1)^2+3 \)}{}{}}
\End

\Data\ID{1103083103}
\Updated{2021-04-07 07:04:15}
\Level{A}
\SubArea{Quadratic Functions}
\IMG{1103083103.pdf}
\LANGen{
\Question{
The quadratic function \( f \) is graphed in the picture. Which of the following is its equation?}
{\( f(x)=-x^2-2x+1  \)}{\( f(x)=x^2+2x+1 \)}{\( f(x)=-x^2+2x+1 \)}{\( f(x)=-x^2-4x-5 \)}{}{}}
\LANGpl{
\Question{
Funkcję kwadratową \( f \) przedstawiono za pomocą wykresu poniżej. Które z podanych równań jest jej wzorem?}
{\( f(x)=-x^2-2x+1  \)}{\( f(x)=x^2+2x+1 \)}{\( f(x)=-x^2+2x+1 \)}{\( f(x)=-x^2-4x-5 \)}{}{}}
\LANGcs{
\Question{
Na obrázku je graf kvadratické funkce  \( f \). Z nabídnutých možností vyberte její předpis.}
{\( f(x)=-x^2-2x+1  \)}{\( f(x)=x^2+2x+1 \)}{\( f(x)=-x^2+2x+1 \)}{\( f(x)=-x^2-4x-5 \)}{}{}}
\LANGsk{
\Question{
Na obrázku je graf kvadratickej funkcie \( f \). Aký predpis zodpovedá funkcii \( f \)?}
{\( f(x)=-x^2-2x+1  \)}{\( f(x)=x^2+2x+1 \)}{\( f(x)=-x^2+2x+1 \)}{\( f(x)=-x^2-4x-5 \)}{}{}}
\LANGes{
\Question{
En el dibujo aparece la gráfica de una función cuadrática \( f \). ¿Cuál es su expresión analítica?}
{\( f(x)=-x^2-2x+1  \)}{\( f(x)=x^2+2x+1 \)}{\( f(x)=-x^2+2x+1 \)}{\( f(x)=-x^2-4x-5 \)}{}{}}
\End

\Data\ID{1103083104}
\Updated{2021-04-07 07:04:54}
\Level{A}
\SubArea{Quadratic Functions}
\IMG{1103083104.pdf}
\LANGen{
\Question{
The quadratic function \( f \) is graphed in the picture. Which of the following is its equation?}
{\( f(x)=\frac23x^2-4x+4 \)}{\( f(x)=x^2-6x+7 \)}{\( f(x)=\frac13x^2-2x+1 \)}{\( f(x)=x^2-2x+4 \)}{}{}}
\LANGpl{
\Question{
Funkcję kwadratową  \( f \) przedstawiono poniżej za pomocą wykresu. Które z poniższych równań jest jej wzorem?}
{\( f(x)=\frac23x^2-4x+4 \)}{\( f(x)=x^2-6x+7 \)}{\( f(x)=\frac13x^2-2x+1 \)}{\( f(x)=x^2-2x+4 \)}{}{}}
\LANGcs{
\Question{
Na obrázku je graf kvadratické funkce  \( f \). Z nabídnutých možností vyberte její předpis.}
{\( f(x)=\frac23x^2-4x+4 \)}{\( f(x)=x^2-6x+7 \)}{\( f(x)=\frac13x^2-2x+1 \)}{\( f(x)=x^2-2x+4 \)}{}{}}
\LANGsk{
\Question{
Na obrázku je graf kvadratickej funkcie \( f \). Aký predpis zodpovedá funkcii \( f \)?}
{\( f(x)=\frac23x^2-4x+4 \)}{\( f(x)=x^2-6x+7 \)}{\( f(x)=\frac13x^2-2x+1 \)}{\( f(x)=x^2-2x+4 \)}{}{}}
\LANGes{
\Question{
En el dibujo aparece la gráfica de una función cuadrática \( f \). ¿Cuál es su expresión analítica?}
{\( f(x)=\frac23x^2-4x+4 \)}{\( f(x)=x^2-6x+7 \)}{\( f(x)=\frac13x^2-2x+1 \)}{\( f(x)=x^2-2x+4 \)}{}{}}
\End

\Data\ID{1103083105}
\Updated{2021-04-07 07:04:23}
\Level{A}
\SubArea{Quadratic Functions}
\IMG{1103083105.pdf}
\LANGen{
\Question{
The quadratic function \( f \) is graphed in the picture. Which of the following is its equation?}
{\( f(x)=-2x^2-4x-5 \)}{\( f(x)=-2x^2+12x-19 \)}{\( f(x)=-x^2-2x-4 \)}{\( f(x)=x^2-2x-5 \)}{}{}}
\LANGpl{
\Question{
Funkcję kwadratową \( f \) przedstawiono za pomocą wykresu poniżej. Które z poniższych równań jest jej wzorem?}
{\( f(x)=-2x^2-4x-5 \)}{\( f(x)=-2x^2+12x-19 \)}{\( f(x)=-x^2-2x-4 \)}{\( f(x)=x^2-2x-5 \)}{}{}}
\LANGcs{
\Question{
Na obrázku je graf kvadratické funkce  \( f \). Z nabídnutých možností vyberte její předpis.}
{\( f(x)=-2x^2-4x-5 \)}{\( f(x)=-2x^2+12x-19 \)}{\( f(x)=-x^2-2x-4 \)}{\( f(x)=x^2-2x-5 \)}{}{}}
\LANGsk{
\Question{
Na obrázku je graf kvadratickej funkcie \( f \). Aký predpis zodpovedá funkcii \( f \)?}
{\( f(x)=-2x^2-4x-5 \)}{\( f(x)=-2x^2+12x-19 \)}{\( f(x)=-x^2-2x-4 \)}{\( f(x)=x^2-2x-5 \)}{}{}}
\LANGes{
\Question{
En el dibujo aparece la gráfica de una función cuadrática \( f \). ¿Cuál es su expresión analítica?}
{\( f(x)=-2x^2-4x-5 \)}{\( f(x)=-2x^2+12x-19 \)}{\( f(x)=-x^2-2x-4 \)}{\( f(x)=x^2-2x-5 \)}{}{}}
\End

\Data\ID{1103083106}
\Updated{2021-04-07 07:04:40}
\Level{A}
\SubArea{Quadratic Functions}
\IMG{1103083106.pdf}
\LANGen{
\Question{
The quadratic function \( f \) is graphed in the picture. Which of the following is its equation?}
{\( f(x)=x^2+6x+9 \)}{\( f(x)=x^2-6x+9 \)}{\( f(x)=x^2-3 \)}{\( f(x)=x^2+3 \)}{}{}}
\LANGpl{
\Question{
Funkcję kwadratową  \( f \) przedstawiono za pomocą wykresu.  Które z poniższych jest jej równaniem?}
{\( f(x)=x^2+6x+9 \)}{\( f(x)=x^2-6x+9 \)}{\( f(x)=x^2-3 \)}{\( f(x)=x^2+3 \)}{}{}}
\LANGcs{
\Question{
Na obrázku je graf kvadratické funkce  \( f \). Z nabídnutých možností vyberte její předpis.}
{\( f(x)=x^2+6x+9 \)}{\( f(x)=x^2-6x+9 \)}{\( f(x)=x^2-3 \)}{\( f(x)=x^2+3 \)}{}{}}
\LANGsk{
\Question{
Na obrázku je graf kvadratickej funkcie \( f \). Aký predpis zodpovedá funkcii \( f \)?}
{\( f(x)=x^2+6x+9 \)}{\( f(x)=x^2-6x+9 \)}{\( f(x)=x^2-3 \)}{\( f(x)=x^2+3 \)}{}{}}
\LANGes{
\Question{
En el dibujo aparece la gráfica de una función cuadrática \( f \). ¿Cuál es su expresión analítica?}
{\( f(x)=x^2+6x+9 \)}{\( f(x)=x^2-6x+9 \)}{\( f(x)=x^2-3 \)}{\( f(x)=x^2+3 \)}{}{}}
\End

\Data\ID{1103083111}
\Updated{2020-10-09 12:10:40}
\Level{A}
\SubArea{Quadratic Functions}
\LANGen{
\Question{
Find the graph of the quadratic function \( f(x)=2x^2+12x+16 \).}
{}{}{}{}{}{}}
\LANGpl{
\Question{
Wyznacz wykres funkcji kwadratowej \( f(x)=2x^2+12x+16 \).}
{}{}{}{}{}{}}
\LANGcs{
\Question{
Je dána kvadratická funkce \( f(x)=2x^2+12x+16 \). Vyberte obrázek, na kterém je graf této funkce.}
{}{}{}{}{}{}}
\LANGsk{
\Question{
Ktorý z nasledujúcich grafov je graf kvadratickej funkcie \( f(x)=2x^2+12x+16 \)?}
{}{}{}{}{}{}}
\LANGes{
\Question{
Encuentra la gráfica de la función cuadrática \( f(x)=2x^2+12x+16 \).}
{}{}{}{}{}{}}

\IMGanswer{1103083111a.pdf}
\IMGanswer{1103083111b.pdf}
\IMGanswer{1103083111c.pdf}
\IMGanswer{1103083111d.pdf}\End

\Data\ID{1103083112}
\Updated{2020-10-09 12:10:53}
\Level{A}
\SubArea{Quadratic Functions}
\LANGen{
\Question{
Find the graph of the quadratic function \( f(x)=\frac15x^2-2x+6 \).}
{}{}{}{}{}{}}
\LANGpl{
\Question{
Wyznacz wykres funkcji kwadratowej \( f(x)=\frac15x^2-2x+6 \).}
{}{}{}{}{}{}}
\LANGcs{
\Question{
Je dána kvadratická funkce \( f(x)=\frac15x^2-2x+6 \). Vyberte obrázek, na kterém je graf této funkce.}
{}{}{}{}{}{}}
\LANGsk{
\Question{
Ktorý z nasledujúcich grafov je graf kvadratickej funkcie \( f(x)=\frac15x^2-2x+6 \)?}
{}{}{}{}{}{}}
\LANGes{
\Question{
Encuentra la gráfica de la función cuadrática \( f(x)=\frac15x^2-2x+6 \).}
{}{}{}{}{}{}}

\IMGanswer{1103083112a.pdf}
\IMGanswer{1103083112b.pdf}
\IMGanswer{1103083112c.pdf}
\IMGanswer{1103083112d.pdf}\End

\Data\ID{1103120004}
\Updated{2022-04-23 05:04:54}
\Level{A}
\SubArea{Quadratic Functions}
\IMG{1103120004.pdf}
\LANGen{
\Question{
Let \( f(x)=x^2 \). Given the graph of the function \( f \) and the graph of a function \( g \) which was obtained as a vertical shift  of the graph of \( f \) (see the picture), choose the function \( g \).}
{\( g(x) = x^2-3 \)}{\( g(x) = (x+3)^2 \)}{\( g(x) = x^2+3 \)}{\( g(x) = (x-3)^2 \)}{}{}}
\LANGpl{
\Question{
Niech  \( f(x)=x^2 \). Dany jest wykres funkcji \( f \) i wykres funkcji  \( g \), który uzyskano poprzez pionowe przesunięcie wykresu funkcji  \( f \) (rysunek), wybierz funkcję  \( g \).}
{\( g(x) = x^2-3 \)}{\( g(x) = (x+3)^2 \)}{\( g(x) = x^2+3 \)}{\( g(x) = (x-3)^2 \)}{}{}}
\LANGcs{
\Question{
Na obrázku je nakreslen graf funkce \( f(x)=x^2 \) a graf funkce \( g \), který vznikl posunem grafu funkce \( f \). Vyberte předpis funkce \( g \).}
{\( g(x) = x^2-3 \)}{\( g(x) = (x+3)^2 \)}{\( g(x) = x^2+3 \)}{\( g(x) = (x-3)^2 \)}{}{}}
\LANGsk{
\Question{
Na obrázku je nakreslený graf funkcie \( f(x)=x^2 \) a graf funkcie \( g \), ktorý vznikol posunutím grafu funkcie \( f \). Vyberte predpis funkcie \( g \).}
{\( g(x) = x^2-3 \)}{\( g(x) = (x+3)^2 \)}{\( g(x) = x^2+3 \)}{\( g(x) = (x-3)^2 \)}{}{}}
\LANGes{
\Question{
Sea \( f(x)=x^2 \). Dada la gráfica de la función \( f \) y la gráfica de la función \( g \) que fue obtenida por un movimiento vertical de la gráfica de \( f \) (en el dibujo), elije la función \( g \).}
{\( g(x) = x^2-3 \)}{\( g(x) = (x+3)^2 \)}{\( g(x) = x^2+3 \)}{\( g(x) = (x-3)^2 \)}{}{}}
\End

\Data\ID{2000004301}
\Updated{2022-02-20 08:02:18}
\Level{A}
\SubArea{Quadratic Functions}
\LANGen{
\Question{
Find the intervals of monotonicity of the quadratic function \( f(x)=4-3x^2\).}
{The function is increasing on \( (-\infty; 0 ]\) and decreasing on \( [ 0 ; +\infty)\).}{The function is increasing on \( (-\infty; 4 ]\) and decreasing on \( [ 4 ; +\infty)\).}{The function is decreasing on \( (-\infty; 0 ]\) and increasing on \( [ 0 ; +\infty)\).}{The function is decreasing on \( (-\infty; 4 ]\) and increasing on \( [ 4 ; +\infty)\).}{}{}}
\LANGpl{
\Question{
Znajdź przedziały monotoniczności funkcji kwadratowej  \( f(x)=4-3x^2\).}
{Funkcja rośnie w przedziale \( (-\infty; 0 \rangle\) i maleje w przedziale \( \langle 0 ; +\infty)\).}{Funkcja rośnie w przedziale \( (-\infty; 4 \rangle\) i maleje w przedziale \( \langle 4 ; +\infty)\).}{Funkcja maleje w przedziale \( (-\infty; 0 \rangle\) i rośnie w przedziale \( \langle 0 ; +\infty)\).}{Funkcja maleje w przedziale \( (-\infty; 4 \rangle\) i rośnie w przedziale \( \langle 4 ; +\infty)\).}{}{}}
\LANGcs{
\Question{
Určete intervaly monotónie kvadratické funkce \( f: y =4-3x^2\).}
{Funkce roste na intervalu  \( (-\infty; 0 \rangle\) a klesá na intervalu  \( \langle 0 ; +\infty)\).}{Funkce roste na intervalu  \( (-\infty; 4 \rangle\) a klesá na intervalu  \( \langle 4 ; +\infty)\).}{Funkce klesá na intervalu  \( (-\infty; 0 \rangle\) a roste na intervalu \( \langle 0 ; +\infty)\).}{Funkce klesá na intervalu  \( (-\infty; 4 \rangle\) a roste na intervalu  \( \langle 4 ; +\infty)\).}{}{}}
\LANGsk{
\Question{
Určte intervaly monotónnosti kvadratickej funkcie \( f: y =4-3x^2\).}
{Funkcia rastie na intervale  \( (-\infty; 0 \rangle\) a klesá na intervale  \( \langle 0 ; +\infty)\).}{Funkcia rastie na intervale  \( (-\infty; 4 \rangle\) a klesá na intervale  \( \langle 4 ; +\infty)\).}{Funkcia klesá na intervale  \( (-\infty; 0 \rangle\) a rastie na intervale \( \langle 0 ; +\infty)\).}{Funkcia klesá na intervale  \( (-\infty; 4 \rangle\) a rastie na intervale  \( \langle 4 ; +\infty)\).}{}{}}
\LANGes{
\Question{
Averigua los intervalos de monotonía de la función cuadrática  \( f: y =4-3x^2\).}
{La función crece en el intervalo \( (-\infty; 0 ]\) y decrece en el intervalo \( [ 0 ; +\infty)\).}{La función crece en el intervalo \( (-\infty; 4 ]\) y decrece en el intervalo \( [ 4 ; +\infty)\).}{La función decrece en el intervalo  \( (-\infty; 0 ]\) y crece en el intervalo \( [ 0 ; +\infty)\).}{La función decrece en el intervalo\( (-\infty; 4 ]\) y crece en el intervalo \( [ 4 ; +\infty)\).}{}{}}
\End

\Data\ID{2010012202}
\Updated{2022-11-02 09:11:23}
\Level{A}
\SubArea{Quadratic Functions}
\LANGen{
\Question{
Find all the values of the real parameter \( a \) such that \( f(x)=ax^2+2 \) is increasing on \( (0;\infty) \).}
{\( a\in(0;+\infty) \)}{\( a\in(-\infty;0) \)}{\( a\in [ 2;+\infty) \)}{\( a\in(-\infty;2 ] \)}{}{}}
\LANGpl{
\Question{
Znajdź wszystkie wartości rzeczywistego parametru \( a \), dla których funkcja \( f(x)=ax^2+2 \) rośnie w przedziale \( (0;\infty) \).}
{\( a\in(0;+\infty) \)}{\( a\in(-\infty;0) \)}{\( a\in \langle 2;+\infty) \)}{\( a\in(-\infty;2 \rangle \)}{}{}}
\LANGcs{
\Question{
Najděte všechny hodnoty reálného parametru  \( a \), pro které je funkce  \( f(x)=ax^2+2 \) rostoucí na intervalu  \( (0;\infty) \).}
{\( a\in(0;+\infty) \)}{\( a\in(-\infty;0) \)}{\( a\in \langle 2;+\infty) \)}{\( a\in(-\infty;2 \rangle \)}{}{}}
\LANGsk{
\Question{
Nájdite všetky hodnoty reálneho parametra  \( a \), pre ktoré je funkcia  \( f(x)=ax^2+2 \) rastúca na intervale  \( (0;\infty) \).}
{\( a\in(0;+\infty) \)}{\( a\in(-\infty;0) \)}{\( a\in \langle 2;+\infty) \)}{\( a\in(-\infty;2 \rangle \)}{}{}}
\LANGes{
\Question{
Halla todos los valores del parámetro real \( a \) tales que \( f(x)=ax^2+2 \) sea creciente en \( (0;\infty) \).}
{\( a\in(0;+\infty) \)}{\( a\in(-\infty;0) \)}{\( a\in [ 2;+\infty) \)}{\( a\in(-\infty;2 ] \)}{}{}}
\End

\Data\ID{2010012301}
\Updated{2022-11-02 09:11:23}
\Level{A}
\SubArea{Quadratic Functions}
\LANGen{
\Question{
Find the \(x\)-intercepts
of the function \(f(x)= 2x^{2} + 2x - 12\).}
{\([-3;0]\) and
\([2;0]\)}{\([0;-12]\) and
\([2;0]\)}{\([-3;2]\) and
\([-3;-2]\)}{The function \(f\) does
not have \(x\)-intercepts.}{}{}}
\LANGpl{
\Question{
Wyznacz punkty przecięcia wykresu funkcji \(f(x)= 2x^{2} + 2x - 12\) z osią \(x\).}
{\([-3;0]\) i
\([2;0]\)}{\([0;-12]\) i
\([2;0]\)}{\([-3;2]\) i
\([-3;-2]\)}{Funkcja \(f\) nie przecina osi \(x\).}{}{}}
\LANGcs{
\Question{
Je dána funkce \(f\):  \(y= 2x^{2} + 2x - 12\). Určete průsečíky grafu funkce \(f\) s osou \(x\).}
{\([-3;0]\) a \([2;0]\)}{\([0;-12]\) a
\([2;0]\)}{\([-3;2]\) a
\([-3;-2]\)}{Graf funkce \(f\) neprotíná osu \(x\).}{}{}}
\LANGsk{
\Question{
Je daná funkcia \(f\):  \(y= 2x^{2} + 2x - 12\). Určte priesečníky grafu funkcie \(f\) s osou \(x\).}
{\([-3;0]\) a
\([2;0]\)}{\([0;-12]\) a
\([2;0]\)}{\([-3;2]\) a
\([-3;-2]\)}{Funkcia \(f\) nemá s osou \(x\)-priesečníky.}{}{}}
\LANGes{
\Question{
Halla los puntos de corte con el eje \(x\)
de la función \(f(x)= 2x^{2} + 2x - 12\).}
{\([-3;0]\) y
\([2;0]\)}{\([0;-12]\) y
\([2;0]\)}{\([-3;2]\) y
\([-3;-2]\)}{La función \(f\) no tiene puntos de corte con el eje \(x\).}{}{}}
\End

\Data\ID{2010012302}
\Updated{2022-11-02 09:11:23}
\Level{A}
\SubArea{Quadratic Functions}
\LANGen{
\Question{
Find the intervals of monotonicity of the quadratic function
\(f(x) = -3x^{2} + 2\).}
{The function is increasing on \( (- \infty ;0 ] \) 
and decreasing on \( [ 0;\infty ) \).}{The function is increasing on \((-\infty;2) \) 
and decreasing on \( ( 2;\infty) \).}{The function is increasing on \(\left(-\infty;\frac23 \right] \) 
and decreasing on \( \left[ \frac23;\infty\right) \).}{The function is decreasing on its domain.}{}{}}
\LANGpl{
\Question{
Wyznacz przedziały monotoniczności funkcji kwadratowej
\(f(x) = -3x^{2} + 2\).}
{Funkcja rośnie w przedziale \( (- \infty ;0 \rangle \) 
i maleje w przedziale \( \langle 0;\infty ) \).}{Funkcja rośnie w przedziale \((-\infty;2) \) 
i maleje w przedziale \( ( 2;\infty) \).}{Funkcja rośnie w przedziale\(\left(-\infty;\frac23 \right\rangle \) 
 i maleje w przedziale \( \left\langle \frac23;\infty\right) \).}{Funkcja maleje w całej swojej dziedzinie.}{}{}}
\LANGcs{
\Question{
Určete intervaly monotonie kvadratické funkce 
\(f: y = -3x^{2} + 2\).}
{Funkce roste na intervalu  \( (- \infty ;0 \rangle \) 
a klesá na intervalu \( \langle 0;\infty ) \).}{Funkce roste na intervalu  \((-\infty;2) \) 
a klesá na intervalu \( ( 2;\infty) \).}{Funkce roste na intervalu  \(\left(-\infty;\frac23 \right\rangle \) 
a klesá na intervalu \( \left\langle \frac23;\infty\right) \).}{Funkce je klesající na celém \(D(f)\).}{}{}}
\LANGsk{
\Question{
Určte intervaly monotónnosti kvadratickej funkcie 
\(f: y = -3x^{2} + 2\).}
{Funkcia je rastúca na \( (- \infty ;0 \rangle \) 
a klesajúca na \( \langle 0;\infty ) \).}{Funkcia je rastúca na \((-\infty;2) \) 
a klesajúca na \( ( 2;\infty) \).}{Funkcia je rastúca na \(\left(-\infty;\frac23 \right\rangle \) 
a klesajúca na \( \left\langle \frac23;\infty\right) \).}{Funkcia je klesajúca na celom definičnom obore.}{}{}}
\LANGes{
\Question{
Halla los intervalos de monotonía de la función cuadrática
\(f(x) = -3x^{2} + 2\).}
{La función es creciente en \( (- \infty ;0 ] \) 
y decreciente en \( [ 0;\infty ) \).}{La función es creciente en \((-\infty;2) \) 
y decreciente en \( ( 2;\infty) \).}{La función es creciente en \(\left(-\infty;\frac23 \right] \) 
y decreciente en \( \left[ \frac23;\infty\right) \).}{La función es decreciente en todo su dominio.}{}{}}
\End

\Data\ID{2010012303}
\Updated{2024-05-31 11:05:45}
\Level{A}
\SubArea{Quadratic Functions}
\LANGen{
\Question{
Find the \(y\)-intercept
of the following function. \[f(x) = 6x^{2} +12x - 7.2\]}
{\([0;-7.2]\)}{\([-7.2;0]\)}{\([6;0]\)}{There is no \(y\)-intercept.}{}{}}
\LANGpl{
\Question{
Wyznacz punkty przecięcia wykresu funkcji \[f(x) = 6x^{2} +12x - 7{,}2\] z osią \(y\).}
{\([0;-7{,}2]\)}{\([-7{,}2;0]\)}{\([6;0]\)}{Funkcja nie przecina osi \(y\).}{}{}}
\LANGcs{
\Question{
Určete průsečíky grafu funkce \[f: y = 6x^{2} +12x - 7{,}2\] s osou \(y\).}
{\([0;-7{,}2]\)}{\([-7{,}2;0]\)}{\([6;0]\)}{Funkce \(f\) osu \(y\) neprotíná.}{}{}}
\LANGsk{
\Question{
Určte priesečníky grafu funkcie \[f: y = 6x^{2} +12x - 7{,}2\] s osou \(y\).}
{\([0;-7{,}2]\)}{\([-7{,}2;0]\)}{\([6;0]\)}{Funkcia \(f\) os \(y\) nepretína.}{}{}}
\LANGes{
\Question{
Halla los puntos de corte con el eje \(y\)
de la siguiente función. \[f(x) = 6x^{2} +12x - 7.2\]}
{\([0;-7.2]\)}{\([-7.2;0]\)}{\([6;0]\)}{No presenta puntos de corte con el eje \(y\).}{}{}}
\End

\Data\ID{2010012304}
\Updated{2022-11-02 09:11:23}
\Level{A}
\SubArea{Quadratic Functions}
\IMG{2010012301-figure0.pdf}
\LANGen{
\Question{
Let \( f(x)=-x^2 \). Given the graph of the function \( f \) and the graph of a function \( g \) which was obtained as a vertical shift  of the graph of \( f \) (see the picture), choose the function \( g \).}
{\( g(x) = -x^2+2 \)}{\( g(x) = (x-2)^2 \)}{\( g(x) = -x^2-2 \)}{\( g(x) = (x+2)^2 \)}{}{}}
\LANGpl{
\Question{
Rysunek przedstawia wykres funkcji \( f(x)=-x^2 \) oraz wykres funkcji \( g \),  który został utworzony przez przesunięcie wykresu funkcji \( f \). Wybierz wzór funkcji \( g \) (patrz rysunek).}
{\( g(x) = -x^2+2 \)}{\( g(x) = (x-2)^2 \)}{\( g(x) = -x^2-2 \)}{\( g(x) = (x+2)^2 \)}{}{}}
\LANGcs{
\Question{
Na obrázku je nakreslen graf funkce \( f(x)=-x^2 \) a graf funkce \(g\), který vznikl posunutím grafu funkce \( f \). Vyberte předpis funkce  \( g \) (viz obrázek).}
{\( g(x) = -x^2+2 \)}{\( g(x) = (x-2)^2 \)}{\( g(x) = -x^2-2 \)}{\( g(x) = (x+2)^2 \)}{}{}}
\LANGsk{
\Question{
Na obrázku je nakreslený graf funkcie \( f(x)=-x^2 \) a graf funkcie \(g\), ktorý vznikol posunom grafu funkcie \( f \). Vyberte predpis funkcie  \( g \), pozri obrázok.}
{\( g(x) = -x^2+2 \)}{\( g(x) = (x-2)^2 \)}{\( g(x) = -x^2-2 \)}{\( g(x) = (x+2)^2 \)}{}{}}
\LANGes{
\Question{
Sea \( f(x)=-x^2 \). Dada la gráfica de la función \( f \) y la de la función \( g \) que obtenemos por un desplazamiento vertical de la gráfica de \( f \) (ver imagen), elige la expresión analítica de la función \( g \).}
{\( g(x) = -x^2+2 \)}{\( g(x) = (x-2)^2 \)}{\( g(x) = -x^2-2 \)}{\( g(x) = (x+2)^2 \)}{}{}}
\End

\Data\ID{9000014801}
\Updated{2020-11-20 10:11:35}
\Level{A}
\SubArea{Quadratic Functions}
\LANGen{
\Question{
Identify a point which is on the graph of the function
\(f(x) = 3x^{2} + 3x - 2\).}
{\(B = [2;16]\)}{\(A = [0;3]\)}{\(C = [-1;0]\)}{\(D = [5;-8]\)}{}{}}
\LANGpl{
\Question{
Określ, który punkt znajduje się na wykresie funkcji
\(f\colon y = 3x^{2} + 3x - 2\).}
{\(B = [2;16]\)}{\(A = [0;3]\)}{\(C = [-1;0]\)}{\(D = [5;-8]\)}{}{}}
\LANGcs{
\Question{
Který z následujících bodů leží na grafu funkce
\(f\colon y = 3x^{2} + 3x - 2\)?}
{\(B = [2;16]\)}{\(A = [0;3]\)}{\(C = [-1;0]\)}{\(D = [5;-8]\)}{}{}}
\LANGsk{
\Question{
Ktorý z následujúcich bodov leží na grafe funkcie
\(f\colon y = 3x^{2} + 3x - 2\)?}
{\(B = [2;16]\)}{\(A = [0;3]\)}{\(C = [-1;0]\)}{\(D = [5;-8]\)}{}{}}
\LANGes{
\Question{
¿Cuál de los puntos pertenece a la gráfica de la función
\(f(x) = 3x^{2} + 3x - 2\)?}
{\(B = [2;16]\)}{\(A = [0;3]\)}{\(C = [-1;0]\)}{\(D = [5;-8]\)}{}{}}
\End

\Data\ID{9000014802}
\Updated{2020-11-20 10:11:54}
\Level{A}
\SubArea{Quadratic Functions}
\LANGen{
\Question{
Let \(f(x) = -x^{2} + 11x - 2\).
Which of the following statements is true?}
{\(f(-2) = -28\)}{\(f(0) = 2\)}{\(f(3.5) = 12.25\)}{\(f\left (\frac{1}
{2}\right ) = \frac{15} 
 {4} \)}{}{}}
\LANGpl{
\Question{
Jeśli  \(f\colon y = -x^{2} + 11x - 2\).
Które z podanych stwierdzeń jest prawdziwe?}
{\(f(-2) = -28\)}{\(f(0) = 2\)}{\(f(3.5) = 12.25\)}{\(f\left (\frac{1}
{2}\right ) = \frac{15} 
 {4} \)}{}{}}
\LANGcs{
\Question{
Je dána funkce \(f\colon y = -x^{2} + 11x - 2\).
Označte, který výrok platí.}
{\(f(-2) = -28\)}{\(f(0) = 2\)}{\(f(3{,}5) = 12{,}25\)}{\(f\left (\frac{1}
{2}\right ) = \frac{15} 
 {4} \)}{}{}}
\LANGsk{
\Question{
Je daná funkcia \(f\colon y = -x^{2} + 11x - 2\).
Ktoré z nasledujúcich tvrdení o funkcii \(f\)  je pravdivé?}
{\(f(-2) = -28\)}{\(f(0) = 2\)}{\(f(3{,}5) = 12{,}25\)}{\(f\left (\frac{1}
{2}\right ) = \frac{15} 
 {4} \)}{}{}}
\LANGes{
\Question{
Sea \(f(x) = -x^{2} + 11x - 2\).
¿Cuál de las declaraciones es correcta?}
{\(f(-2) = -28\)}{\(f(0) = 2\)}{\(f(3.5) = 12.25\)}{\(f\left (\frac{1}
{2}\right ) = \frac{15} 
 {4} \)}{}{}}
\End

\Data\ID{9000014807}
\Updated{2022-04-23 05:04:54}
\Level{A}
\SubArea{Quadratic Functions}
\LANGen{
\Question{
Find the \(x\)-intercepts
of the function \(f(x)= 3x^{2} + 6x - 9\).}
{\([-3;0]\) and
\([1;0]\)}{\([0;9]\) and
\([1;0]\)}{\([-3;2]\) and
\([-3;-2]\)}{The function \(f\) does
not have \(x\)-intercepts.}{}{}}
\LANGpl{
\Question{
Znajdź punkty przecięcia prostej z osią współrzędnych \(x\)
funkcji \(f\colon y = 3x^{2} + 6x - 9\).}
{\([-3;0]\) i \([1;0]\)}{\([0;9]\) i \([1;0]\)}{\([-3;2]\) i \([-3;-2]\)}{Funkcja \(f\) nie ma punktów przecięcia prostej z osią współrzędnych  \(x\).}{}{}}
\LANGcs{
\Question{
Je dána funkce \(f\colon y = 3x^{2} + 6x - 9\).
Určete průsečíky grafu funkce s osou
\(x\).}
{\([-3;0]\) a
\([1;0]\)}{\([0;9]\) a
\([1;0]\)}{\([-3;2]\) a
\([-3;-2]\)}{Graf funkce \(f\) 
neprotíná osu \(x\).}{}{}}
\LANGsk{
\Question{
Daná je funkcia \(f\colon y = 3x^{2} + 6x - 9\).
Určte priesečníky grafu funkcie s osou
\(x\).}
{\([-3;0]\) a
\([1;0]\)}{\([0;9]\) a
\([1;0]\)}{\([-3;2]\) a
\([-3;-2]\)}{Graf funkcie \(f\) 
nepretína os \(x\).}{}{}}
\LANGes{
\Question{
Halla los puntos de intersección de la función \(f(x)= 3x^{2} + 6x - 9\) con el eje \(x\).}
{\([-3;0]\) y
\([1;0]\)}{\([0;9]\) y
\([1;0]\)}{\([-3;2]\) y
\([-3;-2]\)}{La función \(f\) no interseca el eje \(x\).}{}{}}
\End

\Data\ID{9000014808}
\Updated{2022-04-23 05:04:54}
\Level{A}
\SubArea{Quadratic Functions}
\LANGen{
\Question{
Find the intervals of monotonicity of the quadratic function
\(f(x) = 2x^{2} + 3\).}
{The function is increasing on \(\left [ 0;\infty \right )\) 
and decreasing on \(\left (-\infty ;0\right ] \).}{The function is increasing on \(\left (3;\infty \right )\) 
and decreasing on \(\left (-\infty ;3\right )\).}{The function is increasing on \(\left [ -\frac{3}
{2};\infty \right )\) 
and decreasing on \(\left (-\infty ;-\frac{3}
{2}\right ] \).}{The function is increasing on its domain.}{}{}}
\LANGpl{
\Question{
Określ przedziały monotoniczności funkcji kwadratowej
\(f\colon y = 2x^{2} + 3\).}
{Funkcja jest rosnąca w przedziale \(\left [ 0;\infty \right )\) 
i malejąca w przedziale \(\left (-\infty ;0\right ] \).}{Funkcja jest rosnąca w przedziale \(\left (3;\infty \right )\) 
i malejąca w przedziale \(\left (-\infty ;3\right )\).}{Funkcja jest rosnąca w przedziale \(\left [ -\frac{3}
{2};\infty \right )\) 
i malejąca w przedziale  \(\left (-\infty ;-\frac{3}
{2}\right ] \).}{Funkcja jest rosnąca w swojej dziedzinie.}{}{}}
\LANGcs{
\Question{
Určete intervaly monotonie kvadratické funkce
\(f\colon y = 2x^{2} + 3\).}
{Funkce roste na intervalu \(\left \langle 0;\infty \right )\) 
a klesá na intervalu \(\left (-\infty ;0\right \rangle \).}{Funkce roste na intervalu \(\left (3;\infty \right )\) 
a klesá na intervalu \(\left (-\infty ;3\right )\).}{Funkce roste na intervalu \(\left \langle -\frac{3}
{2};\infty \right )\) 
a klesá na intervalu \(\left (-\infty ;-\frac{3}
{2}\right \rangle \).}{Funkce je rostoucí na celém \(D(f)\).}{}{}}
\LANGsk{
\Question{
Určte intervaly monotónnosti kvadratickej funkcie
\(f\colon y = 2x^{2} + 3\).}
{Funkcia rastie na intervale \(\left \langle 0;\infty \right )\) 
a klesá na intervale \(\left (-\infty ;0\right \rangle \).}{Funkcia rastie na intervale \(\left (3;\infty \right )\) 
a klesá na intervale \(\left (-\infty ;3\right )\).}{Funkcia rastie na intervale \(\left \langle -\frac{3}
{2};\infty \right )\) 
a klesá na intervale \(\left (-\infty ;-\frac{3}
{2}\right \rangle \).}{Funkcia je rastúca na celom \(D(f)\).}{}{}}
\LANGes{
\Question{
Encuentra los intervalos de monotonía de la función
\(f(x) = 2x^{2} + 3\).}
{La función es creciente en \(\left [ 0;\infty \right )\) 
y decreciente en \(\left (-\infty ;0\right ] \).}{La función es creciente en \(\left (3;\infty \right )\) 
y decreciente en \(\left (-\infty ;3\right )\).}{La función es creciente en \(\left [ -\frac{3}
{2};\infty \right )\) 
y decreciente en \(\left (-\infty ;-\frac{3}
{2}\right ] \).}{La función es creciente en todo su Dominio.}{}{}}
\End

\Data\ID{9000014809}
\Updated{2022-04-23 05:04:54}
\Level{A}
\SubArea{Quadratic Functions}
\LANGen{
\Question{
Find the \(y\)-intercept
of the following function: \[f(x) = 10x^{2} - 18x - 6.3\]}
{\([0;-6.3]\)}{\([10;0]\)}{\([0.3;0]\)}{There is no \(y\)-intercept.}{}{}}
\LANGpl{
\Question{
Znajdź punkt przecięcia prostej z osią współrzędnych  \(y\)
funkcji \(f\colon y = 10x^{2} - 18x - 6.3\).}
{\([0;-6.3]\)}{\([10;0]\)}{\([0.3;0]\)}{Nie ma takiego punktu.}{}{}}
\LANGcs{
\Question{
Určete průsečíky grafu funkce \(f\colon y = 10x^{2} - 18x - 6{,}3\) 
s osou \(y\).}
{\([0;-6{,}3]\)}{\([10;0]\)}{\([0{,}3;0]\)}{Graf funkce \(f\) 
neprotíná osu \(y\).}{}{}}
\LANGsk{
\Question{
Daná je funkcia \(f\colon y = 10x^{2} - 18x - 6{,}3\). Určte priesečníky grafu funkcie s osou \(y\).}
{\([0;-6{,}3]\)}{\([10;0]\)}{\([0{,}3;0]\)}{Graf funkcie \(f\) 
nepretína os \(y\).}{}{}}
\LANGes{
\Question{
Encuentra el punto de intersección del eje  \(y\)
con la función \(f\colon y = 10x^{2} - 18x - 6.3\).}
{\([0;-6.3]\)}{\([10;0]\)}{\([0.3;0]\)}{La función no interseca el eje \(y\)}{}{}}
\End

\Data\ID{9000014810}
\Updated{2020-11-20 12:11:22}
\Level{A}
\SubArea{Quadratic Functions}
\IMG{000148_10-figure0.pdf}
\LANGen{
\Question{
Find the domain and range of the quadratic function
\(f\) 
graphed in the picture.}
{\(\begin{aligned}[t] &\mathop{\mathrm{Dom}}(f) =\mathbb{R} &
\\&\mathop{\mathrm{Ran}}(f) = \left (-\infty ;2\right ] 
\\ \end{aligned}\)}{\(\begin{aligned}[t] &\mathop{\mathrm{Dom}}(f) =\mathbb{R} &
\\&\mathop{\mathrm{Ran}}(f) = \left [ 2;\infty \right )
\\ \end{aligned}\)}{\(\begin{aligned}[t] &\mathop{\mathrm{Dom}}(f) = \left [ 0;\infty \right )&
\\&\mathop{\mathrm{Ran}}(f) = \left [ 2;4\right ] 
\\ \end{aligned}\)}{\(\begin{aligned}[t] &\mathop{\mathrm{Dom}}(f) = \left (-\infty ;0\right ] &
\\&\mathop{\mathrm{Ran}}(f) =\mathbb{R}
\\ \end{aligned}\)}{}{}}
\LANGpl{
\Question{
Znajdź dziedzinę i zakres funkcji kwadratowej
\(f\) 
przedstawionej na rysunku poniżej.}
{\(\begin{aligned}[t] &\mathop{\mathrm{Dom}}(f) =\mathbb{R} &
\\&\mathop{\mathrm{Ran}}(f) = \left (-\infty ;2\right ] 
\\ \end{aligned}\)}{\(\begin{aligned}[t] &\mathop{\mathrm{Dom}}(f) =\mathbb{R} &
\\&\mathop{\mathrm{Ran}}(f) = \left [ 2;\infty \right )
\\ \end{aligned}\)}{\(\begin{aligned}[t] &\mathop{\mathrm{Dom}}(f) = \left [ 0;\infty \right )&
\\&\mathop{\mathrm{Ran}}(f) = \left [ 2;4\right ] 
\\ \end{aligned}\)}{\(\begin{aligned}[t] &\mathop{\mathrm{Dom}}(f) = \left (-\infty ;0\right ] &
\\&\mathop{\mathrm{Ran}}(f) =\mathbb{R}
\\ \end{aligned}\)}{}{}}
\LANGcs{
\Question{
Je dán graf kvadratické funkce \(f\).
Která z následujících možností správně popisuje vlastnosti funkce na
obrázku?}
{\(\begin{aligned}[t] &D(f) =\mathbb{R} &
\\&H(f) = \left (-\infty ;2\right \rangle 
\\ \end{aligned}\)}{\(\begin{aligned}[t] &D(f) =\mathbb{R} &
\\&H(f) = \left \langle 2;\infty \right )
\\ \end{aligned}\)}{\(\begin{aligned}[t] &D(f) = \left \langle 0;\infty \right ) &
\\&H(f) = \left \langle 2;4\right \rangle 
\\ \end{aligned}\)}{\(\begin{aligned}[t] &D(f) = \left (-\infty ;0\right \rangle &
\\&H(f) =\mathbb{R}
\\ \end{aligned}\)}{}{}}
\LANGsk{
\Question{
Určte definičný obor a obor hodnôt kvadratickej funkcie \(f\), ktorej graf je na obrázku.}
{\(\begin{aligned}[t] &D(f) =\mathbb{R} &
\\&H(f) = \left (-\infty ;2\right \rangle 
\\ \end{aligned}\)}{\(\begin{aligned}[t] &D(f) =\mathbb{R} &
\\&H(f) = \left \langle 2;\infty \right )
\\ \end{aligned}\)}{\(\begin{aligned}[t] &D(f) = \left \langle 0;\infty \right ) &
\\&H(f) = \left \langle 2;4\right \rangle 
\\ \end{aligned}\)}{\(\begin{aligned}[t] &D(f) = \left (-\infty ;0\right \rangle &
\\&H(f) =\mathbb{R}
\\ \end{aligned}\)}{}{}}
\LANGes{
\Question{
Encuentra el Dominio y el Rango de la función cuadrática
\(f\) 
cuya gráfica está en el dibujo.}
{\(\begin{aligned}[t] &\mathop{\mathrm{Dom}}(f) =\mathbb{R} &
\\&\mathop{\mathrm{Ran}}(f) = \left (-\infty ;2\right ] 
\\ \end{aligned}\)}{\(\begin{aligned}[t] &\mathop{\mathrm{Dom}}(f) =\mathbb{R} &
\\&\mathop{\mathrm{Ran}}(f) = \left [ 2;\infty \right )
\\ \end{aligned}\)}{\(\begin{aligned}[t] &\mathop{\mathrm{Dom}}(f) = \left [ 0;\infty \right )&
\\&\mathop{\mathrm{Ran}}(f) = \left [ 2;4\right ] 
\\ \end{aligned}\)}{\(\begin{aligned}[t] &\mathop{\mathrm{Dom}}(f) = \left (-\infty ;0\right ] &
\\&\mathop{\mathrm{Ran}}(f) =\mathbb{R}
\\ \end{aligned}\)}{}{}}
\End

\Data\ID{9000022302}
\Updated{2021-04-07 07:04:28}
\Level{A}
\SubArea{Quadratic Functions}
\LANGen{
\Question{
Complete the following statement, if possible: „The function \(f\colon y = -x^{2} - 2x + 15\) attains only
... values on the interval \([- 5;3] \).”}
{nonnegative}{positive}{negative}{neither of the above is valid}{}{}}
\LANGpl{
\Question{
Jeśli to możliwe dokończ następujące zdanie: „Funkcja \(f\colon y = -x^{2} - 2x + 15\) przyjmuje jedynie
... wartości w przedziale \([- 5;3] \).”}
{nieujemne}{dodatnie}{ujemne}{żadne z powyższych nie jest prawdziwe}{}{}}
\LANGcs{
\Question{
Na intervalu \(\langle - 5;3\rangle \) jsou všechny
funkční hodnoty funkce \(f\colon y = -x^{2} - 2x + 15\):}
{nezáporné}{kladné}{záporné}{žádná z možností}{}{}}
\LANGsk{
\Question{
Doplňte nasledujúce tvrdenie:  „Funkcia \(f\colon y = -x^{2} - 2x + 15\)  dosahuje len ... hodnoty na intervale \(\langle - 5;3\rangle \)."}
{nezáporné}{kladné}{záporné}{žiadna z možností}{}{}}
\LANGes{
\Question{
Completa la siguiente frase, si es posible: "La función \(f\colon y = -x^{2} - 2x + 15\) alcanza solamente valores ... en el intervalo \([- 5;3] \).”}
{no negativos}{positivos}{negativos}{ninguna respuesta está correcta}{}{}}
\End

\Data\ID{1003083113}
\Updated{2021-02-15 07:02:04}
\Level{B}
\SubArea{Quadratic Functions}
\LANGen{
\Question{
Determine the coordinates of the parabola vertex, if  the parabola is the graph of \( f(x)=-\frac43 x^2+8x-13 \).}
{\( [3;-1] \)}{\( [3;-22] \)}{\( [-3;-1] \)}{\( [-3;22] \)}{}{}}
\LANGpl{
\Question{
Określ współrzędne wierzchołka paraboli, jeśli parabola jest wykresem \( f(x)=-\frac43 x^2+8x-13 \).}
{\( [3;-1] \)}{\( [3;-22] \)}{\( [-3;-1] \)}{\( [-3;22] \)}{}{}}
\LANGcs{
\Question{
Určete souřadnice vrcholu paraboly, která je grafem funkce \( f(x)=-\frac43 x^2+8x-13 \).}
{\( [3;-1] \)}{\( [3;-22] \)}{\( [-3;-1] \)}{\( [-3;22] \)}{}{}}
\LANGsk{
\Question{
Vypočítajte súradnice vrcholu paraboly, ktorá je daná grafom funkcie \( f(x)=-\frac43 x^2+8x-13 \).}
{\( [3;-1] \)}{\( [3;-22] \)}{\( [-3;-1] \)}{\( [-3;22] \)}{}{}}
\LANGes{
\Question{
Determina las coordenadas del vértice de la parábola, siendo la parábola es la gráfica de \( f(x)=-\frac43 x^2+8x-13 \)}
{\( [3;-1] \)}{\( [3;-22] \)}{\( [-3;-1] \)}{\( [-3;22] \)}{}{}}
\End

\Data\ID{1003083114}
\Updated{2022-04-23 05:04:17}
\Level{B}
\SubArea{Quadratic Functions}
\LANGen{
\Question{
The graph of the function \( f(x)=4x^2+16x+11 \) is a parabola. Which of the following points is the vertex of this parabola?}
{\( [-2;-5] \)}{\( [-2;7] \)}{\( [2;11] \)}{\( [-2;11] \)}{}{}}
\LANGpl{
\Question{
Określ współrzędne wierzchołka paraboli, jeśli parabola jest wykresem \( f(x)=4x^2+16x+11 \).}
{\( [-2;-5] \)}{\( [-2;7] \)}{\( [2;11] \)}{\( [-2;11] \)}{}{}}
\LANGcs{
\Question{
Určete souřadnice vrcholu paraboly, která je grafem funkce \( f(x)=4x^2+16x+11 \).}
{\( [-2;-5] \)}{\( [-2;7] \)}{\( [2;11] \)}{\( [-2;11] \)}{}{}}
\LANGsk{
\Question{
Vypočítajte súradnice vrcholu paraboly, ktorá je daná grafom funkcie \( f(x)=4x^2+16x+11 \).}
{\( [-2;-5] \)}{\( [-2;7] \)}{\( [2;11] \)}{\( [-2;11] \)}{}{}}
\LANGes{
\Question{
Halla el vértice de la parábola definida por  \( f(x)=4x^2+16x+11 \).}
{\( [-2;-5] \)}{\( [-2;7] \)}{\( [2;11] \)}{\( [-2;11] \)}{}{}}
\End

\Data\ID{1003108301}
\Updated{2020-09-30 03:09:54}
\Level{B}
\SubArea{Quadratic Functions}
\LANGen{
\Question{
The graph of the quadratic function \( f \) intersects the coordinate axes at the points \( [1;0] \), \( [4;0] \), \( [0;8] \). Find the function \( f \).}
{\( f(x)=2x^2-10x+8 \)}{\( f(x)=x^2-5x+4 \)}{\( f(x)=-2x^2+10x-8 \)}{\( f(x)=2x^2+10x+8 \)}{}{}}
\LANGpl{
\Question{
Wykres funkcji kwadratowej \( f \) przecina osie współrzędnych w punktach  \( [1;0] \), \( [4;0] \), \( [0;8] \). Znajdź funkcję \( f \).}
{\( f(x)=2x^2-10x+8 \)}{\( f(x)=x^2-5x+4 \)}{\( f(x)=-2x^2+10x-8 \)}{\( f(x)=2x^2+10x+8 \)}{}{}}
\LANGcs{
\Question{
Graf kvadratické funkce \( f \) protíná souřadnicové osy v bodech \( [1;0] \), \( [4;0] \), \( [0;8] \). Z nabídnutých možností vyberte předpis této funkce.}
{\( f(x)=2x^2-10x+8 \)}{\( f(x)=x^2-5x+4 \)}{\( f(x)=-2x^2+10x-8 \)}{\( f(x)=2x^2+10x+8 \)}{}{}}
\LANGsk{
\Question{
Graf kvadratickej funkcie \( f \) pretína súradnicové osi v bodoch \( [1;0] \), \( [4;0] \), \( [0;8] \). Nájdite predpis funkcie \( f \).}
{\( f(x)=2x^2-10x+8 \)}{\( f(x)=x^2-5x+4 \)}{\( f(x)=-2x^2+10x-8 \)}{\( f(x)=2x^2+10x+8 \)}{}{}}
\LANGes{
\Question{
La gráfica de la función cuadrática \( f \) interseca los ejes de coordenadas en los puntos \( [1;0] \), \( [4;0] \), \( [0;8] \). Encuentra la función \( f \).}
{\( f(x)=2x^2-10x+8 \)}{\( f(x)=x^2-5x+4 \)}{\( f(x)=-2x^2+10x-8 \)}{\( f(x)=2x^2+10x+8 \)}{}{}}
\End

\Data\ID{1003108302}
\Updated{2020-09-30 03:09:45}
\Level{B}
\SubArea{Quadratic Functions}
\LANGen{
\Question{
The graph of the quadratic function \( f \) is a parabola with the vertex \( [2;5] \). The parabola intersects the \(  y \)-axis at the point \( [0;3] \). Find the function \( f \).}
{\( f(x)=-\frac12(x-2)^2+5 \)}{\( f(x)=-\frac12(x+2)^2+5 \)}{\( f(x)=-2(x-2)^2+5 \)}{\( f(x)=-2(x+2)^2+5 \)}{}{}}
\LANGpl{
\Question{
Wykres funkcji kwadratowej  \( f \) jest parabolą z wierzchołkiem \( [2;5] \). Parabola przecina oś współrzędnych \(  y \) w punkcie \( [0;3] \). Wyznacz funkcję \( f \).}
{\( f(x)=-\frac12(x-2)^2+5 \)}{\( f(x)=-\frac12(x+2)^2+5 \)}{\( f(x)=-2(x-2)^2+5 \)}{\( f(x)=-2(x+2)^2+5 \)}{}{}}
\LANGcs{
\Question{
Grafem kvadratické funkce \( f \) je parabola s vrcholem \( [2;5] \), která protíná osu \(  y \) v bodě \( [0;3] \). Z nabídnutých možností vyberte předpis této funkce.}
{\( f(x)=-\frac12(x-2)^2+5 \)}{\( f(x)=-\frac12(x+2)^2+5 \)}{\( f(x)=-2(x-2)^2+5 \)}{\( f(x)=-2(x+2)^2+5 \)}{}{}}
\LANGsk{
\Question{
Grafom kvadratickej funkcie \( f \) je parabola s vrcholom \( [2;5] \), ktorá pretína os \(  y \) v bode \( [0;3] \). Nájdite predpis funkcie \( f \).}
{\( f(x)=-\frac12(x-2)^2+5 \)}{\( f(x)=-\frac12(x+2)^2+5 \)}{\( f(x)=-2(x-2)^2+5 \)}{\( f(x)=-2(x+2)^2+5 \)}{}{}}
\LANGes{
\Question{
La gráfica de una función cuadrática \( f \) es una parábola con el vértice en \( [2;5] \).  La parábola interseca al eje \(  y \) en el punto \( [0;3] \) Encuentra la función \( f \).}
{\( f(x)=-\frac12(x-2)^2+5 \)}{\( f(x)=-\frac12(x+2)^2+5 \)}{\( f(x)=-2(x-2)^2+5 \)}{\( f(x)=-2(x+2)^2+5 \)}{}{}}
\End

\Data\ID{1003108303}
\Updated{2021-02-15 07:02:54}
\Level{B}
\SubArea{Quadratic Functions}
\LANGen{
\Question{
Maximum value of the quadratic function \( f \) is \( 2 \). The graph of \( f \) intersects the \( x \)-axis at the points       \( [-1;0] \) and \( [3;0] \). Find the function \( f \).}
{\( f(x)=-\frac12x^2+x+\frac32 \)}{\( f(x)=x^2-2x+3 \)}{\( f(x)=x^2-2x-3 \)}{\( f(x)=-\frac12x^2-x+\frac32 \)}{}{}}
\LANGpl{
\Question{
Maksymalna wartość funkcji kwadratowej \( f \) wynosi \( 2 \). wykres funkcji  \( f \) przecina oś współrzędnych \( x \) w punktach \( [-1;0] \) i \( [3;0] \). Wyznacz funkcję \( f \).}
{\( f(x)=-\frac12x^2+x+\frac32 \)}{\( f(x)=x^2-2x+3 \)}{\( f(x)=x^2-2x-3 \)}{\( f(x)=-\frac12x^2-x+\frac32 \)}{}{}}
\LANGcs{
\Question{
Hodnota maxima kvadratické funkce \( f \) je \( 2 \). Z nabídnutých možností vyberte předpis této funkce, víme-li, že graf této funkce protíná osu \( x \) v bodech \( [-1;0] \) a \( [3;0] \).}
{\( f(x)=-\frac12x^2+x+\frac32 \)}{\( f(x)=x^2-2x+3 \)}{\( f(x)=x^2-2x-3 \)}{\( f(x)=-\frac12x^2-x+\frac32 \)}{}{}}
\LANGsk{
\Question{
Maximum kvadratickej funkcie \( f \) má hodnotu \( 2 \) a pretína os \( x \) v bodoch \( [-1;0] \) a \( [3;0] \). Nájdite predpis funkcie \( f \).}
{\( f(x)=-\frac12x^2+x+\frac32 \)}{\( f(x)=x^2-2x+3 \)}{\( f(x)=x^2-2x-3 \)}{\( f(x)=-\frac12x^2-x+\frac32 \)}{}{}}
\LANGes{
\Question{
El valor máximo de la función cuadrática \( f \) es \( 2 \). La gráfica de \( f \) interseca al eje  \( x \) en los puntos  \( [-1;0] \) y \( [3;0] \). Encuentra la función \( f \).}
{\( f(x)=-\frac12x^2+x+\frac32 \)}{\( f(x)=x^2-2x+3 \)}{\( f(x)=x^2-2x-3 \)}{\( f(x)=-\frac12x^2-x+\frac32 \)}{}{}}
\End

\Data\ID{1003108304}
\Updated{2020-09-30 03:09:04}
\Level{B}
\SubArea{Quadratic Functions}
\LANGen{
\Question{
The graph of the function \( f \) intersects the \( x \)-axis at the points \( [-5;0] \) and \( [-1;0] \). Given \( f(-2)=6 \), find the function \( f \).}
{\( f(x)=-2x^2-12x-10 \)}{\( f(x)=x^2+6x+5 \)}{\( f(x)=-2(x-5)(x-1) \)}{\( f(x)=-2(x+3)^2+6 \)}{}{}}
\LANGpl{
\Question{
Wykres funkcji \( f \) przecina oś współrzędnych \( x \) w punktach \( [-5;0] \) i \( [-1;0] \).  Mając dane \( f(-2)=6 \), wyznacz funkcję \( f \).}
{\( f(x)=-2x^2-12x-10 \)}{\( f(x)=x^2+6x+5 \)}{\( f(x)=-2(x-5)(x-1) \)}{\( f(x)=-2(x+3)^2+6 \)}{}{}}
\LANGcs{
\Question{
Graf kvadratické funkce \( f \) protíná osu \( x \) v bodech \( [-5;0] \) a \( [-1;0] \). Z nabídnutých možností vyberte předpis této funkce, víme-li, že platí \( f(-2)=6 \).}
{\( f(x)=-2x^2-12x-10 \)}{\( f(x)=x^2+6x+5 \)}{\( f(x)=-2(x-5)(x-1) \)}{\( f(x)=-2(x+3)^2+6 \)}{}{}}
\LANGsk{
\Question{
Graf kvadratickej funkcie \( f \) pretína os \( x \) v bodoch \( [-5;0] \) a \( [-1;0] \). Nájdite predpis funkcie \( f \), ak vieme, že platí \( f(-2)=6 \).}
{\( f(x)=-2x^2-12x-10 \)}{\( f(x)=x^2+6x+5 \)}{\( f(x)=-2(x-5)(x-1) \)}{\( f(x)=-2(x+3)^2+6 \)}{}{}}
\LANGes{
\Question{
La gráfica de la función \( f \) interseca el eje \( x \) en los puntos \( [-5;0] \) y \( [-1;0] \). Dado \( f(-2)=6 \). Encuentra la función \( f \).}
{\( f(x)=-2x^2-12x-10 \)}{\( f(x)=x^2+6x+5 \)}{\( f(x)=-2(x-5)(x-1) \)}{\( f(x)=-2(x+3)^2+6 \)}{}{}}
\End

\Data\ID{1003108305}
\Updated{2020-09-30 03:09:01}
\Level{B}
\SubArea{Quadratic Functions}
\LANGen{
\Question{
Find the quadratic function \( f \) whose graph is passing through the points \( [0;-9] \), \( [2;-3] \), \( [6;-3] \).}
{\( f(x)=-\frac12x^2+4x-9 \)}{\( f(x)=-\frac12x^2+4x-3 \)}{\( f(x)=x^2-8x-9 \)}{\( f(x)=-2x^2+7x-9 \)}{}{}}
\LANGpl{
\Question{
Wyznacz funkcję kwadratową  \( f \), której wykres przechodzi przez punkty \( [0;-9] \), \( [2;-3] \), \( [6;-3] \).}
{\( f(x)=-\frac12x^2+4x-9 \)}{\( f(x)=-\frac12x^2+4x-3 \)}{\( f(x)=x^2-8x-9 \)}{\( f(x)=-2x^2+7x-9 \)}{}{}}
\LANGcs{
\Question{
Graf kvadratické funkce \( f \) prochází body \( [0;-9] \), \( [2;-3] \), \( [6;-3] \). Z nabídnutých možností vyberte předpis této funkce.}
{\( f(x)=-\frac12x^2+4x-9 \)}{\( f(x)=-\frac12x^2+4x-3 \)}{\( f(x)=x^2-8x-9 \)}{\( f(x)=-2x^2+7x-9 \)}{}{}}
\LANGsk{
\Question{
Nájdite predpis kvadratickej funkcie \( f \), ktorej graf prechádza bodmi \( [0;-9] \), \( [2;-3] \), \( [6;-3] \).}
{\( f(x)=-\frac12x^2+4x-9 \)}{\( f(x)=-\frac12x^2+4x-3 \)}{\( f(x)=x^2-8x-9 \)}{\( f(x)=-2x^2+7x-9 \)}{}{}}
\LANGes{
\Question{
Encuentra la función cuadrática \( f \) cuya gráfica pasa por los puntos \( [0;-9] \), \( [2;-3] \), \( [6;-3] \).}
{\( f(x)=-\frac12x^2+4x-9 \)}{\( f(x)=-\frac12x^2+4x-3 \)}{\( f(x)=x^2-8x-9 \)}{\( f(x)=-2x^2+7x-9 \)}{}{}}
\End

\Data\ID{1003108306}
\Updated{2021-02-15 07:02:04}
\Level{B}
\SubArea{Quadratic Functions}
\LANGen{
\Question{
The \( x \)-axis is the tangent line of the graph of the quadratic function \( f \). The point of tangency has the coordinates \( [-2;0] \). Given \( f(-1)=-4 \), find the function \( f \).}
{\( f(x)=-4x^2-16x-16 \)}{\( f(x)=-4x^2-16x+16 \)}{\( f(x)=-\frac49x^2+\frac{16}9x-\frac{16}9 \)}{\( f(x)=4x^2-16x+16 \)}{}{}}
\LANGpl{
\Question{
Oś współrzędnych \( x \) jest styczną wykresu funkcji kwadratowej \( f \).  Punkt styczności ma współrzędne \( [-2;0] \). Mając dane \( f(-1)=-4 \), wyznacz funkcję \( f \).}
{\( f(x)=-4x^2-16x-16 \)}{\( f(x)=-4x^2-16x+16 \)}{\( f(x)=-\frac49x^2+\frac{16}9x-\frac{16}9 \)}{\( f(x)=4x^2-16x+16 \)}{}{}}
\LANGcs{
\Question{
Osa \( x \) je tečnou grafu kvadratické funkce \( f \). Bod dotyku má souřadnice \( [-2;0] \). Z nabídnutých možností vyberte předpis této funkce, víme-li, že platí \( f(-1)=-4 \).}
{\( f(x)=-4x^2-16x-16 \)}{\( f(x)=-4x^2-16x+16 \)}{\( f(x)=-\frac49x^2+\frac{16}9x-\frac{16}9 \)}{\( f(x)=4x^2-16x+16 \)}{}{}}
\LANGsk{
\Question{
Os \( x \) je dotyčnicou ku grafu kvadratickej funkcie \( f \) v bode \( [-2;0] \). Nájdite predpis funkcie \( f \), ak vieme, že platí \( f(-1)=-4 \).}
{\( f(x)=-4x^2-16x-16 \)}{\( f(x)=-4x^2-16x+16 \)}{\( f(x)=-\frac49x^2+\frac{16}9x-\frac{16}9 \)}{\( f(x)=4x^2-16x+16 \)}{}{}}
\LANGes{
\Question{
El eje \( x \) es tangente a la gráfica de la función cuadrática \( f \). El punto de tangencia tiene coordenadas \( [-2;0] \). Sabiendo que \( f(-1)=-4 \), encuentra la función \( f \).}
{\( f(x)=-4x^2-16x-16 \)}{\( f(x)=-4x^2-16x+16 \)}{\( f(x)=-\frac49x^2+\frac{16}9x-\frac{16}9 \)}{\( f(x)=4x^2-16x+16 \)}{}{}}
\End

\Data\ID{1003108308}
\Updated{2021-02-15 07:02:33}
\Level{B}
\SubArea{Quadratic Functions}
\LANGen{
\Question{
Which of the following information are insufficient to determine quadratic function uniquely?}
{two intersection points with the \( x \)-axis and \( x \)-coordinate of the vertex}{two intersection points with the \( x \)-axis and \( y \)-coordinate of the vertex}{two intersection points with the \( x \)-axis and any other point of the function}{coordinates of the vertex and the intersection with the \( y \)-axis}{}{}}
\LANGpl{
\Question{
Które z poniższych informacji są niewystarczające do jednoznacznego określenia funkcji kwadratowej?}
{dwa punkty przecięcia z osią współrzędnych  \( x \)  i \( x \)-współrzędną wierzchołka}{dwa punkty przecięcia z osią współrzędnych  \( x \) i  \( y \) współrzędną wierzchołka}{dwa punkty przecięcia z osią współrzędnych \( x \) i każdym innym punktem funkcji}{współrzędne wierzchołka i przecięcia z osią współrzędnych  \( y \)}{}{}}
\LANGcs{
\Question{
Které z následujících informací jsou nedostatečné k jednoznačnému určení kvadratické funkce?}
{dva průsečíky s osou \( x \) a \( x \)-souřadnice vrcholu}{dva průsečíky s osou  \( x \) a \( y \)-souřadnice vrcholu}{dva průsečíky s osou  \( x \) a libovolný další bod funkce}{souřadnice vrcholu a průsečík s osou \( y \)}{}{}}
\LANGsk{
\Question{
Ktorá z nasledujúcich možností obsahuje nedostatočné informácie k jednoznačnému určeniu kvadratickej funkcie?}
{dva priesečníky s osou \( x \) a \( x \)-súradnica vrchola}{dva priesečníky s osou \( x \) a \( y \)-súradnica vrchola}{dva priesečníky s osou \( x \) a ľubovoľný ďalší bod funkcie}{súradnice vrchola a priesečník s osou \( y \)}{}{}}
\LANGes{
\Question{
¿Cuál de las siguientes informaciones no basta para determinar una función cuadrática?}
{dos puntos de intersección con el eje \( x \) y la coordenada \( x \) del vértice.}{dos puntos de intersección con el eje \( x \) y la coordenada \( y \) del vértice.}{dos puntos de intersección con el eje \( x \) y cualquier punto de la gráfica}{coordinadas del vértice y la intersección con el eje \( y \)}{}{}}
\End

\Data\ID{1003108309}
\Updated{2021-02-15 07:02:01}
\Level{B}
\SubArea{Quadratic Functions}
\LANGen{
\Question{
The graph of the quadratic function \( f \) intersects coordinate axes at the points \( [-3;0] \), \( [1;0] \), \( \left[0;\frac32\right] \). Find the function \( f \).}
{\( f(x)=-\frac12(x+1)^2+2 \)}{\( f(x)=-\frac12(x+1)^2+\frac12 \)}{\( f(x)=-\frac12(x-1)^2+2 \)}{\( f(x)=\frac12(x-1)^2+2 \)}{}{}}
\LANGpl{
\Question{
Wykres funkcji kwadratowej  \( f \) przecina osie współrzędnych i punkty  \( [-3;0] \), \( [1;0] \), \( \left[0;\frac32\right] \). Wyznacz funkcję \( f \).}
{\( f(x)=-\frac12(x+1)^2+2 \)}{\( f(x)=-\frac12(x+1)^2+\frac12 \)}{\( f(x)=-\frac12(x-1)^2+2 \)}{\( f(x)=\frac12(x-1)^2+2 \)}{}{}}
\LANGcs{
\Question{
Graf kvadratické funkce \( f \) protíná souřadnicové osy v bodech \( [-3;0] \), \( [1;0] \), \( \left[0;\frac32\right] \). Určete funkci \( f \).}
{\( f(x)=-\frac12(x+1)^2+2 \)}{\( f(x)=-\frac12(x+1)^2+\frac12 \)}{\( f(x)=-\frac12(x-1)^2+2 \)}{\( f(x)=\frac12(x-1)^2+2 \)}{}{}}
\LANGsk{
\Question{
Graf kvadratickej funkcie \( f \) pretína súradnicové osi v bodoch \( [-3;0] \), \( [1;0] \), \( \left[0;\frac32\right] \). Nájdite predpis funkcie \( f \).}
{\( f(x)=-\frac12(x+1)^2+2 \)}{\( f(x)=-\frac12(x+1)^2+\frac12 \)}{\( f(x)=-\frac12(x-1)^2+2 \)}{\( f(x)=\frac12(x-1)^2+2 \)}{}{}}
\LANGes{
\Question{
La gráfica de la función cuadrátia \( f \) corta a los ejes de coordenadas en los puntos \( [-3;0] \), \( [1;0] \), \( \left[0;\frac32\right] \). Encuentra la función \( f \).}
{\( f(x)=-\frac12(x+1)^2+2 \)}{\( f(x)=-\frac12(x+1)^2+\frac12 \)}{\( f(x)=-\frac12(x-1)^2+2 \)}{\( f(x)=\frac12(x-1)^2+2 \)}{}{}}
\End

\Data\ID{1003108310}
\Updated{2021-02-15 07:02:25}
\Level{B}
\SubArea{Quadratic Functions}
\LANGen{
\Question{
The graph of the quadratic function \( f \) has the vertex at the point \( [3; -1] \) and it passes through the point \( [ -1; 3] \).  Find the function \( f \).}
{\( f(x)=\frac14x^2-\frac32x+\frac54 \)}{\( f(x)=\frac14x^2+\frac32x+\frac54 \)}{\( f(x)=-\frac14x^2+\frac32x-\frac{13}4 \)}{\( f(x)=x^2+6x+8 \)}{}{}}
\LANGpl{
\Question{
Wykres funkcji kwadratowej  \( f \) ma wierzchołek w punkcie \( [3; -1] \) i przechodzi przez punkt \( [ -1; 3] \).  Wyznacz  \( f \).}
{\( f(x)=\frac14x^2-\frac32x+\frac54 \)}{\( f(x)=\frac14x^2+\frac32x+\frac54 \)}{\( f(x)=-\frac14x^2+\frac32x-\frac{13}4 \)}{\( f(x)=x^2+6x+8 \)}{}{}}
\LANGcs{
\Question{
Graf kvadratické funkce \( f \) má vrchol v bodě \( [3; -1] \) a prochází bodem \( [ -1; 3] \).  Určete předpis funkce \( f \).}
{\( f(x)=\frac14x^2-\frac32x+\frac54 \)}{\( f(x)=\frac14x^2+\frac32x+\frac54 \)}{\( f(x)=-\frac14x^2+\frac32x-\frac{13}4 \)}{\( f(x)=x^2+6x+8 \)}{}{}}
\LANGsk{
\Question{
Graf kvadratickej funkcie \( f \) má vrchol v bode \( [3; -1] \) a prechádza bodom \( [ -1; 3] \).  Nájdite predpis funkcie \( f \).}
{\( f(x)=\frac14x^2-\frac32x+\frac54 \)}{\( f(x)=\frac14x^2+\frac32x+\frac54 \)}{\( f(x)=-\frac14x^2+\frac32x-\frac{13}4 \)}{\( f(x)=x^2+6x+8 \)}{}{}}
\LANGes{
\Question{
La gráfica de la función cuadrática \( f \) tiene el vértice en el punto \( [3; -1] \) y pasa por el punto \( [ -1; 3] \). Halla la función \( f \).}
{\( f(x)=\frac14x^2-\frac32x+\frac54 \)}{\( f(x)=\frac14x^2+\frac32x+\frac54 \)}{\( f(x)=-\frac14x^2+\frac32x-\frac{13}4 \)}{\( f(x)=x^2+6x+8 \)}{}{}}
\End

\Data\ID{1003108311}
\Updated{2021-02-15 07:02:49}
\Level{B}
\SubArea{Quadratic Functions}
\LANGen{
\Question{
The quadratic function \( f \) has the minimum at \( x=-2 \) and its graph passes through the points \( [0;13] \), \( [-1; 4] \). Find the function \( f \).}
{\( f(x)=3(x+2)^2+1 \)}{\( f(x)=-\frac59(x-2)^2+9 \)}{\( f(x)=\frac59(x-2)^2+9 \)}{\( f(x)=3(x+2)^2-1 \)}{}{}}
\LANGpl{
\Question{
Funkcja kwadratowa  \( f \) ma minimum w  \( x=-2 \), a jej wykres przechodzi przez punkty  \( [0;13] \), \( [-1; 4] \). Wyznacz funkcję \( f \).}
{\( f(x)=3(x+2)^2+1 \)}{\( f(x)=-\frac59(x-2)^2+9 \)}{\( f(x)=\frac59(x-2)^2+9 \)}{\( f(x)=3(x+2)^2-1 \)}{}{}}
\LANGcs{
\Question{
Kvadratická funkce \( f \) má minimum, které nastává v \( x=-2 \) a její graf prochází body \( [0;13] \), \( [-1; 4] \). Určete funkci \( f \).}
{\( f(x)=3(x+2)^2+1 \)}{\( f(x)=-\frac59(x-2)^2+9 \)}{\( f(x)=\frac59(x-2)^2+9 \)}{\( f(x)=3(x+2)^2-1 \)}{}{}}
\LANGsk{
\Question{
Graf kvadratickej funkcie \( f \) má minimum v bode \( x=-2 \) a prechádza bodmi \( [0;13] \), \( [-1; 4] \). Nájdite predpis funkcie \( f \).}
{\( f(x)=3(x+2)^2+1 \)}{\( f(x)=-\frac59(x-2)^2+9 \)}{\( f(x)=\frac59(x-2)^2+9 \)}{\( f(x)=3(x+2)^2-1 \)}{}{}}
\LANGes{
\Question{
La función cuadrática \( f \) alcanza su valor mínimo en 0 para \( x=-2 \) y su gráfica pasa por los puntos \( [0;13] \), \( [-1; 4] \). Halla la función \( f \).}
{\( f(x)=3(x+2)^2+1 \)}{\( f(x)=-\frac59(x-2)^2+9 \)}{\( f(x)=\frac59(x-2)^2+9 \)}{\( f(x)=3(x+2)^2-1 \)}{}{}}
\End

\Data\ID{1003108312}
\Updated{2022-04-23 05:04:17}
\Level{B}
\SubArea{Quadratic Functions}
\LANGen{
\Question{
The graph of the function \( f \) is a parabola, vertex of which is \( [6;0] \) and \( f(2)= 8 \) is given. Find the function \( f \).}
{\( f(x)=\frac12(x-6)^2 \)}{\( f(x)=-\frac12(x-6)^2 \)}{\( f(x)=\frac12(x+6)^2 \)}{\( f(x)=\frac12x^2+6 \)}{}{}}
\LANGpl{
\Question{
Wykres funkcji  \( f \) jest parabolą, której wierzchołkiem jest \( [6;0] \), a \( f(2)= 8 \) . Wyznacz funkcję  \( f \).}
{\( f(x)=\frac12(x-6)^2 \)}{\( f(x)=-\frac12(x-6)^2 \)}{\( f(x)=\frac12(x+6)^2 \)}{\( f(x)=\frac12x^2+6 \)}{}{}}
\LANGcs{
\Question{
Grafem funkce \( f \) je parabola s vrcholem v bodě \( [6;0] \) a dále platí \( f(2)= 8 \). Určete funkci \( f \).}
{\( f(x)=\frac12(x-6)^2 \)}{\( f(x)=-\frac12(x-6)^2 \)}{\( f(x)=\frac12(x+6)^2 \)}{\( f(x)=\frac12x^2+6 \)}{}{}}
\LANGsk{
\Question{
Grafom kvadratickej funkcie \( f \) je parabola s vrcholom v bode \( [6;0] \) a platí \( f(2)= 8 \). Nájdite predpis funkcie \( f \).}
{\( f(x)=\frac12(x-6)^2 \)}{\( f(x)=-\frac12(x-6)^2 \)}{\( f(x)=\frac12(x+6)^2 \)}{\( f(x)=\frac12x^2+6 \)}{}{}}
\LANGes{
\Question{
La gráfica de la función \( f \) es una parábola, cuyo vértice es el punto \( [6;0] \) y sabemos que \( f(2)= 8 \). Encuentra la función \( f \).}
{\( f(x)=\frac12(x-6)^2 \)}{\( f(x)=-\frac12(x-6)^2 \)}{\( f(x)=\frac12(x+6)^2 \)}{\( f(x)=\frac12x^2+6 \)}{}{}}
\End

\Data\ID{1103034601}
\Updated{2021-02-15 07:02:37}
\Level{B}
\SubArea{Quadratic Functions}
\IMG{1103034601.pdf}
\LANGen{
\Question{
Given graphs of the functions \( f(x)=x^2-6x+8\) and \( g(x)=-x^2-2x+24 \), find the solution set of the following inequality.
\[ x^2-6x+8\leq-x^2-2x+24 \]}
{\( [-2;4] \)}{\( [2;4] \)}{\( [0;24] \)}{\( [-6;4] \)}{}{}}
\LANGpl{
\Question{
Dane są wykresy funkcji \( f(x)=x^2-6x+8\) i \( g(x)=-x^2-2x+24 \), określ zbiór rozwiązań następującej nierówności.
\[ x^2-6x+8\leq-x^2-2x+24 \]}
{\( \langle-2;4\rangle \)}{\( \langle2;4\rangle \)}{\( \langle0;24\rangle \)}{\( \langle-6;4\rangle \)}{}{}}
\LANGcs{
\Question{
Pomocí grafů funkcí \( f(x)=x^2-6x+8\) a \( g(x)=-x^2-2x+24 \) určete množinu řešení dané nerovnice.
\[ x^2-6x+8\leq-x^2-2x+24 \]}
{\( \langle-2;4\rangle \)}{\( \langle2;4\rangle \)}{\( \langle0;24\rangle \)}{\( \langle-6;4\rangle \)}{}{}}
\LANGsk{
\Question{
Pomocou grafov funkcií \( f(x)=x^2-6x+8\) a \( g(x)=-x^2-2x+24 \) určte množinu riešení danej nerovnice.
\[ x^2-6x+8\leq-x^2-2x+24 \]}
{\( \langle-2;4\rangle \)}{\( \langle2;4\rangle \)}{\( \langle0;24\rangle \)}{\( \langle-6;4\rangle \)}{}{}}
\LANGes{
\Question{
Dadas las gráficas de las funciones \( f(x)=x^2-6x+8\) y \( g(x)=-x^2-2x+24 \), encuentra el conjunto solución de la siguiente inecuación.
\[ x^2-6x+8\leq-x^2-2x+24 \]}
{\( [-2;4] \)}{\( [2;4] \)}{\( [0;24] \)}{\( [-6;4] \)}{}{}}
\End

\Data\ID{1103034602}
\Updated{2021-02-15 07:02:02}
\Level{B}
\SubArea{Quadratic Functions}
\IMG{1103034602_0.pdf}
\LANGen{
\Question{
Given graphs of the functions \( f(x)=x^2-6x+8\) and \( g(x)=2x+1 \), find the solution set of the following inequality.
\[ x^2-6x+8>2x+1 \]}
{\( (-\infty;1)\cup(7;\infty) \)}{\( (1;7) \)}{\( (-\infty;2)\cup(4;\infty) \)}{\( (1;\infty) \)}{}{}}
\LANGpl{
\Question{
Dane są wykresy funkcji  \( f(x)=x^2-6x+8\) i \( g(x)=2x+1 \), określ zbiór rozwiązań następującej nierówności.
\[ x^2-6x+8>2x+1 \]}
{\( (-\infty;1)\cup(7;\infty) \)}{\( (1;7) \)}{\( (-\infty;2)\cup(4;\infty) \)}{\( (1;\infty) \)}{}{}}
\LANGcs{
\Question{
Pomocí grafů funkcí \( f(x)=x^2-6x+8\) a \( g(x)=2x+1 \) určete množinu řešení dané nerovnice.
\[ x^2-6x+8>2x+1 \]}
{\( (-\infty;1)\cup(7;\infty) \)}{\( (1;7) \)}{\( (-\infty;2)\cup(4;\infty) \)}{\( (1;\infty) \)}{}{}}
\LANGsk{
\Question{
Pomocou grafov funkcií \( f(x)=x^2-6x+8\) a \( g(x)=2x+1 \) určte množinu riešení danej nerovnice.
\[ x^2-6x+8>2x+1 \]}
{\( (-\infty;1)\cup(7;\infty) \)}{\( (1;7) \)}{\( (-\infty;2)\cup(4;\infty) \)}{\( (1;\infty) \)}{}{}}
\LANGes{
\Question{
Dadas las gráficas de las funciones \( f(x)=x^2-6x+8\) y \( g(x)=2x+1 \), encuentra el conjunto solución de la siguiente inecuación.
\[ x^2-6x+8>2x+1 \]}
{\( (-\infty;1)\cup(7;\infty) \)}{\( (1;7) \)}{\( (-\infty;2)\cup(4;\infty) \)}{\( (1;\infty) \)}{}{}}
\End

\Data\ID{1103034603}
\Updated{2021-02-15 07:02:23}
\Level{B}
\SubArea{Quadratic Functions}
\IMG{1103034603.pdf}
\LANGen{
\Question{
Given graphs of the functions \( f(x)=x^2-6x+8\) and \( g(x)=-2x+4 \), find the solution set of the following inequality.
\[ x^2-6x+8\geq-2x+4 \]}
{\( \mathbb{R} \)}{\( (-\infty;2]\cup[4;\infty) \)}{\( \{2\} \)}{\( [2;4] \)}{}{}}
\LANGpl{
\Question{
Dane są wykresy funkcji  \( f(x)=x^2-6x+8\) i \( g(x)=-2x+4 \), określ zbiór rozwiązań następującej nierówności.
\[ x^2-6x+8\geq-2x+4 \]}
{\( \mathbb{R} \)}{\( (-\infty;2\rangle\cup\langle4;\infty) \)}{\( \{2\} \)}{\( \langle2;4\rangle \)}{}{}}
\LANGcs{
\Question{
Pomocí grafů funkcí \( f(x)=x^2-6x+8\) a \( g(x)=-2x+4 \) určete množinu řešení dané nerovnice.
\[ x^2-6x+8\geq-2x+4 \]}
{\( \mathbb{R} \)}{\( (-\infty;2\rangle\cup\langle4;\infty) \)}{\( \{2\} \)}{\( \langle2;4\rangle \)}{}{}}
\LANGsk{
\Question{
Pomocou grafov funkcií \( f(x)=x^2-6x+8\) a \( g(x)=-2x+4 \) určte množinu riešení danej nerovnice.
\[ x^2-6x+8\geq-2x+4 \]}
{\( \mathbb{R} \)}{\( (-\infty;2\rangle\cup\langle4;\infty) \)}{\( \{2\} \)}{\( \langle2;4\rangle \)}{}{}}
\LANGes{
\Question{
Dadas las gráficas de las funciones \( f(x)=x^2-6x+8\) y \( g(x)=-2x+4 \), encuentra el conjunto solución de la siguiente inecuación.
\[ x^2-6x+8\geq-2x+4 \]}
{\( \mathbb{R} \)}{\( (-\infty;2]\cup[4;\infty) \)}{\( \{2\} \)}{\( [2;4] \)}{}{}}
\End

\Data\ID{1103034604}
\Updated{2021-02-15 07:02:06}
\Level{B}
\SubArea{Quadratic Functions}
\IMG{1103034604_0.pdf}
\LANGen{
\Question{
Given graphs of the functions \( f(x)=-2x^2+5x+3\) and \( g(x)=2x+1 \), find the solution set of the following inequality.
\[ -2x^2+5x+3 < 2x+1 \]}
{\( (-\infty;-0.5)\cup(2;\infty) \)}{\( (-0.5;2) \)}{\( (-\infty;-0.5)\cup(3;\infty) \)}{\( (-0.5;3) \)}{}{}}
\LANGpl{
\Question{
Dane są wykresy funkcji  \( f(x)=-2x^2+5x+3\) i \( g(x)=2x+1 \), określ zbiór rozwiązań następującej nierówności.
\[ -2x^2+5x+3 < 2x+1 \]}
{\( (-\infty;-0{,}5)\cup(2;\infty) \)}{\( (-0{,}5;2) \)}{\( (-\infty;-0{,}5)\cup(3;\infty) \)}{\( (-0{,}5;3) \)}{}{}}
\LANGcs{
\Question{
Pomocí grafů funkcí \( f(x)=-2x^2+5x+3\) a \( g(x)=2x+1 \) určete množinu řešení dané nerovnice.
\[ -2x^2+5x+3 < 2x+1 \]}
{\( (-\infty;-0{,}5)\cup(2;\infty) \)}{\( (-0{,}5;2) \)}{\( (-\infty;-0{,}5)\cup(3;\infty) \)}{\( (-0{,}5;3) \)}{}{}}
\LANGsk{
\Question{
Pomocou grafov funkcií \( f(x)=-2x^2+5x+3\) a \( g(x)=2x+1 \) určte množinu riešení danej nerovnice.
\[ -2x^2+5x+3 < 2x+1 \]}
{\( (-\infty;-0{,}5)\cup(2;\infty) \)}{\( (-0{,}5;2) \)}{\( (-\infty;-0{,}5)\cup(3;\infty) \)}{\( (-0{,}5;3) \)}{}{}}
\LANGes{
\Question{
Dadas las gráficas de las funciones \( f(x)=-2x^2+5x+3\) y \( g(x)=2x+1 \), encuentra el conjunto solución de la siguiente inecuación.
\[ -2x^2+5x+3 < 2x+1 \]}
{\( (-\infty;-0.5)\cup(2;\infty) \)}{\( (-0.5;2) \)}{\( (-\infty;-0.5)\cup(3;\infty) \)}{\( (-0.5;3) \)}{}{}}
\End

\Data\ID{1103034605}
\Updated{2021-04-07 05:04:53}
\Level{B}
\SubArea{Quadratic Functions}
\IMG{1103034605_0.pdf}
\LANGen{
\Question{
Given graphs of the functions \( f(x)=x^2+2x-3\) and \( g(x)=-x^2+3x-4 \), find the solution set of the following inequality.
\[ x^2+2x-3\leq-x^2+3x-4  \]}
{\( \emptyset \)}{\( \mathbb{R} \)}{\( (-\infty;0] \)}{\( [-3;1] \)}{}{}}
\LANGpl{
\Question{
Dane są wykresy funkcji \( f(x)=x^2+2x-3\) i \( g(x)=-x^2+3x-4 \), określ zbiór rozwiązań następującej nierówności.
\[ x^2+2x-3\leq-x^2+3x-4  \]}
{\( \emptyset \)}{\( \mathbb{R} \)}{\( (-\infty;0\rangle \)}{\( \langle-3;1\rangle \)}{}{}}
\LANGcs{
\Question{
Pomocí grafů funkcí \( f(x)=x^2+2x-3\) a \( g(x)=-x^2+3x-4 \) určete množinu řešení dané nerovnice.
\[ x^2+2x-3\leq-x^2+3x-4  \]}
{\( \emptyset \)}{\( \mathbb{R} \)}{\( (-\infty;0\rangle \)}{\( \langle-3;1\rangle \)}{}{}}
\LANGsk{
\Question{
Pomocou grafov funkcií \( f(x)=x^2+2x-3\) a \( g(x)=-x^2+3x-4 \) určte množinu riešení danej nerovnice.
\[ x^2+2x-3\leq-x^2+3x-4  \]}
{\( \emptyset \)}{\( \mathbb{R} \)}{\( (-\infty;0\rangle \)}{\( \langle-3;1\rangle \)}{}{}}
\LANGes{
\Question{
Dadas las gráficas de las funciones \( f(x)=x^2+2x-3\) y \( g(x)=-x^2+3x-4 \), encuentra el conjunto solución de la siguiente inecuación.
\[ x^2+2x-3\leq-x^2+3x-4  \]}
{\( \emptyset \)}{\( \mathbb{R} \)}{\( (-\infty;0] \)}{\( [-3;1] \)}{}{}}
\End

\Data\ID{1103034606}
\Updated{2021-04-07 05:04:20}
\Level{B}
\SubArea{Quadratic Functions}
\IMG{1103034606_0.pdf}
\LANGen{
\Question{
Given graphs of the functions \( f(x)=x^2-4x\) and \( g(x)=4x^2-16x+12 \), find the solution set of the following inequality.
\[ 4x^2-16x+12\leq x^2-4x \]}
{\( \{2\} \)}{\( \mathbb{R} \)}{\( \{-4\} \)}{\( [1;3] \)}{}{}}
\LANGpl{
\Question{
Dane są wykresy funkcji \( f(x)=x^2-4x\) i \( g(x)=4x^2-16x+12 \), określ zbiór rozwiązań następującej nierówności.
\[ 4x^2-16x+12\leq x^2-4x \]}
{\( \{2\} \)}{\( \mathbb{R} \)}{\( \{-4\} \)}{\( \langle1;3\rangle \)}{}{}}
\LANGcs{
\Question{
Pomocí grafů funkcí \( f(x)=x^2-4x\) a \( g(x)=4x^2-16x+12 \) určete množinu řešení dané nerovnice.
\[ 4x^2-16x+12\leq x^2-4x \]}
{\( \{2\} \)}{\( \mathbb{R} \)}{\( \{-4\} \)}{\( \langle1;3\rangle \)}{}{}}
\LANGsk{
\Question{
Pomocou grafov funkcií \( f(x)=x^2-4x\) a \( g(x)=4x^2-16x+12 \) určte množinu riešení danej nerovnice.
\[ 4x^2-16x+12\leq x^2-4x \]}
{\( \{2\} \)}{\( \mathbb{R} \)}{\( \{-4\} \)}{\( \langle1;3\rangle \)}{}{}}
\LANGes{
\Question{
Dadas las gráficas de las funciones \( f(x)=x^2-4x\) y \( g(x)=4x^2-16x+12 \), encuentra el conjunto solución de la siguiente inecuación.
\[ 4x^2-16x+12\leq x^2-4x \]}
{\( \{2\} \)}{\( \mathbb{R} \)}{\( \{-4\} \)}{\( [1;3] \)}{}{}}
\End

\Data\ID{1103034701}
\Updated{2021-04-07 05:04:47}
\Level{B}
\SubArea{Quadratic Functions}
\IMG{1103034701_0.pdf}
\LANGen{
\Question{
Given graphs of the functions \( f(x)=2x^2-2x-4 \) and \( g(x)=2x+2 \), find the solution set of the following equation.
\[ 2x^2-2x-4=2x+2 \]}
{\( \{-1;3\} \)}{\( \{-1;2\} \)}{\( \{0;8\} \)}{\( \{-4;0\} \)}{}{}}
\LANGpl{
\Question{
Dane są wykresy funkcji \( f(x)=2x^2-2x-4 \) i \( g(x)=2x+2 \), określ zbiór rozwiązań następującego równania.
\[ 2x^2-2x-4=2x+2 \]}
{\( \{-1;3\} \)}{\( \{-1;2\} \)}{\( \{0;8\} \)}{\( \{-4;0\} \)}{}{}}
\LANGcs{
\Question{
Pomocí grafů funkcí \( f(x)=2x^2-2x-4 \) a \( g(x)=2x+2 \) určete množinu řešení dané rovnice.
\[ 2x^2-2x-4=2x+2 \]}
{\( \{-1;3\} \)}{\( \{-1;2\} \)}{\( \{0;8\} \)}{\( \{-4;0\} \)}{}{}}
\LANGsk{
\Question{
Pomocou grafov funkcií \( f(x)=2x^2-2x-4 \) a \( g(x)=2x+2 \) určte množinu riešení danej rovnice.
\[ 2x^2-2x-4=2x+2 \]}
{\( \{-1;3\} \)}{\( \{-1;2\} \)}{\( \{0;8\} \)}{\( \{-4;0\} \)}{}{}}
\LANGes{
\Question{
Dadas las gráficas de las funciones \( f(x)=2x^2-2x-4 \) y \( g(x)=2x+2 \), encuentra el conjunto  solución de la siguiente ecuación.
\[ 2x^2-2x-4=2x+2 \]}
{\( \{-1;3\} \)}{\( \{-1;2\} \)}{\( \{0;8\} \)}{\( \{-4;0\} \)}{}{}}
\End

\Data\ID{1103034702}
\Updated{2021-04-07 05:04:01}
\Level{B}
\SubArea{Quadratic Functions}
\IMG{1103034702_0.pdf}
\LANGen{
\Question{
Given graphs of the functions \( f(x)=x^2-4x \) and \( g(x)=4x^2-16x+12 \), find the solution set of the following equation.
\[ 4x^2-16x+12=x^2-4x \]}
{\( \{2\} \)}{\( \{-4\} \)}{\( \{0;4\} \)}{\( \{0\} \)}{}{}}
\LANGpl{
\Question{
Dane są wykresy funkcji \( f(x)=x^2-4x \) i \( g(x)=4x^2-16x+12 \), określ zbiór rozwiązań następującego równania.
\[ 4x^2-16x+12=x^2-4x \]}
{\( \{2\} \)}{\( \{-4\} \)}{\( \{0;4\} \)}{\( \{0\} \)}{}{}}
\LANGcs{
\Question{
Pomocí grafů funkcí \( f(x)=x^2-4x \) a \( g(x)=4x^2-16x+12 \) určete množinu řešení dané rovnice.
\[ 4x^2-16x+12=x^2-4x \]}
{\( \{2\} \)}{\( \{-4\} \)}{\( \{0;4\} \)}{\( \{0\} \)}{}{}}
\LANGsk{
\Question{
Pomocou grafov funkcií \( f(x)=x^2-4x \) a \( g(x)=4x^2-16x+12 \) určte množinu riešení danej rovnice.
\[ 4x^2-16x+12=x^2-4x \]}
{\( \{2\} \)}{\( \{-4\} \)}{\( \{0;4\} \)}{\( \{0\} \)}{}{}}
\LANGes{
\Question{
dadas las gráficas de las funciones \( f(x)=x^2-4x \) y \( g(x)=4x^2-16x+12 \), encuentra el conjunto  solución  la siguiente ecuación.
\[ 4x^2-16x+12=x^2-4x \]}
{\( \{2\} \)}{\( \{-4\} \)}{\( \{0;4\} \)}{\( \{0\} \)}{}{}}
\End

\Data\ID{1103034703}
\Updated{2021-04-07 05:04:14}
\Level{B}
\SubArea{Quadratic Functions}
\IMG{1103034703_0.pdf}
\LANGen{
\Question{
Given the graph of the function \( f(x)=2x^2-2x-4 \) and the points \( A \), \( B \), \( C \), \( D \), \( E \), find the solution set of the following equation.
\[ 2x^2-2x-4=1 \]}
{\( \{a;e\} \)}{\( \{b;d\} \)}{\( \{c;e\} \)}{\( \{ c \} \)}{}{}}
\LANGpl{
\Question{
Dany jest wykres funkcji \( f(x)=2x^2-2x-4 \) i punkty \( A \), \( B \), \( C \), \( D \), \( E \), określ zbiór rozwiązań następującego równania.
\[ 2x^2-2x-4=1 \]}
{\( \{a;e\} \)}{\( \{b;d\} \)}{\( \{c;e\} \)}{\( \{ c \} \)}{}{}}
\LANGcs{
\Question{
Pomocí grafu funkce \( f(x)=2x^2-2x-4 \) a bodů \( A \), \( B \), \( C \), \( D \), \( E \) určete množinu řešení dané rovnice.
\[ 2x^2-2x-4=1 \]}
{\( \{a;e\} \)}{\( \{b;d\} \)}{\( \{c;e\} \)}{\( \{ c \} \)}{}{}}
\LANGsk{
\Question{
Pomocou grafu funkcie \( f(x)=2x^2-2x-4 \) a bodov \( A \), \( B \), \( C \), \( D \), \( E \) určte množinu riešení danej rovnice.
\[ 2x^2-2x-4=1 \]}
{\( \{a;e\} \)}{\( \{b;d\} \)}{\( \{c;e\} \)}{\( \{ c \} \)}{}{}}
\LANGes{
\Question{
Dada la gráfica de la función \( f(x)=2x^2-2x-4 \) y los puntos \( A \), \( B \), \( C \), \( D \), \( E \), encuentra el conjunto solución de la siguiente ecuación.
\[ 2x^2-2x-4=1 \]}
{\( \{a;e\} \)}{\( \{b;d\} \)}{\( \{c;e\} \)}{\( \{ c \} \)}{}{}}
\End

\Data\ID{1103034704}
\Updated{2021-04-07 05:04:29}
\Level{B}
\SubArea{Quadratic Functions}
\IMG{1103034704_0.pdf}
\LANGen{
\Question{
Given graphs of the functions \( f(x)=-x^2+4\) and \( g(x)=x+2 \), find the solution set of the following equation.
\[ -x^2+4=x+2 \]}
{\( \{-2;1\} \)}{\( \{0;3\} \)}{\( \{-2;2\} \)}{\( \{0;4\} \)}{}{}}
\LANGpl{
\Question{
Dane są wykresy funkcji \( f(x)=-x^2+4\) i \( g(x)=x+2 \), określ zbiór rozwiązań następującego równania.
\[ -x^2+4=x+2 \]}
{\( \{-2;1\} \)}{\( \{0;3\} \)}{\( \{-2;2\} \)}{\( \{0;4\} \)}{}{}}
\LANGcs{
\Question{
Pomocí grafů funkcí \( f(x)=-x^2+4\) a \( g(x)=x+2 \) určete množinu řešení dané rovnice.
\[ -x^2+4=x+2 \]}
{\( \{-2;1\} \)}{\( \{0;3\} \)}{\( \{-2;2\} \)}{\( \{0;4\} \)}{}{}}
\LANGsk{
\Question{
Pomocou grafov funkcií \( f(x)=-x^2+4\) a \( g(x)=x+2 \) určte množinu riešení danej rovnice.
\[ -x^2+4=x+2 \]}
{\( \{-2;1\} \)}{\( \{0;3\} \)}{\( \{-2;2\} \)}{\( \{0;4\} \)}{}{}}
\LANGes{
\Question{
Dadas las gráficas de las funciones \( f(x)=-x^2+4\) y \( g(x)=x+2 \), encuentra el conjunto solución de la siguiente ecuación.
\[ -x^2+4=x+2 \]}
{\( \{-2;1\} \)}{\( \{0;3\} \)}{\( \{-2;2\} \)}{\( \{0;4\} \)}{}{}}
\End

\Data\ID{1103034705}
\Updated{2021-04-07 06:04:28}
\Level{B}
\SubArea{Quadratic Functions}
\IMG{1103034705_0.pdf}
\LANGen{
\Question{
Given graphs of the functions \( f(x)=-x^2-x+6 \) and \( g(x)=x^2-4x+4 \), find the solution set of the following equation.
\[ x^2-4x+4 = -x^2-x+6 \]}
{\( \{-0.5;2\} \)}{\( \{2\} \)}{\( \{-3;2\} \)}{\( \{0;6.25\} \)}{}{}}
\LANGpl{
\Question{
Dane są wykresy funkcji \( f(x)=-x^2-x+6 \) i \( g(x)=x^2-4x+4 \), określ zbiór rozwiązań następującego równania.
\[ x^2-4x+4 = -x^2-x+6 \]}
{\( \{-0{,}5;2\} \)}{\( \{2\} \)}{\( \{-3;2\} \)}{\( \{0;6{,}25\} \)}{}{}}
\LANGcs{
\Question{
Pomocí grafů funkcí \( f(x)=-x^2-x+6 \) a \( g(x)=x^2-4x+4 \) určete množinu řešení dané rovnice.
\[ x^2-4x+4 = -x^2-x+6 \]}
{\( \{-0{,}5;2\} \)}{\( \{2\} \)}{\( \{-3;2\} \)}{\( \{0;6{,}25\} \)}{}{}}
\LANGsk{
\Question{
Pomocou grafov funkcií \( f(x)=-x^2-x+6 \) a \( g(x)=x^2-4x+4 \) určte množinu riešení danej rovnice.
\[ x^2-4x+4 = -x^2-x+6 \]}
{\( \{-0{,}5;2\} \)}{\( \{2\} \)}{\( \{-3;2\} \)}{\( \{0;6{,}25\} \)}{}{}}
\LANGes{
\Question{
Dadas las gráficas de las funciones \( f(x)=-x^2-x+6 \) y \( g(x)=x^2-4x+4 \), encuentra el conjunto solución de la siguiente ecuación.
\[ x^2-4x+4 = -x^2-x+6 \]}
{\( \{-0.5;2\} \)}{\( \{2\} \)}{\( \{-3;2\} \)}{\( \{0;6.25\} \)}{}{}}
\End

\Data\ID{1103034706}
\Updated{2021-04-07 06:04:53}
\Level{B}
\SubArea{Quadratic Functions}
\IMG{1103034706_0.pdf}
\LANGen{
\Question{
Given graphs of the functions \( f(x)=-x^2-x+6 \) and \( g(x)=x^2-4x+4 \), find the solution set of the following equation.
\[ -x^2-x+6=0 \]}
{\( \{-3;2\} \)}{\( \{-0.5;2\} \)}{\( \{2\} \)}{\( \{0;6.25\} \)}{}{}}
\LANGpl{
\Question{
Dane są wykresy funkcji \( f(x)=-x^2-x+6 \) i \( g(x)=x^2-4x+4 \), określ zbiór rozwiązań następującego równania.
\[ -x^2-x+6=0 \]}
{\( \{-3;2\} \)}{\( \{-0{,}5;2\} \)}{\( \{2\} \)}{\( \{0;6{,}25\} \)}{}{}}
\LANGcs{
\Question{
Pomocí grafů funkcí \( f(x)=-x^2-x+6 \) a \( g(x)=x^2-4x+4 \) určete množinu řešení dané rovnice.
\[ -x^2-x+6=0 \]}
{\( \{-3;2\} \)}{\( \{-0{,}5;2\} \)}{\( \{2\} \)}{\( \{0;6{,}25\} \)}{}{}}
\LANGsk{
\Question{
Pomocou grafov funkcií \( f(x)=-x^2-x+6 \) a \( g(x)=x^2-4x+4 \) určte množinu riešení danej rovnice.
\[ -x^2-x+6=0 \]}
{\( \{-3;2\} \)}{\( \{-0{,}5;2\} \)}{\( \{2\} \)}{\( \{0;6{,}25\} \)}{}{}}
\LANGes{
\Question{
Dadas las gráficas de las funciones \( f(x)=-x^2-x+6 \) y \( g(x)=x^2-4x+4 \), encuentra el conjunto de solución la siguiente equación.
\[ -x^2-x+6=0 \]}
{\( \{-3;2\} \)}{\( \{-0.5;2\} \)}{\( \{2\} \)}{\( \{0;6.25\} \)}{}{}}
\End

\Data\ID{1103083107}
\Updated{2021-04-07 07:04:56}
\Level{B}
\SubArea{Quadratic Functions}
\IMG{1103083107.pdf}
\LANGen{
\Question{
The quadratic functions \( f \) and \( g \) that have the same vertex \( V \) are graphed in the picture. The graph of \( g \) is the reflection of the graph of \( f \) in  the vertex \( V \). Also, both the graphs are symmetric across \( y \)-axis. Identify the true statement about \( f \) and \( g \).}
{The equations of \( f \) and \( g \) differ in the sign of the coefficient at the quadratic term only.}{The equations of \( f \) and \( g \) differ in the sign of the coefficient at the linear term only.}{The equations of \( f \) and \( g \) differ in the sign of the coefficient at the absolute term only.}{None of the statements above is true.}{}{}}
\LANGpl{
\Question{
Funkcje kwadratowe \( f \) i \( g \), które mają ten sam wierzchołek \( V \) przedstawiono na rysunku poniżej. Wykres funkcji \( g \) jest odbiciem wykresu funkcji \( f \) w wierzchołku \( V \). Obydwa wykresy są symetryczne względem osi \( y \). Wybierz prawdziwe stwierdzenie dotyczące funkcji \( f \) i  \( g \).}
{Wzory funkcji \( f \) i \( g \) różnią się znakiem współczynnika tylko w stosunku kwadratowym.}{Wzory funkcji  \( f \) i \( g \) różnią się znakiem współczynnika tylko w stosunku liniowym.}{Wzory funkcji  \( f \) i \( g \) różnią się znakiem współczynnika tylko w stosunku bezwzględnym.}{Żadne z powyższych stwierdzeń nie jest prawdziwe.}{}{}}
\LANGcs{
\Question{
Na obrázku jsou grafy kvadratických funkcí \( f \) a \( g \) -- paraboly se společným vrcholem \( V \). Graf funkce \( g \) je obrazem grafu funkce \( f \) ve středové souměrnosti dané bodem \( V \) a současně jsou oba grafy osově souměrné podle osy \( y \). Vyberte pravdivé tvrzení o předpisech těchto funkcí.}
{Předpisy funkcí \( f \) a \( g \) se liší pouze znaménkem koeficientu u kvadratického členu.}{Předpisy funkcí \( f \) a \( g \) se liší pouze znaménkem koeficientu u lineárního členu.}{Předpisy funkcí \( f \) a \( g \) se liší pouze znaménkem koeficientu u absolutního členu.}{Neplatí žádné z výše uvedených tvrzení.}{}{}}
\LANGsk{
\Question{
Na obrázku sú grafy kvadratických funkcií \( f \) a \( g \), ktoré majú spoločný vrchol \( V \). Grafy funkcií \( f \) a \( g \) sú stredovo súmerné podľa vrcholu \( V \) a súčasne obidve funkcie sú osovo súmerné podľa osi \( y \).  Vyberte pravdivé tvrdenie o predpisoch funkcii \( f \) a \( g \).}
{Predpisy funkcií \( f \) a \( g \) sa líšia len znamienkom koeficientu kvadratického člena.}{Predpisy funkcií \( f \) a \( g \) sa líšia len znamienkom koeficientu lineárneho člena.}{Predpisy funkcií \( f \) a \( g \) sa líšia len znamienkom absolútneho člena.}{Žiadne tvrdenie nie je pravdivé.}{}{}}
\LANGes{
\Question{
Las gráficas de las  funciones cuadráticas \( f \) y \( g \) que tienen el mismo vértice \( V \) aparecen en el dibujo. La gráfica de \( g \) es la reflexión de la gráfica de \( f \) respecto al vértice \( V \). Además ambas gráficas son simétricas respecto el eje \( y \). Identifica la declaración correcta sobre \( f \) y \( g \).}
{Las ecuaciones de \( f \) y \( g \) difieren solamente en el signo del coeficiente del término cuadrático.}{Las ecuaciones de \( f \) y \( g \) difieren solamente en el signo del coeficiente del término lineal.}{Las ecuaciones de \( f \) y \( g \) difieren solamente en el signo del coeficiente del término absoluto.}{Ninguna de las declaraciones es correcta.}{}{}}
\End

\Data\ID{1103083109}
\Updated{2021-07-24 10:07:47}
\Level{B}
\SubArea{Quadratic Functions}
\IMG{1103083109.pdf}
\LANGen{
\Question{
The graphs of the quadratic functions \( f \) and \( g \) are shown in the picture. The graph of \( g \) is the reflection of the graph of \( f \) about \( y \)-axis. Identify which of the following statements about \( f \) and \( g \) is true.}
{The equations of \( f \) and \( g \) differ in the sign of the coefficient at the linear term only.}{The equations of \( f \) and \( g \) differ in the sign of the coefficient at the quadratic term only.}{The equations of \( f \) and \( g \) differ in in the sign of the coefficient at the absolute term  only.}{None of the statements above is true.}{}{}}
\LANGpl{
\Question{
Wykresy funkcji kwadratowych  \( f \) i \( g \) przedstawiono na rysunku poniżej.  Wykres funkcji  \( g \) jest odbiciem wykresu funkcji \( f \) względem osi \( y \). Określ, które z poniższych stwierdzeń dotyczący wykresów funkcji \( f \) i \( g \) jest prawdziwe.}
{Wzory funkcji \( f \) i \( g \) różnią się znakiem współczynnika tylko w stosunku liniowym.}{Wzory funkcji  \( f \) i \( g \) różnią się znakiem współczynnika tylko w stosunku kwadratowym.}{Wzory funkcji \( f \) i \( g \) różnią się znakiem współczynnika tylko w stosunku bezwzględnym.}{Żadne z powyższych stwierdzeń nie jest prawdziwe.}{}{}}
\LANGcs{
\Question{
Na obrázku jsou grafy kvadratických funkcí \( f \) a \( g \) - paraboly s různými vrcholy. Oba grafy jsou osově souměrné podle osy \( y \). Vyberte pravdivé tvrzení o předpisech těchto funkcí.}
{Předpisy funkcí \( f \) a \( g \) se liší pouze znaménkem koeficientu u lineárního členu.}{Předpisy funkcí \( f \) a \( g \) se liší pouze znaménkem koeficientu u kvadratického členu.}{Předpisy funkcí \( f \) a \( g \) se liší pouze znaménkem koeficientu u absolutního členu.}{Neplatí žádné z výše uvedených tvrzení.}{}{}}
\LANGsk{
\Question{
Na obrázku sú grafy kvadratických funkcií \( f \) a \( g \). Grafy funkcií \( f \) a \( g \) sú osovo súmerné podľa osi \( y \). Vyberte pravdivé tvrdenie o predpisoch funkcii \( f \) a \( g \).}
{Predpisy funkcií \( f \) a \( g \) sa líšia len znamienkom koeficientu lineárneho člena.}{Predpisy funkcií \( f \) a \( g \) sa líšia len znamienkom koeficientu kvadratického člena.}{Predpisy funkcií \( f \) a \( g \) sa líšia len znamienkom absolútneho člena.}{Žiadne tvrdenie nie je pravdivé.}{}{}}
\LANGes{
\Question{
En el dibujo aparecen las gráficas de las funciones cuadráticas \( f \) y \( g \) . La gráfica de \( g \) es la reflexión de la gráfica de \( f \) respecto al eje \( y \). Identifica cuál de las declaraciones sobre \( f \) y \( g \) es correcta.}
{Las ecuaciones de \( f \) y \( g \) difieren solamente en el signo del coeficiente del término lineal.}{Las ecuaciones de \( f \) y \( g \) difieren solamente en el signo del coeficiente del término cuadrático.}{Las ecuaciones de \( f \) y \( g \) difieren solamente en el signo del coeficiente del término absoluto.}{Ninguna de las declaraciones es correcta.}{}{}}
\End

\Data\ID{1103120001}
\Updated{2021-04-07 07:04:09}
\Level{B}
\SubArea{Quadratic Functions}
\IMG{1103120001a_0.pdf}
\LANGen{
\Question{
In the picture A,  we are given the graph of the quadratic function \( f(x)=\frac12x^2 \).  Use the graph of \( f \) as help to identify which of the graphs given in the picture B is the graph of \( g(x) =\frac12 x^2-2 \).
Choose what is the colour of the graph of \( g \). (Note: The graphs in the picture B were obtained  by shifting  the graph of \( f \).)}
{blue}{green}{red}{yellow}{}{}}
\LANGpl{
\Question{
Na rysunku A,  przedstawiono wykres funkcji kwadratowej \( f(x)=\frac12x^2 \).  Używając wykresu funkcji \( f \) określ, który z wykresów przedstawionych na rysunku B jest wykresem funkcji \( g(x) =\frac12 x^2-2 \).
Wybierz kolor wykresu funkcji \( g \). (Wskazówka: Wykresy na rysunku  B uzyskano poprzez przesunięcie wykresu \( f \).)}
{niebieski}{zielony}{czerwony}{żółty}{}{}}
\LANGcs{
\Question{
Na obrázku A je nakreslen graf kvadratické funkce \( f(x)=\frac12x^2 \). Grafy na obrázku B vznikly posunutím grafu funkce \( f \). Určete barvu grafu funkce \( g(x) =\frac12 x^2-2 \).}
{modrá}{zelená}{červená}{žlutá}{}{}}
\LANGsk{
\Question{
Na obrázku A je nakreslený graf funkcie \( f(x)=\frac12x^2 \). Grafy na obrázku B vznikli posunutím grafu funkcie \( f \). Určte farbu grafu funkcie \( g(x) =\frac12 x^2-2 \).}
{modrá}{zelená}{červená}{žltá}{}{}}
\LANGes{
\Question{
En el dibujo A, tenemos la gráfica de la función cuadrática \( f(x)=\frac12x^2 \).  Usa la gráfica de \( f \) como ayuda para identificar cuál de las gráficas en el dibujo B es la gráfica de \( g(x) =\frac12 x^2-2 \). Elije cuál es el color de la gráfica de \( g \). (Nota: Las gráficas en el dibujo B fueron obtenidas trasladando la gráfica de \( f \).)}
{azul}{verde}{rojo}{amarillo}{}{}}
\End

\Data\ID{1103120002}
\Updated{2022-04-23 05:04:17}
\Level{B}
\SubArea{Quadratic Functions}
\IMG{1103120002.pdf}
\LANGen{
\Question{
Let \( f(x)=2x^2 \). Given the graph of the function \( f \) and the graph of a function \( g \) which was obtained as a right shift  of the graph of \( f \) (see the picture), choose the function \( g \).}
{\( g(x) = 2(x-3)^2 \)}{\( g(x) = 2(x+3)^2 \)}{\( g(x) = 2x^2+3 \)}{\( g(x) = 2x^2-3 \)}{}{}}
\LANGpl{
\Question{
Niech  \( f(x)=2x^2 \). Dany jest wykres funkcji  \( f \) i wykres funkcji \( g \), który powstał poprzez przesuniecie w prawo wykresu funkcji \( f \) (rysunek poniżej), wybierz funkcję  \( g \).}
{\( g(x) = 2(x-3)^2 \)}{\( g(x) = 2(x+3)^2 \)}{\( g(x) = 2x^2+3 \)}{\( g(x) = 2x^2-3 \)}{}{}}
\LANGcs{
\Question{
Na obrázku je nakreslen graf funkce \( f(x)=2x^2 \) a graf funkce \( g \), který vznikl posunem grafu funkce \( f \). Vyberte předpis funkce \( g \).}
{\( g(x) = 2(x-3)^2 \)}{\( g(x) = 2(x+3)^2 \)}{\( g(x) = 2x^2+3 \)}{\( g(x) = 2x^2-3 \)}{}{}}
\LANGsk{
\Question{
Na obrázku je nakreslený graf funkcie \( f(x)=2x^2 \) a graf funkcie \( g \), ktorý vznikol posunutím grafu funkcie \( f \).  Vyberte predpis funkcie \( g \).}
{\( g(x) = 2(x-3)^2 \)}{\( g(x) = 2(x+3)^2 \)}{\( g(x) = 2x^2+3 \)}{\( g(x) = 2x^2-3 \)}{}{}}
\LANGes{
\Question{
Sea \( f(x)=2x^2 \). Dadas las gráficas de la función \( f \) y de la función \( g \) que fue obtenida moviendo la gráfica de \( f \) a la derecha (mira el dibujo), elije la función \( g \).}
{\( g(x) = 2(x-3)^2 \)}{\( g(x) = 2(x+3)^2 \)}{\( g(x) = 2x^2+3 \)}{\( g(x) = 2x^2-3 \)}{}{}}
\End

\Data\ID{1103120003}
\Updated{2021-04-07 07:04:59}
\Level{B}
\SubArea{Quadratic Functions}
\IMG{1103120003a.pdf}
\LANGen{
\Question{
In the picture A,  we are given the graph of the quadratic function \( f(x)=-2x^2 \).  Use the graph of \( f \) as help to identify which of the graphs given in the picture B is the graph of  \( g(x)=-2(x+4)^2 \). Choose what is the colour of the graph of \( g \). (Note: The graphs in the picture B were obtained  by  shifting  the graph of \(  f \).)}
{red}{blue}{green}{yellow}{}{}}
\LANGpl{
\Question{
Na rysunku A,  przedstawiono wykres funkcji kwadratowej \( f(x)=-2x^2 \).  Użyj wykresu funkcji \( f \), aby określić, który z wykresów przedstawionych na rysunku B jest wykresem funkcji  \( g(x)=-2(x+4)^2 \).  Wybierz kolor przedstawiający wykres funkcji \( g \). (Wskazówka: Wykresy na rysunku B powstały w wyniku przesunięcia wykresu funkcji \(  f \).)}
{czerwony}{niebieski}{zielony}{żółty}{}{}}
\LANGcs{
\Question{
Na obrázku A je nakreslen graf kvadratické funkce \( f(x)=-2x^2 \). Grafy na obrázku B vznikly posunutím grafu funkce \( f \). Určete barvu grafu funkce \( g(x)=-2(x+4)^2 \).}
{červená}{modrá}{zelená}{žlutá}{}{}}
\LANGsk{
\Question{
Na obrázku A je nakreslený graf funkcie \( f(x)=-2x^2 \).  Grafy na obrázku B vznikli posunutím grafu funkcie \( f \).  Určte farbu grafu funkcie \( g(x)=-2(x+4)^2 \).}
{červená}{modrá}{zelená}{žltá}{}{}}
\LANGes{
\Question{
En el dibujo A tenemos la gráfica de la función cuadrática \( f(x)=-2x^2 \).  Usa la gráfica de \( f \) como ayuda para determinar cuál de las gráficas en el dibujo B es la de \( g(x)=-2(x+4)^2 \). Elije qué color tiene la gráfica de la función \( g \). (Nota: Las gráficas en el dibujo B fueron obtenidas moviendo la gráfica de \(  f \).)}
{rojo}{azul}{verde}{amarillo}{}{}}
\End

\Data\ID{1103120005}
\Updated{2021-04-07 07:04:12}
\Level{B}
\SubArea{Quadratic Functions}
\IMG{1103120005a.pdf}
\LANGen{
\Question{
In the picture A,  we are given the graph of the quadratic function \( f(x)=x^2 \).  Use the graph of \( f \) as help to identify which of the graphs given in the picture B is the graph of \( g(x)=\frac32(x+3)^2-2 \). Choose what is the colour of the graph of \( g \). (Note: The graphs in the picture B were obtained  by shifting and stretching  the graph of \( f \).)}
{green}{red}{blue}{yellow}{}{}}
\LANGpl{
\Question{
Na rysunku A,  przedstawiono wykres funkcji kwadratowej  \( f(x)=x^2 \). Używając wykresu funkcji \( f \) określ, który z wykresów przedstawionych na rysunku B jest wykresem funkcji \( g(x)=\frac32(x+3)^2-2 \). Wybierz kolor przedstawiający kolor wykresu funkcji \( g \). (Wskazówka: Wykresy na rysunku  B powstały w wyniku  przesunięcia i  rozciągnięcia  wykresu funkcji \( f \).)}
{zielony}{czerwony}{niebieski}{żółty}{}{}}
\LANGcs{
\Question{
Na obrázku A je nakreslen graf kvadratické funkce \( f(x)=x^2 \).  Grafy na obrázku B vznikly posunutím a protažením grafu funkce \( f \). Určete barvu grafu funkce \( g(x)=\frac32(x+3)^2-2 \).}
{zelená}{červená}{modrá}{žlutá}{}{}}
\LANGsk{
\Question{
Na obrázku A je nakreslený graf kvadratickej funkcie \( f(x)=x^2 \). Grafy na obrázku B vznikli posunutím a natiahnutím grafu funkcie \( f \). Určte farbu grafu funkcie \( g(x)=\frac32(x+3)^2-2 \).}
{zelená}{červená}{modrá}{žltá}{}{}}
\LANGes{
\Question{
En el dibujo A tenemos la gráfica de la función cuadrática \( f(x)=x^2 \).  Usa la gráfica de \( f \) como ayuda para identificar cuál de las gráficas dadas en el dibujo B es la gráfica de \( g(x)=\frac32(x+3)^2-2 \). Elije el color de la gráfica de \( g \). (Nota: Las gráficas en el dibujo B fueron obtenidas moviendo y modificando la gráfica de \( f \).)}
{verde}{rojo}{azul}{amarillo}{}{}}
\End

\Data\ID{1103120006}
\Updated{2021-04-07 07:04:29}
\Level{B}
\SubArea{Quadratic Functions}
\IMG{1103120006a.pdf}
\LANGen{
\Question{
In the picture A,  we are given the graph of the quadratic function \( f(x)=x^2 \).  Use the graph of \( f \) as help to identify which of the graphs given in the picture B is the graph of \( g(x)=-(x+1)^2-3 \). Choose what is the colour of the graph of \( g \).}
{blue}{red}{yellow}{green}{}{}}
\LANGpl{
\Question{
Na rysunku A,  przestawiono wykres funkcji kwadratowej \( f(x)=x^2 \).  Użyj wykresu funkcji \( f \), aby określić, który z wykresów przedstawionych na rysunku  B jest wykresem \( g(x)=-(x+1)^2-3 \). Wybierz kolor wykresu \( g \).}
{niebieski}{czerwony}{żółty}{zielony}{}{}}
\LANGcs{
\Question{
Na obrázku A je nakreslen graf kvadratické funkce \( f(x)=x^2 \).  Jeden z grafů na obrázku B je graf funkce \( g(x)=-(x+1)^2-3 \). Určete, jakou barvu má graf funkce \( g \).}
{modrá}{červená}{žlutá}{zelená}{}{}}
\LANGsk{
\Question{
Na obrázku A je nakreslený graf kvadratickej funkcie \( f(x)=x^2 \).  Využitím grafu funkcie \( f \) určte farbu grafu funkcie \( g(x)=-(x+1)^2-3 \), ktorý je jedným z grafov na obrázku B.}
{modrá}{červená}{žltá}{zelená}{}{}}
\LANGes{
\Question{
En el dibujo A tenemos la gráfica de la función cuadrática \( f(x)=x^2 \). Usa la gráfica de \( f \) como ayuda para identificar cuál de las gráficas dadas en el dibujo B es la de \( g(x)=-(x+1)^2-3 \). Elije qué color tiene la gráfica de \( g \).}
{azul}{rojo}{amarillo}{verde}{}{}}
\End

\Data\ID{1103120007}
\Updated{2021-04-07 07:04:06}
\Level{B}
\SubArea{Quadratic Functions}
\IMG{1103120007.pdf}
\LANGen{
\Question{
Let \( f(x)=x^2 \). Given the graph of the function \( f \) and the graph of a function \( g \) which was obtained by shifting  the graph of \( f \) (see the picture), choose the function \( g \).}
{\( g(x) = (x+2)^2-4 \)}{\( g(x) = (x-2)^2-4 \)}{\( g(x)=(x-4)^2-2 \)}{\( g(x) = (x-2)^2+4 \)}{}{}}
\LANGpl{
\Question{
Niech  \( f(x)=x^2 \). Dany jest wykres funkcji \( f \) i wykres funkcji \( g \), który uzyskano poprzez przesunięcie wykresu funkcji  \( f \) (rysunek poniżej), wybierz funkcję \( g \).}
{\( g(x) = (x+2)^2-4 \)}{\( g(x) = (x-2)^2-4 \)}{\( g(x)=(x-4)^2-2 \)}{\( g(x) = (x-2)^2+4 \)}{}{}}
\LANGcs{
\Question{
Na obrázku je nakreslen graf funkce \( f(x)=x^2 \) a graf funkce \( g \), který vznikl posunem grafu funkce \( f \). Vyberte předpis funkce \( g \).}
{\( g(x) = (x+2)^2-4 \)}{\( g(x) = (x-2)^2-4 \)}{\( g(x)=(x-4)^2-2 \)}{\( g(x) = (x-2)^2+4 \)}{}{}}
\LANGsk{
\Question{
Na obrázku je nakreslený graf funkcie \( f(x)=x^2 \) a graf funkcie \( g \), ktorý vznikol posunutím grafu funkcie \( f \). Vyberte predpis funkcie \( g \).}
{\( g(x) = (x+2)^2-4 \)}{\( g(x) = (x-2)^2-4 \)}{\( g(x)=(x-4)^2-2 \)}{\( g(x) = (x-2)^2+4 \)}{}{}}
\LANGes{
\Question{
Sea \( f(x)=x^2 \). Dada la gráfica de la función \( f \) y la gráfica de la función \( g \) obtenida moviendo la gráfica de \( f \) (mira el dibujo), elije la función \( g \).}
{\( g(x) = (x+2)^2-4 \)}{\( g(x) = (x-2)^2-4 \)}{\( g(x)=(x-4)^2-2 \)}{\( g(x) = (x-2)^2+4 \)}{}{}}
\End

\Data\ID{1103120008}
\Updated{2020-10-09 08:10:41}
\Level{B}
\SubArea{Quadratic Functions}
\IMG{1103120008.pdf}
\LANGen{
\Question{
Let \( f(x)=x^2 \). Given the graph of the function \( f \) and the graph of a function \( g \) (see the picture), choose the function \( g \).}
{\( g(x) = -(x+2)^2+4 \)}{\( g(x)=(x+2)^2+4 \)}{\( g(x)=-(x-2)^2+4 \)}{\( g(x)=(x-2)^2-4 \)}{}{}}
\LANGpl{
\Question{
Niech \( f(x)=x^2 \). Dany jest  wykres funkcji  \( f \)  i wykres funkcji \( g \) (rysunek poniżej), wybierz funkcję \( g \).}
{\( g(x) = -(x+2)^2+4 \)}{\( g(x)=(x+2)^2+4 \)}{\( g(x)=-(x-2)^2+4 \)}{\( g(x)=(x-2)^2-4 \)}{}{}}
\LANGcs{
\Question{
Na obrázku je nakreslen graf funkce \( f(x)=x^2 \) a graf funkce \( g \). Vyberte předpis funkce \( g \).}
{\( g(x) = -(x+2)^2+4 \)}{\( g(x)=(x+2)^2+4 \)}{\( g(x)=-(x-2)^2+4 \)}{\( g(x)=(x-2)^2-4 \)}{}{}}
\LANGsk{
\Question{
Na obrázku je nakreslený graf funkcie \( f(x)=x^2 \)  a graf funkcie \( g \). Vyberte predpis funkcie \( g \).}
{\( g(x) = -(x+2)^2+4 \)}{\( g(x)=(x+2)^2+4 \)}{\( g(x)=-(x-2)^2+4 \)}{\( g(x)=(x-2)^2-4 \)}{}{}}
\LANGes{
\Question{
Sea \( f(x)=x^2 \). Dadas las gráficas de las funciones \( f \) y \( g \) (en el dibujo), elije la función \( g \).}
{\( g(x) = -(x+2)^2+4 \)}{\( g(x)=(x+2)^2+4 \)}{\( g(x)=-(x-2)^2+4 \)}{\( g(x)=(x-2)^2-4 \)}{}{}}
\End

\Data\ID{1103123808}
\Updated{2021-04-07 07:04:41}
\Level{B}
\SubArea{Quadratic Functions}
\IMG{1103123808_0.pdf}
\LANGen{
\Question{
Choose the equation of which the solution is graphed in red in the picture.}
{\( -x^2+2=0 \)}{\( -x^2-2=0 \)}{\( x^2+\sqrt2=0 \)}{\( (x-\sqrt2)(x+\sqrt2)=0 \)}{}{}}
\LANGpl{
\Question{
Wybierz równanie, którego rozwiązanie czerwono przedstawiono na rysunku.}
{\( -x^2+2=0 \)}{\( -x^2-2=0 \)}{\( x^2+\sqrt2=0 \)}{\( (x-\sqrt2)(x+\sqrt2)=0 \)}{}{}}
\LANGcs{
\Question{
Vyberte rovnici, jejíž řešení je červeně znázorněno na obrázku.}
{\( -x^2+2=0 \)}{\( -x^2-2=0 \)}{\( x^2+\sqrt2=0 \)}{\( (x-\sqrt2)(x+\sqrt2)=0 \)}{}{}}
\LANGsk{
\Question{
Vyberte rovnicu, ktorej riešenie je znázornené červenou na obrázku.}
{\( -x^2+2=0 \)}{\( -x^2-2=0 \)}{\( x^2+\sqrt2=0 \)}{\( (x-\sqrt2)(x+\sqrt2)=0 \)}{}{}}
\LANGes{
\Question{
Elige la ecuación cuya solución está en rojo el dibujo.}
{\( -x^2+2=0 \)}{\( -x^2-2=0 \)}{\( x^2+\sqrt2=0 \)}{\( (x-\sqrt2)(x+\sqrt2)=0 \)}{}{}}
\End

\Data\ID{1103123809}
\Updated{2021-04-07 07:04:03}
\Level{B}
\SubArea{Quadratic Functions}
\IMG{1103123809_0.pdf}
\LANGen{
\Question{
Choose the equation of which the solution is graphed in the picture.}
{\( x^2-4x=0 \)}{\( x^2-4=0 \)}{\( x^2-2x=0 \)}{\( x^2-2=0 \)}{}{}}
\LANGpl{
\Question{
Wybierz równanie, którego rozwiązanie przedstawiono na rysunku poniżej.}
{\( x^2-4x=0 \)}{\( x^2-4=0 \)}{\( x^2-2x=0 \)}{\( x^2-2=0 \)}{}{}}
\LANGcs{
\Question{
Vyberte rovnici, jejíž grafické řešení je znázorněno na obrázku.}
{\( x^2-4x=0 \)}{\( x^2-4=0 \)}{\( x^2-2x=0 \)}{\( x^2-2=0 \)}{}{}}
\LANGsk{
\Question{
Vyberte rovnicu, ktorej grafické riešenie je znázornené na obrázku.}
{\( x^2-4x=0 \)}{\( x^2-4=0 \)}{\( x^2-2x=0 \)}{\( x^2-2=0 \)}{}{}}
\LANGes{
\Question{
Elige la ecuación cuya solución está en el dibujo.}
{\( x^2-4x=0 \)}{\( x^2-4=0 \)}{\( x^2-2x=0 \)}{\( x^2-2=0 \)}{}{}}
\End

\Data\ID{1103148604}
\Updated{2021-04-07 08:04:02}
\Level{B}
\SubArea{Quadratic Functions}
\LANGen{
\Question{
The total mechanical energy \( E \) of an object is given by the formula \( E=mgh+\frac12mv^2 \), where \( m \) is the mass of the object, \( g \) is the acceleration due to gravity (about \( 10\,\mathrm{m}/\mathrm{s}^2 \)), \( h \) is the height of the object above the ground and \( v \) is the velocity of that object. Suppose that an object of the fixed mass \( m \) moves horizontally in the constant height \( h \) above the ground. Choose the graph, which could represent the dependence between the total mechanical energy (\( E \)) and velocity (\( v \)) of the object.}
{}{}{}{}{}{}}
\LANGpl{
\Question{
Całkowita energia mechaniczna \( E \) przedmiotu wyrażona jest za pomocą wzoru \( E=mgh+\frac12mv^2 \), gdzie \( m \) stanowi masę przedmiotu, \( g \) jest przyspieszeniem spowodowanym grawitacją (około\( 10\,\mathrm{m}/\mathrm{s}^2 \)), \( h \) jest wysokością przedmiotu nad ziemią, a \( v \) jest prędkością przedmiotu. Załóżmy, ze przedmiot o określonej masie \( m \) porusza się poziomo na stałej wysokości \( h \) nad ziemią.  Wybierz wykres przedstawiający zależność pomiędzy całkowitą energią mechaniczną  (\( E \)) a prędkością (\( v \)) tego przedmiotu.}
{}{}{}{}{}{}}
\LANGcs{
\Question{
Celková mechanická energie tělesa \( E \) je určena vztahem  \( E=mgh+\frac12mv^2 \), kde \( m \) je hmotnost tělesa, \( g \) je tíhové zrychlení (cca \( 10\,\mathrm{m}/\mathrm{s}^2 \)), \( h \) je výška tělesa nad povrchem Země a \( v \) je rychlost tělesa.  Vyberte obrázek, který může vyjadřovat závislost celkové mechanické energie (\( E \)) na rychlosti tělesa (\( v \)). Předpokládáme pohyb tělesa ve stálé výšce \( h \) a  neměnnou hmotnost tělesa \( m \).}
{}{}{}{}{}{}}
\LANGsk{
\Question{
Celková mechanická energie telesa \( E \) je určená vzťahom  \( E=mgh+\frac12mv^2 \), kde \( m \) je hmotnosť telesa, \( g \) je tiažové zrýchlenie (cca \( 10\,\mathrm{m}/\mathrm{s}^2 \)), \( h \) je výška telesa nad povrchom Zeme a \( v \) je rýchlosť telesa.  Vyberte obrázok, ktorý môže vyjadrovať závislosť celkovej mechanickej energie (\( E \)) na rýchlosti telesa (\( v \)). Predpokladáme pohyb telesa v stálej výške \( h \) s nemennou hmotnosť telesa \( m \).}
{}{}{}{}{}{}}
\LANGes{
\Question{
La energía mecánica total \( E \) de un objeto viene dada por la fórmula \( E=mgh+\frac12mv^2 \), dónde \( m \) es la masa del objeto, \( g \) es la aceleración de la gravedad (aproximadamente \( 10\,\mathrm{m}/\mathrm{s}^2 \)), \( h \) es la altura del objeto sobre la tierra y \( v \) es la velocidad del objeto. Supongamos que un objeto de masa fija \( m \) se mueve horizontalmente a una altura constante \( h \) sobre la tierra. Elige la gráfica que puede representar la dependencia entre la energía mecánica total (\( E \)) y la velocidad (\( v \)) del objeto.}
{}{}{}{}{}{}}

\IMGanswer{1103148604a_0.pdf}
\IMGanswer{1103148604b_0.pdf}
\IMGanswer{1103148604c_0.pdf}
\IMGanswer{1103148604d_0.pdf}\End

\Data\ID{2000004302}
\Updated{2022-02-20 08:02:58}
\Level{B}
\SubArea{Quadratic Functions}
\IMG{2000004401-priklad2.pdf}
\LANGen{
\Question{
In the picture A, the graph of the quadratic function \( f(x) = x^2\) is shown. With the aid of the graph of \(f\) identify which of the graphs shown in the picture B is the graph of \( g(x) = -\frac{1}{2} x^2\). What is the color of the graph of \(g\)? (Note: Each graph in the picture B is the graph of some transformation of \(f\).)}
{green}{blue}{yellow}{red}{}{}}
\LANGpl{
\Question{
Na rysunku A przedstawiony jest wykres funkcji kwadratowej \( f(x) = x^2\). Z pomocą wykresu funkcji \(f\) wskaż, który z wykresów na rysunku B jest wykresem funkcji  \( g(x) = -\frac{1}{2} x^2\). Jaki kolor ma funkcja \(g\)? (Uwaga: każdy wykres na rysunku B jest pewnym przekształceniem funkcji \(f\).)}
{zielony}{niebieski}{żółty}{czerwony}{}{}}
\LANGcs{
\Question{
Na obrázku A je nakreslen graf kvadratické funkce  \( f(x) = x^2\).  Grafy na obrázku B vznikly transformací grafu funkce \(f\).  Určete barvu funkce  \( g(x) = -\frac{1}{2} x^2\).}
{zelená}{modrá}{žlutá}{červená}{}{}}
\LANGsk{
\Question{
Na obrázku A je nakreslený graf kvadratickej funkcie  \( f(x) = x^2\).  Grafy na obrázku B vznikli transformáciou grafu funkcie \(f\).  Určte farbu funkcie  \( g(x) = -\frac{1}{2} x^2\).}
{zelená}{modrá}{žltá}{červená}{}{}}
\LANGes{
\Question{
En el dibujo A aparece la gráfica de la función cuadrática \( f(x) = x^2\).  Las gráficas en el dibujo B son transformaciones de \(f\).  Elige el color de la gráfica de \( g(x) = -\frac{1}{2} x^2\).}
{verde}{azul}{amarillo}{rojo}{}{}}
\End

\Data\ID{2000004303}
\Updated{2022-02-20 08:02:21}
\Level{B}
\SubArea{Quadratic Functions}
\IMG{2000004401-figure0_3_0.pdf}
\LANGen{
\Question{
In the picture A, the graph of the quadratic function \( f(x) = x^2\) is shown. With the aid of the graph of \(f\) identify which of the graphs shown in the picture B is the graph of \( g(x) = (x+2)^2\). What is the color of the graph of \(g\)? (Note: Each graph in the picture B was obtained by shifting the graph of \(f\).)}
{red}{green}{yellow}{blue}{}{}}
\LANGpl{
\Question{
Na rysunku A przedstawiony jest wykres funkcji \( f(x) = x^2\). Z pomocą wykresu funkcji \(f\) wskaż, który z wykresów na rynku B  jest wykresem funkcji \( g(x) = (x+2)^2\). Wskaż odpowiedni kolor wykresu funkcji \(g\)? (Uwaga:  Każdy wykres na rysunku B został uzyskany poprzez przesunięcie wykresu funkcji \(f\).)}
{czerwony}{zielony}{żółty}{niebieski}{}{}}
\LANGcs{
\Question{
Na obrázku A je nakreslen graf kvadratické funkce  \( f(x) = x^2\).  Grafy na obrázku B vznikly transformací grafu funkce \(f\).  Určete barvu funkce \( g(x) = (x+2)^2\).}
{červená}{zelená}{žlutá}{modrá}{}{}}
\LANGsk{
\Question{
Na obrázku A je nakreslený graf kvadratickej funkcie  \( f(x) = x^2\).  Grafy na obrázku B vznikli transformáciou grafu funkcie \(f\).  Určte farbu funkcie \( g(x) = (x+2)^2\).}
{červená}{zelená}{žltá}{modrá}{}{}}
\LANGes{
\Question{
En el dibujo A aparece la gráfica de la función cuadrática \( f(x) = x^2\).  Las gráficas en el dibujo B son transformaciones de \(f\).  Elige el color de la gráfica de \( g(x) = (x+2)^2\).}
{rojo}{verde}{amarillo}{azul}{}{}}
\End

\Data\ID{2000004304}
\Updated{2022-02-20 08:02:42}
\Level{B}
\SubArea{Quadratic Functions}
\IMG{2000004401_priklad4.pdf}
\LANGen{
\Question{
In the picture A, the graph of the quadratic function \( f(x) = x^2\) is shown. With the aid of the graph of \(f\) identify which of the graphs shown in the picture B is the graph of \( g(x) = (x-2)^2\). What is the color of the graph of \(g\)? (Note: Each graph in the picture B is the graph of some transformation of \(f\).)}
{yellow}{red}{green}{blue}{}{}}
\LANGpl{
\Question{
Rysunek A jest wykresem funkcji kwadratowej \( f(x) = x^2\). Wykresy na rysunku B zostały utworzone poprzez przekształcenie wykresu funkcji \(f\). Określ kolor funkcji \( g(x) = (x-2)^2\).}
{żółty}{czerwony}{zielony}{niebieski}{}{}}
\LANGcs{
\Question{
Na obrázku A je nakreslen graf kvadratické funkce  \( f(x) = x^2\).  Grafy na obrázku B vznikly transformací grafu funkce \(f\).  Určete barvu funkce \( g(x) = (x-2)^2\).}
{žlutá}{červená}{zelená}{modrá}{}{}}
\LANGsk{
\Question{
Na obrázku A je nakreslený graf kvadratickej funkcie  \( f(x) = x^2\).  Grafy na obrázku B vznikli transformáciou grafu funkcie \(f\).  Určte farbu funkcie \( g(x) = (x-2)^2\).}
{žltý}{červený}{zelený}{modrý}{}{}}
\LANGes{
\Question{
En el dibujo A aparece la gráfica de la función \( f(x) = x^2\).  Las gráficas en el dibujo B son transformaciones de la función \(f\).  Elige el color de la función \( g(x) = (x-2)^2\).}
{amarillo}{rojo}{verde}{azul}{}{}}
\End

\Data\ID{2000004305}
\Updated{2022-02-20 08:02:19}
\Level{B}
\SubArea{Quadratic Functions}
\IMG{2000004401-figure0_4.pdf}
\LANGen{
\Question{
In the picture A, the graph of the quadratic function \( f(x) = x^2\) is shown. With the aid of the graph of \(f\) identify which of the graphs shown in the picture B is the graph of \( g(x) = x^2-2x-3=(x+1)(x-3)\). What is the color of the graph of \(g\)? (Note: Each graph in the picture B was obtained by shifting the graph of \(f\).)}
{blue}{red}{yellow}{green}{}{}}
\LANGpl{
\Question{
Rysunek A jest wykresem funkcji kwadratowej \( f(x) = x^2\). Wykresy na rysunku B zostały utworzone poprzez przekształcenie wykresu funkcji \(f\).  Określ kolor funkcji \(g(x) = x^2-2x-3=(x+1)(x-3)\).}
{niebieski}{czerwony}{żółty}{zielony}{}{}}
\LANGcs{
\Question{
Na obrázku A je nakreslen graf kvadratické funkce  \( f(x) = x^2\).  Grafy na obrázku B vznikly transformací grafu funkce \(f\).  Určete barvu funkce \( g(x) = x^2-2x-3=(x+1)(x-3)\).}
{modrá}{červená}{žlutá}{zelená}{}{}}
\LANGsk{
\Question{
Na obrázku A je nakreslený graf kvadratické funkcie  \( f(x) = x^2\).  Grafy na obrázku B vznikli transformáciou grafu funkcie \(f\).  Určte farbu funkcie \( g(x) = x^2-2x-3=(x+1)(x-3)\).}
{modrý}{červený}{žltý}{zelený}{}{}}
\LANGes{
\Question{
En el dibujo A aparece la gráfica de la función cuadrática \( f(x) = x^2\).  Las gráficas en el dibujo B son transformaciones de la gráfica de \(f\).  Elige el color de la función \( g(x) = x^2-2x-3=(x+1)(x-3)\).}
{azul}{rojo}{amarillo}{verde}{}{}}
\End

\Data\ID{2000004306}
\Updated{2022-02-20 08:02:38}
\Level{B}
\SubArea{Quadratic Functions}
\IMG{2000004401-figure6.pdf}
\LANGen{
\Question{
Let \(f(x) = x^2\). Given the graph of the function \(f\) and the graph of a function \(g\) (see the picture), choose the function \(g\).}
{\( g(x) =-x^2+4\)}{\( g(x) =x^2-4\)}{\( g(x) =x^2-2\)}{\( g(x) =-x^2+2\)}{}{}}
\LANGpl{
\Question{
Dana jest funkcja \(f(x) = x^2\). Na podstawie wykresu funkcji \(f\) oraz funkcji \(g\) (patrz rysunek), wskaż funkcję \(g\).}
{\( g(x) =-x^2+4\)}{\( g(x) =x^2-4\)}{\( g(x) =x^2-2\)}{\( g(x) =-x^2+2\)}{}{}}
\LANGcs{
\Question{
Na obrázku je nakreslen graf funkce \(f(x) = x^2\) a graf funkce \(g(x)\). Vyberte předpis funkce  \(g(x)\).}
{\( g(x) =-x^2+4\)}{\( g(x) =x^2-4\)}{\( g(x) =x^2-2\)}{\( g(x) =-x^2+2\)}{}{}}
\LANGsk{
\Question{
Na obrázku je nakreslený graf funkcie \(f(x) = x^2\) a graf funkcie \(g(x)\). Vyberte prepis funkcie  \(g(x)\).}
{\( g(x) =-x^2+4\)}{\( g(x) =x^2-4\)}{\( g(x) =x^2-2\)}{\( g(x) =-x^2+2\)}{}{}}
\LANGes{
\Question{
En el dibujo aparecen las gráficas de las funciónes \(f(x) = x^2\) y \(g(x)\). Elige la ecuación de la función \(g(x)\).}
{\( g(x) =-x^2+4\)}{\( g(x) =x^2-4\)}{\( g(x) =x^2-2\)}{\( g(x) =-x^2+2\)}{}{}}
\End

\Data\ID{2010017001}
\Updated{2022-11-02 09:11:42}
\Level{B}
\SubArea{Quadratic Functions}
\LANGen{
\Question{
The graph of the function \(f(x) = -4x^{2} + 2\) is a parabola. Which of the following points is the vertex of this parabola?}
{\([0;2]\)}{\(\left[\frac{\sqrt2}2;0\right]\)}{\([2;0]\)}{\([1;-2]\)}{}{}}
\LANGpl{
\Question{
Wykres funkcji \(f(x) = -4x^{2} + 2\) jest parabolą. Który z poniższych punktów jest wierzchołkiem tej paraboli?}
{\([0;2]\)}{\(\left[\frac{\sqrt2}2;0\right]\)}{\([2;0]\)}{\([1;-2]\)}{}{}}
\LANGcs{
\Question{
Grafem funkce  \(f(x) = -4x^{2} + 2\) je parabola. Který z následujících bodů je vrcholem této paraboly?}
{\([0;2]\)}{\(\left[\frac{\sqrt2}2;0\right]\)}{\([2;0]\)}{\([1;-2]\)}{}{}}
\LANGsk{
\Question{
Grafom funkcie  \(f(x) = -4x^{2} + 2\) je parabola. Ktorý z nasledujúcich bodov je vrcholom tejto paraboly?}
{\([0;2]\)}{\(\left[\frac{\sqrt2}2;0\right]\)}{\([2;0]\)}{\([1;-2]\)}{}{}}
\LANGes{
\Question{
La gráfica de la función \(f(x) = -4x^{2} + 2\) es una parábola. ¿Cuál de los siguientes puntos es el vértice de esta parábola?}
{\([0;2]\)}{\(\left[\frac{\sqrt2}2;0\right]\)}{\([2;0]\)}{\([1;-2]\)}{}{}}
\End

\Data\ID{2010017002}
\Updated{2022-11-02 09:11:42}
\Level{B}
\SubArea{Quadratic Functions}
\LANGen{
\Question{
The graph of the function \(f(x) = x^{2} + 4x + 7\) is a parabola. Which of the following points is the vertex of this parabola?}
{\([-2;3]\)}{\([3;-2]\)}{\([0;7]\)}{\([1;12]\)}{}{}}
\LANGpl{
\Question{
Wykres funkcji  \(f(x) = x^{2} + 4x + 7\) jest parabolą. Który z poniższych punktów jest wierzchołkiem tej paraboli?}
{\([-2;3]\)}{\([3;-2]\)}{\([0;7]\)}{\([1;12]\)}{}{}}
\LANGcs{
\Question{
Grafem funkce  \(f(x) = x^{2} + 4x + 7\) je parabola. Který z následujících bodů je vrcholem této paraboly?}
{\([-2;3]\)}{\([3;-2]\)}{\([0;7]\)}{\([1;12]\)}{}{}}
\LANGsk{
\Question{
Grafom funkcie \(f(x) = x^{2} + 4x + 7\) je parabola. Ktorý z nasledujúcich bodov je vrcholom paraboly?}
{\([-2;3]\)}{\([3;-2]\)}{\([0;7]\)}{\([1;12]\)}{}{}}
\LANGes{
\Question{
La gráfica de la función \(f(x) = x^{2} + 4x + 7\) es una parábola. ¿Cuál de los siguientes puntos es el vértice de esta parábola?}
{\([-2;3]\)}{\([3;-2]\)}{\([0;7]\)}{\([1;12]\)}{}{}}
\End

\Data\ID{2010017003}
\Updated{2022-11-02 09:11:42}
\Level{B}
\SubArea{Quadratic Functions}
\LANGen{
\Question{
Find the maximum value of the quadratic function
\(f(x)= -x^{2} +2x +1\).}
{\(2\)}{\(1\)}{The maximum value of the function \(f\) does not exist.}{\(0\)}{}{}}
\LANGpl{
\Question{
Wyznacz maksimum funkcji kwadratowej
\(f(x)= -x^{2} +2x +1\).}
{\(2\)}{\(1\)}{Funkcja nie ma maksimum,}{\(0\)}{}{}}
\LANGcs{
\Question{
Určete největší hodnotu, které nabývá kvadratická funkce 
\(f(x)= -x^{2} +2x +1\).}
{\(2\)}{\(1\)}{Největší hodnota funkce \(f\) neexistuje.}{\(0\)}{}{}}
\LANGsk{
\Question{
Určte najväčšiu hodnotu, ktorú nadobúda kvadratická funkcia 
\(f(x)= -x^{2} +2x +1\).}
{\(2\)}{\(1\)}{Najväčšia hodnota funkcie \(f\) neexistuje.}{\(0\)}{}{}}
\LANGes{
\Question{
Halla el valor máximo de la función cuadrática
\(f(x)= -x^{2} +2x +1\).}
{\(2\)}{\(1\)}{No existe.}{\(0\)}{}{}}
\End

\Data\ID{2010017004}
\Updated{2022-11-02 09:11:42}
\Level{B}
\SubArea{Quadratic Functions}
\LANGen{
\Question{
The graph of the function \( f(x)=3x^2-6x-3\) is a parabola. Which of the following points is the vertex of this parabola?}
{\( [1;-6] \)}{\( [0;-3] \)}{\( [1;-8] \)}{\( [-1;0] \)}{}{}}
\LANGpl{
\Question{
Wykres funkcji \( f(x)=3x^2-6x-3\) jest parabolą. Który z poniższych punktów jest wierzchołkiem tej paraboli?}
{\( [1;-6] \)}{\( [0;-3] \)}{\( [1;-8] \)}{\( [-1;0] \)}{}{}}
\LANGcs{
\Question{
Určete souřadnice vrcholu paraboly, která je grafem funkce  \( f(x)=3x^2-6x-3\).}
{\( [1;-6] \)}{\( [0;-3] \)}{\( [1;-8] \)}{\( [-1;0] \)}{}{}}
\LANGsk{
\Question{
Určte súradnice vrcholu paraboly, ktorá je grafom funkcie  \( f(x)=3x^2-6x-3\).}
{\( [1;-6] \)}{\( [0;-3] \)}{\( [1;-8] \)}{\( [-1;0] \)}{}{}}
\LANGes{
\Question{
La gráfica de la función \( f(x)=3x^2-6x-3\) es una parábola. ¿Cuál de los siguientes puntos es el vértice de esta parábola?}
{\( [1;-6] \)}{\( [0;-3] \)}{\( [1;-8] \)}{\( [-1;0] \)}{}{}}
\End

\Data\ID{2010017005}
\Updated{2022-11-02 09:11:42}
\Level{B}
\SubArea{Quadratic Functions}
\IMG{2010017001-figure0.pdf}
\LANGen{
\Question{
Let \( f(x)=3x^2 \). Given the graph of the function \( f \) and the graph of a function \( g \) which was obtained as a left shift  of the graph of \( f \) (see the picture), choose the function \( g \).}
{\( g(x) = 3(x+2)^2 \)}{\( g(x) = 3(x-2)^2 \)}{\( g(x) = 3x^2+3 \)}{\( g(x) = 3x^2 +12 \)}{}{}}
\LANGpl{
\Question{
Dana jest funkcja \( f(x)=3x^2 \). Mając wykres funkcji \( f \) i wykres funkcji \( g \) otrzymany jako przesunięcie w lewo wykresu \( f \) (patrz rysunek), wskaż funkcję \( g \).}
{\( g(x) = 3(x+2)^2 \)}{\( g(x) = 3(x-2)^2 \)}{\( g(x) = 3x^2+3 \)}{\( g(x) = 3x^2 +12 \)}{}{}}
\LANGcs{
\Question{
Na obrázku je nakreslen graf funkce \( f(x)=3x^2\) a graf funkce \(g\), který vznikl posunem grafu funkce \(f\). Vyberte předpis funkce \(g\).}
{\( g(x) = 3(x+2)^2 \)}{\( g(x) = 3(x-2)^2 \)}{\( g(x) = 3x^2+3 \)}{\( g(x) = 3x^2 +12 \)}{}{}}
\LANGsk{
\Question{
Na obrázku je nakreslený graf funkcie \( f(x)=3x^2\) a graf funkcie \(g\), ktorý vznikol posunutím grafu funkcie \(f\). Vyberte predpis funkcie \(g\).}
{\( g(x) = 3(x+2)^2 \)}{\( g(x) = 3(x-2)^2 \)}{\( g(x) = 3x^2+3 \)}{\( g(x) = 3x^2 +12 \)}{}{}}
\LANGes{
\Question{
Sea \( f(x)=3x^2 \). Dada la gráfica de la función \( f \) y la gráfica de la función \( g \) que obtenemos por un desplazamiento a la izquierda de la gráfica de \( f \) (ver la imagen), elige la función \( g \).}
{\( g(x) = 3(x+2)^2 \)}{\( g(x) = 3(x-2)^2 \)}{\( g(x) = 3x^2+3 \)}{\( g(x) = 3x^2 +12 \)}{}{}}
\End

\Data\ID{2010017006}
\Updated{2022-11-02 09:11:42}
\Level{B}
\SubArea{Quadratic Functions}
\LANGen{
\Question{
The graph of the function \( f \) is a parabola, vertex of which is \( [-4;0] \) and \( f(2)= 12 \) is given. Find the function \( f \).}
{\( f(x)=\frac13(x+4)^2 \)}{\( f(x)=\frac13(x-4)^2 \)}{\( f(x)=-\frac13(x+4)^2 \)}{\( f(x)=-\frac13(x-4)^2 \)}{}{}}
\LANGpl{
\Question{
Funkcja \( f \) jest parabolą, której wierzchołek wynosi \( [-4;0] \) oraz  \( f(2)= 12 \). Wskaż funkcję  \( f \).}
{\( f(x)=\frac13(x+4)^2 \)}{\( f(x)=\frac13(x-4)^2 \)}{\( f(x)=-\frac13(x+4)^2 \)}{\( f(x)=-\frac13(x-4)^2 \)}{}{}}
\LANGcs{
\Question{
Grafem funkce \(f\) je parabola s vrcholem v bodě \( [-4;0] \) a dále platí \( f(2)= 12 \). Určete funkci  \( f \).}
{\( f(x)=\frac13(x+4)^2 \)}{\( f(x)=\frac13(x-4)^2 \)}{\( f(x)=-\frac13(x+4)^2 \)}{\( f(x)=-\frac13(x-4)^2 \)}{}{}}
\LANGsk{
\Question{
Grafom funkcie \(f\) je parabola s vrcholom v bode \( [-4;0] \) a ďalej platí \( f(2)= 12 \). Určte funkciu  \( f \).}
{\( f(x)=\frac13(x+4)^2 \)}{\( f(x)=\frac13(x-4)^2 \)}{\( f(x)=-\frac13(x+4)^2 \)}{\( f(x)=-\frac13(x-4)^2 \)}{}{}}
\LANGes{
\Question{
La gráfica de la función \( f \) es una parábola, cuyo vértice es \( [-4;0] \) y \( f(2)= 12 \). Halla la función \( f \).}
{\( f(x)=\frac13(x+4)^2 \)}{\( f(x)=\frac13(x-4)^2 \)}{\( f(x)=-\frac13(x+4)^2 \)}{\( f(x)=-\frac13(x-4)^2 \)}{}{}}
\End

\Data\ID{9000007106}
\Updated{2021-04-07 07:04:02}
\Level{B}
\SubArea{Quadratic Functions}
\IMG{000071_06-figure0.pdf}
\LANGen{
\Question{
The picture shows a graph of a quadratic function and three points on the graph. In
the following list identify another point on the graph of the same quadratic
function.}
{\([0.4;0.16]\)}{\([0.4;0.2]\)}{\([0.4;0.8]\)}{\([0.04;0.16]\)}{}{}}
\LANGpl{
\Question{
Na rysunku przedstawiono wykres funkcji kwadratowej i trzy punkty na tym wykresie. Z poniższej listy wybierz kolejny punkt, który należy do wykresu tej samej funkcji kwadratowej.}
{\([0.4;0.16]\)}{\([0.4;0.2]\)}{\([0.4;0.8]\)}{\([0.04;0.16]\)}{}{}}
\LANGcs{
\Question{
Který z bodů leží na grafu znázorněné kvadratické funkce?}
{\([0{,}4;0{,}16]\)}{\([0{,}4;0{,}2]\)}{\([0{,}4;0{,}8]\)}{\([0{,}04;0{,}16]\)}{}{}}
\LANGsk{
\Question{
Na obrázku je graf kvadratickej funkcie a tri body, ktoré patria funkcii.
Určte, ktorý ďalší bod leží na grafe danej kvadratickej funkcie.}
{\([0{,}4;0{,}16]\)}{\([0{,}4;0{,}2]\)}{\([0{,}4;0{,}8]\)}{\([0{,}04;0{,}16]\)}{}{}}
\LANGes{
\Question{
En el dibujo aparece la gráfica de una función cuadrática y tres puntos pertenecientes a dicha gráfica. ¿Cuál de los puntos de la lista también pertenece a la gráfica?}
{\([0.4;0.16]\)}{\([0.4;0.2]\)}{\([0.4;0.8]\)}{\([0.04;0.16]\)}{}{}}
\End

\Data\ID{9000007107}
\Updated{2021-04-07 07:04:44}
\Level{B}
\SubArea{Quadratic Functions}
\IMG{000071_07-figure0.pdf}
\LANGen{
\Question{
The picture shows a graph of a quadratic function and three points on the graph. In
the following list identify another point on the graph of the same quadratic
function.}
{\([0.6;0.36]\)}{\([0.18;0.36]\)}{\([0.36;0.6]\)}{\([0.6;3.6]\)}{}{}}
\LANGpl{
\Question{
Na rysunku przedstawiono wykres funkcji kwadratowej i trzy punkty na tym wykresie. Z poniższej listy wybierz kolejny punkt, który należy do wykresu tej samej funkcji kwadratowej.}
{\([0.6;0.36]\)}{\([0.18;0.36]\)}{\([0.36;0.6]\)}{\([0.6;3.6]\)}{}{}}
\LANGcs{
\Question{
Který z bodů leží na grafu znázorněné kvadratické funkce?}
{\([0{,}6;0{,}36]\)}{\([0{,}18;0{,}36]\)}{\([0{,}36;0{,}6]\)}{\([0{,}6;3{,}6]\)}{}{}}
\LANGsk{
\Question{
Na obrázku je graf kvadratickej funkcie a tri body, ktoré patria funkcii. Určte, ktorý ďalší bod leží na grafe danej kvadratickej funkcie.}
{\([0{,}6;0{,}36]\)}{\([0{,}18;0{,}36]\)}{\([0{,}36;0{,}6]\)}{\([0{,}6;3{,}6]\)}{}{}}
\LANGes{
\Question{
En el dibujo aparece la gráfica de una función cuadrática y tres puntos pertenecientes a dicha gráfica. ¿Cuál de los puntos de la lista también pertenece a la gráfica?}
{\([0.6;0.36]\)}{\([0.18;0.36]\)}{\([0.36;0.6]\)}{\([0.6;3.6]\)}{}{}}
\End

\Data\ID{9000007108}
\Updated{2021-04-07 07:04:05}
\Level{B}
\SubArea{Quadratic Functions}
\IMG{000071_08-figure0.pdf}
\LANGen{
\Question{
The picture shows a graph of a quadratic function and three points on the graph. In
the following list identify another point on the graph of the same quadratic
function.}
{\([3;3]\)}{\([3;2]\)}{\([3;4]\)}{\([3;9]\)}{}{}}
\LANGpl{
\Question{
Na rysunku przedstawiono wykres funkcji kwadratowej i trzy punkty na tym wykresie. Z poniższej listy wybierz kolejny punkt, który należy do wykresu tej samej funkcji kwadratowej.}
{\([3;3]\)}{\([3;2]\)}{\([3;4]\)}{\([3;9]\)}{}{}}
\LANGcs{
\Question{
Který z bodů leží na grafu znázorněné kvadratické funkce?}
{\([3;3]\)}{\([3;2]\)}{\([3;4]\)}{\([3;9]\)}{}{}}
\LANGsk{
\Question{
Na obrázku je graf kvadratickej funkcie a tri body, ktoré patria funkcii. Určte, ktorý ďalší bod leží na grafe danej kvadratickej funkcie.}
{\([3;3]\)}{\([3;2]\)}{\([3;4]\)}{\([3;9]\)}{}{}}
\LANGes{
\Question{
En el dibujo aparece la gráfica de una función cuadrática y tres puntos pertenecientes a dicha gráfica. ¿Cuál de los puntos de la lista también pertenece a la gráfica?}
{\([3;3]\)}{\([3;2]\)}{\([3;4]\)}{\([3;9]\)}{}{}}
\End

\Data\ID{9000007109}
\Updated{2021-04-07 07:04:29}
\Level{B}
\SubArea{Quadratic Functions}
\IMG{000071_09-figure0.pdf}
\LANGen{
\Question{
The picture shows a graph of a quadratic function and three points on the graph. In
the following list identify another point on the graph of the same quadratic
function.}
{\([2;3]\)}{\([2;2]\)}{\([2;4]\)}{\([2;5]\)}{}{}}
\LANGpl{
\Question{
Na rysunku przedstawiono wykres funkcji kwadratowej i trzy punkty na tym wykresie. Z poniższej listy wybierz kolejny punkt, który należy do wykresu tej samej funkcji kwadratowej.}
{\([2;3]\)}{\([2;2]\)}{\([2;4]\)}{\([2;5]\)}{}{}}
\LANGcs{
\Question{
Který z bodů leží na grafu znázorněné kvadratické funkce?}
{\([2;3]\)}{\([2;2]\)}{\([2;4]\)}{\([2;5]\)}{}{}}
\LANGsk{
\Question{
Na obrázku je graf kvadratickej funkcie a tri body, ktoré patria funkcii. Určte, ktorý ďalší bod leží na grafe danej kvadratickej funkcie.}
{\([2;3]\)}{\([2;2]\)}{\([2;4]\)}{\([2;5]\)}{}{}}
\LANGes{
\Question{
En el dibujo aparece la gráfica de una función cuadrática y tres puntos pertenecientes a dicha gráfica. ¿Cuál de los puntos de la lista también pertenece a la gráfica?}
{\([2;3]\)}{\([2;2]\)}{\([2;4]\)}{\([2;5]\)}{}{}}
\End

\Data\ID{9000007110}
\Updated{2021-04-07 07:04:37}
\Level{B}
\SubArea{Quadratic Functions}
\IMG{000071_10-figure0.pdf}
\LANGen{
\Question{
The picture shows a graph of a quadratic function and three points on the graph. In
the following list identify another point on the graph of the same quadratic
function.}
{\([-1;-3]\)}{\([-1;-1]\)}{\([-1;-2]\)}{\([-1;-4]\)}{}{}}
\LANGpl{
\Question{
Na rysunku przedstawiono wykres funkcji kwadratowej i trzy punkty na tym wykresie. Z poniższej listy wybierz kolejny punkt, który należy do wykresu tej samej funkcji kwadratowej.}
{\([-1;-3]\)}{\([-1;-1]\)}{\([-1;-2]\)}{\([-1;-4]\)}{}{}}
\LANGcs{
\Question{
Který z bodů leží na grafu znázorněné kvadratické funkce?}
{\([-1;-3]\)}{\([-1;-1]\)}{\([-1;-2]\)}{\([-1;-4]\)}{}{}}
\LANGsk{
\Question{
Na obrázku je graf kvadratickej funkcie a tri body, ktoré patria funkcii. Určte, ktorý ďalší bod leží na grafe danej kvadratickej funkcie.}
{\([-1;-3]\)}{\([-1;-1]\)}{\([-1;-2]\)}{\([-1;-4]\)}{}{}}
\LANGes{
\Question{
En el dibujo aparece la gráfica de una función cuadrática y tres puntos sobre ella, ¿Cuál de los puntos de la lista también pertenece a la gráfica de la misma función?}
{\([-1;-3]\)}{\([-1;-1]\)}{\([-1;-2]\)}{\([-1;-4]\)}{}{}}
\End

\Data\ID{9000014803}
\Updated{2022-04-23 05:04:17}
\Level{B}
\SubArea{Quadratic Functions}
\LANGen{
\Question{
The graph of the function \(f(x) = 6x^{2} + 3\) is a parabola. Which of the following points is the vertex of this parabola?}
{\([0;3]\)}{\([3;0]\)}{\([1;9]\)}{\([1;2]\)}{}{}}
\LANGpl{
\Question{
Rozważmy parabolę zdefiniowaną jako wykres funkcji 
\(f\colon y = 6x^{2} + 3\).
Który z poniższych punktów jest wierzchołkiem tej paraboli?}
{\([0;3]\)}{\([3;0]\)}{\([1;9]\)}{\([1;2]\)}{}{}}
\LANGcs{
\Question{
Grafem funkce \(f\colon y = 6x^{2} + 3\) 
je parabola. Který z následujících bodů je vrcholem této paraboly?}
{\([0;3]\)}{\([3;0]\)}{\([1;9]\)}{\([1;2]\)}{}{}}
\LANGsk{
\Question{
Grafom funkcie \(f\colon y = 6x^{2} + 3\) je parabola. Ktorý z následujúcich bodov je vrcholom tejto paraboly?}
{\([0;3]\)}{\([3;0]\)}{\([1;9]\)}{\([1;2]\)}{}{}}
\LANGes{
\Question{
Considera una parábola definida como la gráfica de la función
\(f(x) = 6x^{2} + 3\).
¿Cuál de los puntos es el vértice de la parábola?}
{\([0;3]\)}{\([3;0]\)}{\([1;9]\)}{\([1;2]\)}{}{}}
\End

\Data\ID{9000014804}
\Updated{2022-04-23 05:04:17}
\Level{B}
\SubArea{Quadratic Functions}
\LANGen{
\Question{
The graph of the function \(f(x) = x^{2} - 4x + 13\) is a parabola. Which of the following points is the vertex of this parabola?}
{\([2;9]\)}{\([-2;13]\)}{\([-4;13]\)}{\([0;13]\)}{}{}}
\LANGpl{
\Question{
Rozważ parabolę zdefiniowaną jako wykres funkcji
\(f\colon y = x^{2} - 4x + 13\).
Który z poniższych punktów jest wierzchołkiem tej paraboli?}
{\([2;9]\)}{\([-2;13]\)}{\([-4;13]\)}{\([0;13]\)}{}{}}
\LANGcs{
\Question{
Grafem funkce \(f\colon y = x^{2} - 4x + 13\) 
je parabola. Který z následujících bodů je vrcholem této paraboly?}
{\([2;9]\)}{\([-2;13]\)}{\([-4;13]\)}{\([0;13]\)}{}{}}
\LANGsk{
\Question{
Grafom funkcie \(f\colon y = x^{2} - 4x + 13\) 
je parabola. Ktorý z následujúcich bodov je vrcholom tejto paraboly?}
{\([2;9]\)}{\([-2;13]\)}{\([-4;13]\)}{\([0;13]\)}{}{}}
\LANGes{
\Question{
Considera una parábola definida como la gráfica de la función
\(f(x) = x^{2} - 4x + 13\).
¿Cuál de los puntos es el vértice de la parábola?}
{\([2;9]\)}{\([-2;13]\)}{\([-4;13]\)}{\([0;13]\)}{}{}}
\End

\Data\ID{9000014805}
\Updated{2022-04-23 05:04:17}
\Level{B}
\SubArea{Quadratic Functions}
\LANGen{
\Question{
Find the minimum value of the quadratic function
\(f(x)= 4x^{2} - 4x + 7\).}
{\(6\)}{\(7\)}{does not exist}{\(- 4\)}{}{}}
\LANGpl{
\Question{
Znajdź minimum wartości funkcji kwadratowej
\(f\colon y = 4x^{2} - 4x + 7\).}
{\(6\)}{\(7\)}{nie istnieje}{\(- 4\)}{}{}}
\LANGcs{
\Question{
Určete nejmenší hodnotu, které nabývá kvadratická funkce
\(f\colon y = 4x^{2} - 4x + 7\).}
{\(6\)}{\(7\)}{neexistuje}{\(- 4\)}{}{}}
\LANGsk{
\Question{
Určte najmenšiu hodnotu kvadratickej funkcie
\(f\colon y = 4x^{2} - 4x + 7\).}
{\(6\)}{\(7\)}{neexistuje}{\(- 4\)}{}{}}
\LANGes{
\Question{
Encuentra el valor mínimo de la función cuadrática
\(f(x)= 4x^{2} - 4x + 7\).}
{\(6\)}{\(7\)}{no existe}{\(- 4\)}{}{}}
\End

\Data\ID{9000014806}
\Updated{2024-09-23 09:09:16}
\Level{B}
\SubArea{Quadratic Functions}
\LANGen{
\Question{
Find the maximum value of the quadratic function
\(f(x) = 0.02x^{2} - 7x + 4\).}
{does not exist}{\(4\)}{\(0.02\)}{\(- 7\)}{}{}}
\LANGpl{
\Question{
Znajdź maksimum wartości funkcji kwadratowej
\(f\colon y = 0.02x^{2} - 7x + 4\).}
{nie istnieje}{\(4\)}{\(0.02\)}{\(- 7\)}{}{}}
\LANGcs{
\Question{
Určete největší hodnotu, které nabývá kvadratická funkce
\(f\colon y = 0{,}02x^{2} - 7x + 4\).}
{neexistuje}{\(4\)}{\(0{,}02\)}{\(- 7\)}{}{}}
\LANGsk{
\Question{
Určte najväčšiu hodnotu kvadratickej funkcie
\(f\colon y = 0{,}02x^{2} - 7x + 4\).}
{neexistuje}{\(4\)}{\(0{,}02\)}{\(- 7\)}{}{}}
\LANGes{
\Question{
Encuentra el valor máximo de la función cuadrática
\(f(x) = 0.02x^{2} - 7x + 4\).}
{no existe}{\(4\)}{\(0.02\)}{\(- 7\)}{}{}}
\End

\Data\ID{9000022306}
\Updated{2020-11-20 12:11:48}
\Level{B}
\SubArea{Quadratic Functions}
\IMG{000223_06-figure0.pdf}
\LANGen{
\Question{
Using the graph of the function \(f(x)= -x^{2} - 2x + 8\) 
solve the following inequality. 
\[ 
 -x^{2} - 2x + 8\leq 5
\]}
{\(\left (-\infty ;-3\right ] \cup \left [ 1;\infty \right )\)}{\(\left (-\infty ;-4\right ] \cup \left [ 2;\infty \right )\)}{\(\left [ -3;1\right ] \)}{\(\left [ -4;2\right ] \)}{}{}}
\LANGpl{
\Question{
Wykorzystując wykres funkcji \(f\colon y = -x^{2} - 2x + 8\) 
rozwiąż podaną nierówność.
\[ 
 -x^{2} - 2x + 8\leq 5
\]}
{\(\left (-\infty ;-3\right ] \cup \left [ 1;\infty \right )\)}{\(\left (-\infty ;-4\right ] \cup \left [ 2;\infty \right )\)}{\(\left [ -3;1\right ] \)}{\(\left [ -4;2\right ] \)}{}{}}
\LANGcs{
\Question{
S využitím grafu funkce \(f\colon y = -x^{2} - 2x + 8\) 
určete řešení nerovnice.
\[ 
 -x^{2} - 2x + 8\leq 5
\]}
{\(\left (-\infty ;-3\right \rangle \cup \left \langle 1;\infty \right )\)}{\(\left (-\infty ;-4\right \rangle \cup \left \langle 2;\infty \right )\)}{\(\left \langle -3;1\right \rangle \)}{\(\left \langle -4;2\right \rangle \)}{}{}}
\LANGsk{
\Question{
S využitím grafu funkcie \(f\colon y = -x^{2} - 2x + 8\) 
vyriešte danú nerovnicu. 
\[ 
 -x^{2} - 2x + 8\leq 5
\]}
{\(\left (-\infty ;-3\right ] \cup \left [ 1;\infty \right )\)}{\(\left (-\infty ;-4\right ] \cup \left [ 2;\infty \right )\)}{\(\left [ -3;1\right ] \)}{\(\left [ -4;2\right ] \)}{}{}}
\LANGes{
\Question{
Usando la gráfica de la función \(f(x)= -x^{2} - 2x + 8\) 
resuelve la siguiente inecuación
\[ 
 -x^{2} - 2x + 8\leq 5
\]}
{\(\left (-\infty ;-3\right ] \cup \left [ 1;\infty \right )\)}{\(\left (-\infty ;-4\right ] \cup \left [ 2;\infty \right )\)}{\(\left [ -3;1\right ] \)}{\(\left [ -4;2\right ] \)}{}{}}
\End

\Data\ID{9000022307}
\Updated{2021-04-07 07:04:15}
\Level{B}
\SubArea{Quadratic Functions}
\IMG{000223_07-figure0.pdf}
\LANGen{
\Question{
Using the graph of the function \(f(x) = x^{2} - x - 6\) 
solve the system of inequalities. 
\[ 
 -4 < x^{2} - x - 6 < 0
\]}
{\((-2;-1)\cup (2;3)\)}{\((-2;3)\)}{\((-\infty ;-2)\cup (3;\infty )\)}{\((-\infty ;-1)\cup (2;\infty )\)}{}{}}
\LANGpl{
\Question{
Używając wykresu podanej funkcji  \(f\colon y = x^{2} - x - 6\) 
rozwiąż układ nierówności.
\[ 
 -4 < x^{2} - x - 6 < 0
\]}
{\((-2;-1)\cup (2;3)\)}{\((-2;3)\)}{\((-\infty ;-2)\cup (3;\infty )\)}{\((-\infty ;-1)\cup (2;\infty )\)}{}{}}
\LANGcs{
\Question{
Využijte graf funkce \(f\colon y = x^{2} - x - 6\) 
k řešení soustavy nerovnic.
\[ 
 -4 < x^{2} - x - 6 < 0
\]}
{\((-2;-1)\cup (2;3)\)}{\((-2;3)\)}{\((-\infty ;-2)\cup (3;\infty )\)}{\((-\infty ;-1)\cup (2;\infty )\)}{}{}}
\LANGsk{
\Question{
Použitím grafu funkcie \(f\colon y = x^{2} - x - 6\) 
vyriešte danú sústavu nerovníc. 
\[ 
 -4 < x^{2} - x - 6 < 0
\]}
{\((-2;-1)\cup (2;3)\)}{\((-2;3)\)}{\((-\infty ;-2)\cup (3;\infty )\)}{\((-\infty ;-1)\cup (2;\infty )\)}{}{}}
\LANGes{
\Question{
Usando la gráfica de la función \(f(x) = x^{2} - x - 6\) 
resuelve el sistema de inecuaciones
\[ 
 -4 < x^{2} - x - 6 < 0
\]}
{\((-2;-1)\cup (2;3)\)}{\((-2;3)\)}{\((-\infty ;-2)\cup (3;\infty )\)}{\((-\infty ;-1)\cup (2;\infty )\)}{}{}}
\End

\Data\ID{9000022308}
\Updated{2021-04-07 07:04:43}
\Level{B}
\SubArea{Quadratic Functions}
\IMG{000223_08-figure0.pdf}
\LANGen{
\Question{
Using graphs of the functions \(f(x)= -2x^{2} + 3x + 4\) 
and \(g(x) = x\) 
solve the following quadratic inequality. 
\[ 
 -2x^{2} + 3x + 4\geq x
\]}
{\(\left [ -1;2\right ] \)}{\(\{ - 1;2\}\)}{\(\left (-1;2\right )\)}{\(\left (-\infty ;-1\right )\cup \left (2;\infty \right )\)}{}{}}
\LANGpl{
\Question{
Wykorzystując wykresy funkcji  \(f\colon y = -2x^{2} + 3x + 4\) 
i \(g\colon y = x\) 
rozwiąż podaną nierówność kwadratową.
\[ 
 -2x^{2} + 3x + 4\geq x
\]}
{\(\left [ -1;2\right ] \)}{\(\{ - 1;2\}\)}{\(\left (-1;2\right )\)}{\(\left (-\infty ;-1\right )\cup \left (2;\infty \right )\)}{}{}}
\LANGcs{
\Question{
S využitím grafů funkcí \(f\colon y = -2x^{2} + 3x + 4\) 
a \(g\colon y = x\) 
určete řešení kvadratické nerovnice.
\[ 
 -2x^{2} + 3x + 4\geq x
\]}
{\(\left \langle -1;2\right \rangle \)}{\(\{ - 1;2\}\)}{\(\left (-1;2\right )\)}{\(\left (-\infty ;-1\right )\cup \left (2;\infty \right )\)}{}{}}
\LANGsk{
\Question{
S využitím grafov funkcií \(f\colon y = -2x^{2} + 3x + 4\) 
a \(g\colon y = x\) 
vyriešte danú kvadratickú nerovnicu. 
\[ 
 -2x^{2} + 3x + 4\geq x
\]}
{\(\left [ -1;2\right ] \)}{\(\{ - 1;2\}\)}{\(\left (-1;2\right )\)}{\(\left (-\infty ;-1\right )\cup \left (2;\infty \right )\)}{}{}}
\LANGes{
\Question{
Usando las gráficas de las funciones \(f(x)= -2x^{2} + 3x + 4\) 
y \(g(x) = x\) ,
resuelve la siguiente inecuación cuadrática
\[ 
 -2x^{2} + 3x + 4\geq x
\]}
{\(\left [ -1;2\right ] \)}{\(\{ - 1;2\}\)}{\(\left (-1;2\right )\)}{\(\left (-\infty ;-1\right )\cup \left (2;\infty \right )\)}{}{}}
\End

\Data\ID{9000022309}
\Updated{2021-04-07 07:04:08}
\Level{B}
\SubArea{Quadratic Functions}
\IMG{000223_09-figure0.pdf}
\LANGen{
\Question{
Using graphs of the functions \(f(x) = x^{2} + x - 1\) 
and \(g(x) = -\frac{1}
{2}x\) 
solve the following quadratic inequality.
\[ 
 x^{2} + x - 1 > -\frac{1}
{2}x
\]}
{\(\left (-\infty ;-2\right )\cup \left (\frac{1}
{2};\infty \right )\)}{\(\left (-2; \frac{1} 
{2}\right )\)}{\(\left [ -2; \frac{1} 
{2}\right ] \)}{\(\left (-\infty ;-2\right ] \cup \left [ \frac{1}
{2};\infty \right )\)}{}{}}
\LANGpl{
\Question{
Wykorzystując wykresy funkcji  \(f\colon y = x^{2} + x - 1\) 
i \(g\colon y = -\frac{1}
{2}x\) 
rozwiąż podaną nierówność kwadratową.
\[ 
 x^{2} + x - 1 > -\frac{1}
{2}x
\]}
{\(\left (-\infty ;-2\right )\cup \left (\frac{1}
{2};\infty \right )\)}{\(\left (-2; \frac{1} 
{2}\right )\)}{\(\left [ -2; \frac{1} 
{2}\right ] \)}{\(\left (-\infty ;-2\right ] \cup \left [ \frac{1}
{2};\infty \right )\)}{}{}}
\LANGcs{
\Question{
S využitím grafů funkcí \(f\colon y = x^{2} + x - 1\) 
a \(g\colon y = -\frac{1}
{2}x\) 
určete řešení kvadratické nerovnice.
\[ 
 x^{2} + x - 1 > -\frac{1}
{2}x
\]}
{\(\left (-\infty ;-2\right )\cup \left (\frac{1}
{2};\infty \right )\)}{\(\left (-2; \frac{1} 
{2}\right )\)}{\(\left \langle -2; \frac{1} 
{2}\right \rangle \)}{\(\left (-\infty ;-2\right \rangle \cup \left \langle \frac{1}
{2};\infty \right )\)}{}{}}
\LANGsk{
\Question{
S využitím grafov funkcií \(f\colon y = x^{2} + x - 1\) 
a \(g\colon y = -\frac{1}
{2}x\) 
 vyriešte danú nerovnicu. 
\[ 
 x^{2} + x - 1 > -\frac{1}
{2}x
\]}
{\(\left (-\infty ;-2\right )\cup \left (\frac{1}
{2};\infty \right )\)}{\(\left (-2; \frac{1} 
{2}\right )\)}{\(\left [ -2; \frac{1} 
{2}\right ] \)}{\(\left (-\infty ;-2\right ] \cup \left [ \frac{1}
{2};\infty \right )\)}{}{}}
\LANGes{
\Question{
Usando las gráficas de las funciones \(f(x) = x^{2} + x - 1\) 
y \(g(x) = -\frac{1}
{2}x\) ,
resuelve la siguiente inecuación cuadrática
\[ 
 x^{2} + x - 1 > -\frac{1}
{2}x
\]}
{\(\left (-\infty ;-2\right )\cup \left (\frac{1}
{2};\infty \right )\)}{\(\left (-2; \frac{1} 
{2}\right )\)}{\(\left [ -2; \frac{1} 
{2}\right ] \)}{\(\left (-\infty ;-2\right ] \cup \left [ \frac{1}
{2};\infty \right )\)}{}{}}
\End

\Data\ID{9000025610}
\Updated{2021-04-07 07:04:37}
\Level{B}
\SubArea{Quadratic Functions}
\IMG{000256_10-figure0.pdf}
\LANGen{
\Question{
Identify a quadratic equation which is solved by a graphical method in the
picture.}
{\(x^{2} - 6x + 9 = 0\)}{\(x^{2} + 9x - 3 = 0\)}{\(x^{2} - 9x - 3 = 0\)}{\(x^{2} + 6x + 9 = 0\)}{}{}}
\LANGpl{
\Question{
Rozwiązanie którego równania kwadratowego przedstawiono na rysunku poniżej?}
{\(x^{2} - 6x + 9 = 0\)}{\(x^{2} + 9x - 3 = 0\)}{\(x^{2} - 9x - 3 = 0\)}{\(x^{2} + 6x + 9 = 0\)}{}{}}
\LANGcs{
\Question{
Vyberte kvadratickou rovnici, jejíž grafické řešení je znázorněno na
obrázku.}
{\(x^{2} - 6x + 9 = 0\)}{\(x^{2} + 9x - 3 = 0\)}{\(x^{2} - 9x - 3 = 0\)}{\(x^{2} + 6x + 9 = 0\)}{}{}}
\LANGsk{
\Question{
Vyberte kvadratickú rovnicu, ktorej grafické riešenie je znázornené na obrázku.}
{\(x^{2} - 6x + 9 = 0\)}{\(x^{2} + 9x - 3 = 0\)}{\(x^{2} - 9x - 3 = 0\)}{\(x^{2} + 6x + 9 = 0\)}{}{}}
\LANGes{
\Question{
¿Cuál de las ecuaciones cuadráticas se puede resolver mediante la gráfica del dibujo?}
{\(x^{2} - 6x + 9 = 0\)}{\(x^{2} + 9x - 3 = 0\)}{\(x^{2} - 9x - 3 = 0\)}{\(x^{2} + 6x + 9 = 0\)}{}{}}
\End

\Data\ID{9100025608}
\Updated{2021-04-07 07:04:15}
\Level{B}
\SubArea{Quadratic Functions}
\LANGen{
\Question{
Identify the picture which shows solution of
\[ 
 \text{$2x^{2} - 5x - 3 = 0$}
\]
 by a graphical method.}
{}{}{}{}{}{}}
\LANGpl{
\Question{
Który z poniższych rysunków jest graficznym rozwiązaniem równania?
\[ 
 \text{$2x^{2} - 5x - 3 = 0$}
\]}
{}{}{}{}{}{}}
\LANGcs{
\Question{
Na kterém obrázku je znázorněn graf funkce \(f\colon y= 2x^2-5x-3\) společně s řešením kvadratické rovnice \(2x^2-5x-3=0\)?}
{}{}{}{}{}{}}
\LANGsk{
\Question{
Na ktorom obrázku je zobrazené grafické riešenie  danej rovnice?
\[ 
 \text{$2x^{2} - 5x - 3 = 0$}
\]}
{}{}{}{}{}{}}
\LANGes{
\Question{
¿Cuál de los dibujos representa gráficamente la solución de la siguiente ecuación?
\[ 
 \text{$2x^{2} - 5x - 3 = 0$}
\]}
{}{}{}{}{}{}}

\IMGanswer{000256_08-figure2.pdf}
\IMGanswer{000256_08-figure0.pdf}
\IMGanswer{000256_08-figure1.pdf}
\IMGanswer{000256_08-figure3.pdf}\End

\Data\ID{9100033706}
\Updated{2020-11-20 12:11:13}
\Level{B}
\SubArea{Quadratic Functions}
\LANGen{
\Question{
Identify a picture which shows graphical solution of the following quadratic
inequality, assuming that the solution is the red set. 
\[ 
 -x^{2} + x + 2 > 2x
\]}
{}{}{}{}{}{}}
\LANGpl{
\Question{
Wybierz obraz, który przedstawia graficzne rozwiązanie podanej nierówności kwadratowej, zakładając, że rozwiązaniem jest zbiór czerwony.
\[ 
 -x^{2} + x + 2 > 2x
\]}
{}{}{}{}{}{}}
\LANGcs{
\Question{
Na kterém z obrázků je červenou barvou ilustrováno grafické řešení nerovnice
\(-x^{2} + x + 2 > 2x\)?}
{}{}{}{}{}{}}
\LANGsk{
\Question{
Na ktorom obrázku je zobrazené červenou farbou grafické riešenie danej nerovnice? 
\[ 
 -x^{2} + x + 2 > 2x
\]}
{}{}{}{}{}{}}
\LANGes{
\Question{
¿Cuál de los dibujos representa la solución gráfica de la siguiente inecuación cuadrática, suponiendo que la solución es el conjunto rojo?
\[ 
 -x^{2} + x + 2 > 2x
\]}
{}{}{}{}{}{}}

\IMGanswer{000337_06-figure2.pdf}
\IMGanswer{000337_06-figure0.pdf}
\IMGanswer{000337_06-figure1.pdf}
\IMGanswer{000337_06-figure3.pdf}
\IMGanswer{000337_06-figure4.pdf}
\IMGanswer{000337_06-figure5.pdf}\End

\Data\ID{1003083108}
\Updated{2021-02-15 07:02:30}
\Level{C}
\SubArea{Quadratic Functions}
\LANGen{
\Question{
The parabolas of the functions \( f \) and \( g \) have the same vertex \( V \) and \( f(x)=ax^2+c \), where \( a \) and \( c \) are nonzero real numbers. Find \( g(x) \) such that the graphs of \( f \) and \( g \) are symmetric about  the vertex \( V \) and that  \( y \)-axis is their line of symmetry.}
{\( g(x)=-ax^2+c\), i.e. the equations of \( f \) and \( g \) differ in the sign of the coefficient at the quadratic term only}{\( g(x)=ax^2-c\), i.e. the equations of \( f \) and \( g \) differ in the sign of the coefficient at the linear term only}{\( g(x)=-ax^2-c \), i.e. \( g(x)=-f(x) \)}{None of the statements above is true.}{}{}}
\LANGpl{
\Question{
Parabole funkcji \( f \) i \( g \) mają ten sam wierzchołek \( V \) i \( f(x)=ax^2+c \), gdzie \( a \) i \( c \) są niezerowymi liczbami rzeczywistymi. Wyznacz \( g(x) \) wykresy funkcji  \( f \) i \( g \) są symetryczne względem wierzchołka \( V \) a oś \( y \) jest ich linią symetrii.}
{\( g(x)=-ax^2+c\), tzn. wzory funkcji  \( f \) i \( g \) różnią się znakiem współczynnika tylko w stosunku kwadratowym}{\( g(x)=ax^2-c\),  tzn. wzory funkcji \( f \) i \( g \) różnią się znakiem współczynnika tylko w stosunku liniowym}{\( g(x)=-ax^2-c \), tj.  \( g(x)=-f(x) \)}{Żadne z powyższych stwierdzeń nie jest prawdziwe.}{}{}}
\LANGcs{
\Question{
Grafy funkcí \( f \) a \( g \) jsou paraboly se společným vrcholem \( V \) a \( f(x)=ax^2+c \), kde \( a \) a \( c \) jsou nenulová reálná čísla. Najděte funkci \( g \) tak, aby grafy funkcí \( f \) a \( g \) byly středově souměrné podle vrcholu \( V \) a byly symetrické podle osy \( y \).}
{\( g(x)=-ax^2+c\), tj. předpisy funkcí \( f \) a \( g \) se liší pouze znaménkem koeficientu u kvadratického členu}{\( g(x)=ax^2-c\), tj. předpisy funkcí \( f \) a \( g \) se liší pouze znaménkem koeficientu u lineárního členu}{\( g(x)=-ax^2-c \), tj. \( g(x)=-f(x) \)}{Žádné z výše uvedených tvrzení není pravdivé.}{}{}}
\LANGsk{
\Question{
Dané sú grafy kvadratických funkcií \( f \) a \( g \), ktoré majú spoločný vrchol \( V \). Daný je predpis funkcie \( f(x)=ax^2+c \), kde \( a \) a \( c \) sú nenulové reálne čísla. Určte predpis funkcie \( g \) tak, aby grafy funkcií \( f \) a \( g \) boli stredovo súmerné podľa vrcholu \( V \) a zároveň osovo súmerné podľa osi \( y \).}
{\( g(x)=-ax^2+c\), tj. predpisy funkcií \( f \) a \( g \) sa líšia len znamienkom koeficientu kvadratického člena}{\( g(x)=ax^2-c\), tj. predpisy funkcií \( f \) a \( g \) sa líšia len znamienkom absolútneho člena}{\( g(x)=-ax^2-c \), tj. \( g(x)=-f(x) \)}{Žiadne tvrdenie nie je pravdivé.}{}{}}
\LANGes{
\Question{
Las parábolas de las funciones \( f \) y \( g \) tienen el mismo vértice \( V \) y \( f(x)=ax^2+c \), dónde \( a \) y \( c \) son números reales distintos de cero. Halla \( g(x) \) de manera que las gráficas de \( f \) y \( g \) sean simétricas respecto al vértice \( V \) y que el eje \( y \) sea su eje de simetría.}
{\( g(x)=-ax^2+c\), es decir las ecuaciones de \( f \) y \( g \) difieren solamente en el signo del coeficiente del término de grado 2.}{\( g(x)=ax^2-c\),es decir las ecuaciones de \( f \) y \( g \) difieren solamente en el signo del coeficiente del término de grado 1.}{\( g(x)=-ax^2-c \), es decir \( g(x)=-f(x) \)}{Ninguna de las afirmaciones es correcta.}{}{}}
\End

\Data\ID{1003083110}
\Updated{2021-02-15 07:02:20}
\Level{C}
\SubArea{Quadratic Functions}
\LANGen{
\Question{
The graphs of the quadratic functions \( f \) and \( g \) have not the same vertex and \( f(x)=ax^2+bx+c \), where \( a \), \( b \), \( c \) are nonzero real numbers. Find \( g(x) \) such that the graph of \( g \) is the reflection of the graph of \( f \) about \( y \)-axis.}
{\( g(x)=ax^2-bx+c \), i.e. the equation of \( f \) and \( g \) differ in the sign of the coefficient at the linear term only}{\( g(x)=-ax^2+bx+c \), i.e  the equation of \( f \) and \( g \) differ in the sign of the coefficient at the quadratic term only}{\( g(x)=ax^2+bx-c \), i.e. the equation of \( f \) and \( g \) differ in the sign of the coefficient at the absolute term  only}{\( g(x)=-ax^2-bx-c \), i.e. \( g(x)=-f(x) \)}{None of the statements above is true.}{}}
\LANGpl{
\Question{
Wykresy funkcji kwadratowych \( f \) i \( g \) nie mają tego samego wierzchołka, a \( f(x)=ax^2+bx+c \), gdzie \( a \), \( b \), \( c \) są niezerowymi liczbami rzeczywistymi.  Wyznacz \( g(x) \) takie, dla którego \( g \) jest odbiciem wykresu funkcji  \( f \) względem osi \( y \).}
{\( g(x)=ax^2-bx+c \), tj. wzory funkcji \( f \) i \( g \) różnią się znakiem współczynnika tylko w stosunku liniowym}{\( g(x)=-ax^2+bx+c \), tj. wzory funkcji \( f \) i \( g \) różnią się znakiem współczynnika tylko w stosunku kwadratowym}{\( g(x)=ax^2+bx-c \),  \( f \) i \( g \) różnią się znakiem współczynnika tylko na poziomie bezwzględnym}{\( g(x)=-ax^2-bx-c \), tj. \( g(x)=-f(x) \)}{Żadne z powyższych stwierdzeń nie jest prawdziwe.}{}}
\LANGcs{
\Question{
Grafy funkcí \( f \) a \( g \) jsou paraboly s různými vrcholy a \( f(x)=ax^2+bx+c \), kde \( a \), \( b \), \( c \) jsou nenulová reálná čísla. Najděte funkci \( g(x) \) tak, aby graf  \( g \) byl obrazem grafu \( f \) v osové souměrnosti dané osou \( y \).}
{\( g(x)=ax^2-bx+c \), tj. předpisy funkcí  \( f \) a \( g \) se liší pouze znaménkem koeficientu u lineárního členu}{\( g(x)=-ax^2+bx+c \), tj. předpisy funkcí  \( f \) a \( g \) se liší pouze znaménkem koeficientu u kvadratického členu}{\( g(x)=ax^2+bx-c \), tj. předpisy funkcí  \( f \) a \( g \) se liší pouze znaménkem koeficientu u absolutního členu}{\( g(x)=-ax^2-bx-c \), tj. \( g(x)=-f(x) \)}{Žádná z možností není správně.}{}}
\LANGsk{
\Question{
Dané sú grafy kvadratických funkcií \( f \) a \( g \), ktoré nemajú spoločný vrchol \( V \) a predpis funkcie \( f(x)=ax^2+bx+c \), kde \( a \), \( b \), \( c \) sú nenulové reálne čísla. Určte predpis funkcie \( g(x) \) tak, aby bol graf funkcie \( g \) obrazom grafu funkcie \( f \) v osovej súmernosti podľa osi \( y \).}
{\( g(x)=ax^2-bx+c \), tj. predpisy funkcií \( f \) a \( g \) sa líšia len znamienkom koeficientu lineárneho člena}{\( g(x)=-ax^2+bx+c \), tj. predpisy funkcií \( f \) a \( g \) sa líšia len znamienkom koeficientu kvadratického člena}{\( g(x)=ax^2+bx-c \), tj. predpisy funkcií \( f \) a \( g \) sa líšia len znamienkom absolútneho člena}{\( g(x)=-ax^2-bx-c \), t.j. \( g(x)=-f(x) \)}{Žiadne tvrdenie nie je pravdivé.}{}}
\LANGes{
\Question{
Las gráficas de las funciones cuadráticas \( f \) y \( g \) no tienen el mismo vértice y \( f(x)=ax^2+bx+c \), dónde \( a \), \( b \), \( c \) son números reales distintos de cero. Encuentra \( g(x) \) tal que la gráfica de \( g \) sea la reflexión de la gráfica de \( f \) respecto al eje \( y \).}
{\( g(x)=ax^2-bx+c \), es decir las ecuaciones de \( f \) y \( g \) difieren solamente en el signo del coeficiente del término de grado 1.}{\( g(x)=-ax^2+bx+c \),es decir las ecuaciones de \( f \) y \( g \) difieren solamente en el signo del coeficiente del término de grado 2.}{\( g(x)=ax^2+bx-c \), es decir las ecuaciones de \( f \) y \( g \) difieren solamente en el signo del coeficiente del término independiente.}{\( g(x)=-ax^2-bx-c \), es decir \( g(x)=-f(x) \)}{Ninguna de las afirmaciones es correcta.}{}}
\End

\Data\ID{1003108307}
\Updated{2021-02-15 07:02:18}
\Level{C}
\SubArea{Quadratic Functions}
\LANGen{
\Question{
Choose the triple of points, such that  the graph of any of the functions \( f(x)=ax^2+c \), where \( a\in\mathbb{R}\setminus{0} \), \( c\in\mathbb{R} \), does not pass through all three points.}
{\( [-2;5] \), \( [2;1] \), \( [0;3] \)}{\( [-2;5] \), \( [2;5] \), \( [0;3] \)}{\( [-2;5] \), \( [2;5] \), \( [0;7] \)}{\( [-2;5] \), \( [0;0] \), \( [1;1] \)}{}{}}
\LANGpl{
\Question{
Wskaż trójki punktów, takich, że żaden z wykresów funkcji  \( f(x)=ax^2+c \), gdzie \( a\in\mathbb{R}\setminus{0} \), \( c\in\mathbb{R} \), nie przechodzi przez te trzy punkty.}
{\( [-2;5] \), \( [2;1] \), \( [0;3] \)}{\( [-2;5] \), \( [2;5] \), \( [0;3] \)}{\( [-2;5] \), \( [2;5] \), \( [0;7] \)}{\( [-2;5] \), \( [0;0] \), \( [1;1] \)}{}{}}
\LANGcs{
\Question{
Vyberte trojici bodů, kterými nemůže procházet graf funkce \( f(x)=ax^2+c \), kde \( a\in\mathbb{R}\setminus{0} \), \( c\in\mathbb{R} \).}
{\( [-2;5] \), \( [2;1] \), \( [0;3] \)}{\( [-2;5] \), \( [2;5] \), \( [0;3] \)}{\( [-2;5] \), \( [2;5] \), \( [0;7] \)}{\( [-2;5] \), \( [0;0] \), \( [1;1] \)}{}{}}
\LANGsk{
\Question{
Vyberte trojicu bodov, ktorými nemôže prechádzať graf kvadratickej funkcie \( f(x)=ax^2+c \), kde \( a\in\mathbb{R}\setminus{0} \), \( c\in\mathbb{R} \).}
{\( [-2;5] \), \( [2;1] \), \( [0;3] \)}{\( [-2;5] \), \( [2;5] \), \( [0;3] \)}{\( [-2;5] \), \( [2;5] \), \( [0;7] \)}{\( [-2;5] \), \( [0;0] \), \( [1;1] \)}{}{}}
\LANGes{
\Question{
Elije los tres puntos para los cuáles ninguna de las funciones \( f(x)=ax^2+c \), dónde \( a\in\mathbb{R}\setminus{0} \), \( c\in\mathbb{R} \),  pasa por los tres puntos.}
{\( [-2;5] \), \( [2;1] \), \( [0;3] \)}{\( [-2;5] \), \( [2;5] \), \( [0;3] \)}{\( [-2;5] \), \( [2;5] \), \( [0;7] \)}{\( [-2;5] \), \( [0;0] \), \( [1;1] \)}{}{}}
\End

\Data\ID{1003124801}
\Updated{2021-04-07 08:04:30}
\Level{C}
\SubArea{Quadratic Functions}
\LANGen{
\Question{
Suppose we want to paint a cube so that there remains an unpainted stripe along all the edges on each face. The width of the stripe should be \( 1\,\mathrm{cm} \). The producer gives the paint consumption \( 100\,\mathrm{ml}/1\,\mathrm{m}^2 \). From the following functions choose the one that describes the dependence of the paint consumption \( V \) on the length of the cube edge \( a \). The paint consumption \( V \) is given in millilitres and the length of the cube edge \( a \) is given in meters.}
{\( V=\left(a-\frac1{50}\right)^2\cdot600 \)}{\( V=\left(a-\frac1{50}\right)^2\cdot\frac3{50} \)}{\( V=\left(a-\frac1{100}\right)^2\cdot600 \)}{\( V=(a-2)^2\cdot100 \)}{}{}}
\LANGpl{
\Question{
Załóżmy, że chcemy pomalować sześcian, tak, że na każdej ścianie wzdłuż brzegów pozostawimy niepomalowany pasek. Szerokość paska powinna wynosić  \( 1\,\mathrm{cm} \). Producent twierdzi, że wydajność farby wynosi \( 100\,\mathrm{ml}/1\,\mathrm{m}^2 \).  Z poniższych funkcji wybierz tę, która opisuje zależność pomiędzy wydajnością farby \( V \), a długością boków sześcianu \( a \). Wydajność farby \( V \) jest wyrażona w mililitrach a długość boku \( a \) metrach.}
{\( V=\left(a-\frac1{50}\right)^2\cdot600 \)}{\( V=\left(a-\frac1{50}\right)^2\cdot\frac3{50} \)}{\( V=\left(a-\frac1{100}\right)^2\cdot600 \)}{\( V=(a-2)^2\cdot100 \)}{}{}}
\LANGcs{
\Question{
Potřebujeme natřít těleso tvaru krychle tak, aby každá stěna měla po obvodu nenatřený pruh široký \( 1\,\mathrm{cm} \). Výrobce uvádí spotřebu barvy \( 100\,\mathrm{ml}/1\,\mathrm{m}^2 \). Z následujících možností vyberte funkci, která vyjadřuje spotřebu barvy v závislosti na velikosti hrany krychle. Spotřebu barvy v mililitrech označte \( V \) a velikost hrany krychle v metrech označte \( a \).}
{\( V=\left(a-\frac1{50}\right)^2\cdot600 \)}{\( V=\left(a-\frac1{50}\right)^2\cdot\frac3{50} \)}{\( V=\left(a-\frac1{100}\right)^2\cdot600 \)}{\( V=(a-2)^2\cdot100 \)}{}{}}
\LANGsk{
\Question{
Potrebujeme natrieť teleso tvaru kocky tak, aby každá stena mala po obvode nenatretý pruh široký \( 1\,\mathrm{cm} \). Výrobca uvádza spotrebu farby \( 100\,\mathrm{ml}/1\,\mathrm{m}^2 \). Z následujúcich možností vyberte funkciu, ktorá vyjadruje spotrebu farby v závislosti na velikosti hrany kocky. Spotrebu farby v mililitroch označte \( V \) a veľkosť hrany kocky v metroch označte \( a \).}
{\( V=\left(a-\frac1{50}\right)^2\cdot600 \)}{\( V=\left(a-\frac1{50}\right)^2\cdot\frac3{50} \)}{\( V=\left(a-\frac1{100}\right)^2\cdot600 \)}{\( V=(a-2)^2\cdot100 \)}{}{}}
\LANGes{
\Question{
Supongamos que queremos pintar un cubo de manera que en cada cara quede una banda sin pintar. La anchura de cada banda debería ser \( 1\,\mathrm{cm} \). La cantidad de pintura necesaria es de \( 100\,\mathrm{ml}/1\,\mathrm{m}^2 \). De las siguientes funciones elije la que describe la dependencia de la cantidad de pintura necesaria \( V \) respecto a la longitud del lado del cubo \( a \). La cantidad de pintura  \( V \) viene dada en milímetros y la longitud del lado del cubo \( a \) en metros.}
{\( V=\left(a-\frac1{50}\right)^2\cdot600 \)}{\( V=\left(a-\frac1{50}\right)^2\cdot\frac3{50} \)}{\( V=\left(a-\frac1{100}\right)^2\cdot600 \)}{\( V=(a-2)^2\cdot100 \)}{}{}}
\End

\Data\ID{1003124802}
\Updated{2021-04-05 08:04:48}
\Level{C}
\SubArea{Quadratic Functions}
\LANGen{
\Question{
We want to plant flowers into rectangular flower bed with longer side by one meter longer than its shorter side. Each flower needs \( 1\,\mathrm{dm}^2 \) of free space.  From the following functions, choose the one that describes the dependence of the number of planted flowers \( n \) on the length \( a \) of the shorter side of the flower bed.  (Assume that the dimensions of the flower bed are given in whole meters.)}
{\( n=\left(a^2+a\right)\cdot100 \)}{\( n=\left(a^2+a\right)\cdot\frac1{100} \)}{\( n=(a+1)^2\cdot100 \)}{\( n=\left(a^2+a\right) \)}{}{}}
\LANGpl{
\Question{
Chcemy zasadzić kwiaty w donicy w kształcie prostokąta o dłuższym boku, który jest o metr dłuższy od jego krótszego boku. Każdy kwiatek potrzebuje \( 1\,\mathrm{dm}^2 \) wolnej przestrzeni. Z poniższych funkcji wybierz jedną, która określa zależność pomiędzy liczbą posadzonych kwiatów \( n \) a długością krótszego boku donicy \( a \). (Załóż, że wymiary donicy podano w pełnych metrach.)}
{\( n=\left(a^2+a\right)\cdot100 \)}{\( n=\left(a^2+a\right)\cdot\frac1{100} \)}{\( n=(a+1)^2\cdot100 \)}{\( n=\left(a^2+a\right) \)}{}{}}
\LANGcs{
\Question{
Sazenicemi rostlin chceme osadit záhon tvaru obdélníku, jehož delší strana je o \( 1\,\mathrm{m} \) delší než jeho kratší strana. Každá sazenice potřebuje 
 \( 1\,\mathrm{dm}^2 \) volné plochy. Z následujících možností vyberte funkci, která vyjadřuje závislost počtu sazenic \( n \) na délce \( a \) kratší strany záhonu. (Poznámka: Rozměry záhonu jsou v celých metrech.)}
{\( n=\left(a^2+a\right)\cdot100 \)}{\( n=\left(a^2+a\right)\cdot\frac1{100} \)}{\( n=(a+1)^2\cdot100 \)}{\( n=\left(a^2+a\right) \)}{}{}}
\LANGsk{
\Question{
Záhon tvaru obdĺžnika chceme vysadiť sadenicami rastlín. Dlhšia strana obdĺžnika má o \( 1\,\mathrm{m} \) viac ako kratšia. Každá sadenica potrebuje 
 \( 1\,\mathrm{dm}^2 \) voľnej plochy. Z následujúcich možností vyberte funkciu, ktorá vyjadruje závislosť počtu sadeníc \( n \) na dĺžke kratšej strany záhonu \( a \). (Poznámka: Rozmery záhonu sú v celých metroch.)}
{\( n=\left(a^2+a\right)\cdot100 \)}{\( n=\left(a^2+a\right)\cdot\frac1{100} \)}{\( n=(a+1)^2\cdot100 \)}{\( n=\left(a^2+a\right) \)}{}{}}
\LANGes{
\Question{
Queremos plantar flores en un florero rectangular cuyo lado largo tiene un metro más que el lado corto. Cada flor necesita \( 1\,\mathrm{dm}^2 \) de espacio. De las funciones siguientes elije la que describe la dependencia del número de flores plantadas \( n \) respecto a la longitud del lado más corto del florero \( a \).  (Supón que las dimensiones del florero vienen dadas en metros enteros)}
{\( n=\left(a^2+a\right)\cdot100 \)}{\( n=\left(a^2+a\right)\cdot\frac1{100} \)}{\( n=(a+1)^2\cdot100 \)}{\( n=\left(a^2+a\right) \)}{}{}}
\End

\Data\ID{1003124803}
\Updated{2021-04-07 08:04:16}
\Level{C}
\SubArea{Quadratic Functions}
\LANGen{
\Question{
The annulus shaped component parts are punched from sheet metal. Diameter of the circular hole is \( 25\,\% \) of the diameter of the whole component part. Choose the function that describes the dependence of the area (\( S \)) of material used to produce one component part on its outside diameter (\( d \)).}
{\( S=\frac{15}{64}\,\pi d^2 \)}{\( S=\frac38\,\pi d^2 \)}{\( S=\frac{15}{32}\,\pi d^2 \)}{\( S=\frac{31}{64}\,\pi d^2 \)}{}{}}
\LANGpl{
\Question{
Części w kształcie pierścienia są dziurkowane z blachy. Średnica okrągłego otworu wynosi \( 25\,\% \)  średnicy całej części składowej. Wybierz funkcję opisującą zależność obszaru  (\( S \))  od materiału użytego do wytworzenia jednej części składowej na jej średnicy zewnętrznej  (\( d \)).}
{\( S=\frac{15}{64}\,\pi d^2 \)}{\( S=\frac38\,\pi d^2 \)}{\( S=\frac{15}{32}\,\pi d^2 \)}{\( S=\frac{31}{64}\,\pi d^2 \)}{}{}}
\LANGcs{
\Question{
Z plechu razíme součástky tvaru mezikruží, přičemž průměr kruhového otvoru je \( 25\,\% \) průměru celé součástky. Z následujících možností vyberte funkci, která vyjadřuje závislost plochy (\( S \)) materiálu spotřebovaného při výrobě součástky na jejím vnějším průměru (\( d \)).}
{\( S=\frac{15}{64}\,\pi d^2 \)}{\( S=\frac38\,\pi d^2 \)}{\( S=\frac{15}{32}\,\pi d^2 \)}{\( S=\frac{31}{64}\,\pi d^2 \)}{}{}}
\LANGsk{
\Question{
Súčiastky tvaru medzikružia razíme z plechu. Priemer kruhového otvoru je \( 25\,\% \) priemeru celej súčiastky. Z následujúcich možností vyberte funkciu, ktorá vyjadruje závislosť plochy (\( S \)) materiálu spotrebovaného pri výrobe súčiastky na jej vonkajšom priemere (\( d \)).}
{\( S=\frac{15}{64}\,\pi d^2 \)}{\( S=\frac38\,\pi d^2 \)}{\( S=\frac{15}{32}\,\pi d^2 \)}{\( S=\frac{31}{64}\,\pi d^2 \)}{}{}}
\LANGes{
\Question{
Un componente en forma de anillo está hecho de un metal. El diámetro del agujero es  \( 25\,\% \)  del diámetro de todo el componente. Elije la función que describe la dependencia del área (\( S \)) del material usado para producir el componente respecto al diámetro exterior (\( d \)).}
{\( S=\frac{15}{64}\,\pi d^2 \)}{\( S=\frac38\,\pi d^2 \)}{\( S=\frac{15}{32}\,\pi d^2 \)}{\( S=\frac{31}{64}\,\pi d^2 \)}{}{}}
\End

\Data\ID{1003124804}
\Updated{2021-04-07 08:04:20}
\Level{C}
\SubArea{Quadratic Functions}
\LANGen{
\Question{
In the centre of a square shaped square there is a water fountain. The fountain has a square ground plan with the side length \( 4.5\,\mathrm{m} \). The square should be paved with cobblestones of size \( 25\,\mathrm{cm} \times 25\,\mathrm{cm} \). Choose the function that describes the dependence of the number of cobblestones needed (\( n \)) on the length of the square (\( a \)) given in meters.}
{\( n=16a^2-324 \)}{\( n=\frac{a^2}{625}-324 \)}{\( n=16a^2-625 \)}{\( n=\frac{a^2}{16}-324 \)}{}{}}
\LANGpl{
\Question{
W centrum kwadratowego rynku znajduje się fontanna. Fontanna ma kwadratowy plan o długości boku \(4 {,}5\,\mathrm {m} \). Kwadrat powinien być wyłożony kostką brukową o wymiarach \( 25\,\mathrm{cm} \times 25\,\mathrm{cm} \).  Wybierz funkcję opisującą zależność liczby potrzebnych kostek brukowych (\(n\)) od długości kwadratu (\(a\)) podanej w metrach.}
{\( n=16a^2-324 \)}{\( n=\frac{a^2}{625}-324 \)}{\( n=16a^2-625 \)}{\( n=\frac{a^2}{16}-324 \)}{}{}}
\LANGcs{
\Question{
Uprostřed náměstí čtvercového tvaru stojí kašna. Kašna má čtvercový půdorys se stranou \( 4{,}5\,\mathrm{m} \). Náměstí má být vydlážděno dlaždicemi o rozměru \( 25\,\mathrm{cm} \times 25\,\mathrm{cm} \). Z následujících možností vyberte funkci, která vyjadřuje závislost počtu dlaždic (\( n \)) na délce náměstí (\( a \)) udávané v celých metrech.}
{\( n=16a^2-324 \)}{\( n=\frac{a^2}{625}-324 \)}{\( n=16a^2-625 \)}{\( n=\frac{a^2}{16}-324 \)}{}{}}
\LANGsk{
\Question{
V strede štvorcového námestia stojí fontána. Fontána má štvorcový pôdorys s dĺžkou strany \( 4{,}5\,\mathrm{m} \). Námestie by mali vydláždiť dlaždicami s rozmermi \( 25\,\mathrm{cm} \times 25\,\mathrm{cm} \). Vyberte funkciu, ktorá popisuje závislosť počtu dlaždíc (\( n \)) na dĺžke strany námestia (\( a \)) v metroch.}
{\( n=16a^2-324 \)}{\( n=\frac{a^2}{625}-324 \)}{\( n=16a^2-625 \)}{\( n=\frac{a^2}{16}-324 \)}{}{}}
\LANGes{
\Question{
En el centro de una plaza cuadrada hay una fuente. La fuente tiene una base cuadrada cuyo lado es de \( 4.5\,\mathrm{m} \). Laplaza debería ser pavimentada con ladrillos de tamaño  \( 25\,\mathrm{cm} \times 25\,\mathrm{cm} \). Elije la función que describe la dependencia del número de ladrillos necesarios (\( n \)) respecto a la longitud de la plaza (\( a \)) dada en metros.}
{\( n=16a^2-324 \)}{\( n=\frac{a^2}{625}-324 \)}{\( n=16a^2-625 \)}{\( n=\frac{a^2}{16}-324 \)}{}{}}
\End

\Data\ID{1003124805}
\Updated{2022-01-22 07:01:48}
\Level{C}
\SubArea{Quadratic Functions}
\LANGen{
\Question{
On a spool of mass \( 0.5\,\mathrm{kg} \) is winded an aluminium wire of length \( 100\,\mathrm{m} \). Choose the function that describes the dependence of a mass of the spool with the wire \( m \) (in kilograms) on a diameter of the wire \( d \) (in millimetres). Wire density is \( 2\,700\frac{kg}{m^3} \).
\[ \]
Hint: The density of an object is defined as the ratio of the mass and the volume of the object.}
{\( m=\frac{27\pi}{400} d^2+0.5 \)}{\( m= 67 500\pi d^2+0.5 \)}{\( m=\frac{27\pi}{400} d^2-0.5 \)}{\( m=\frac{27\pi}{200} d^2+0.5 \)}{}{}}
\LANGpl{
\Question{
Na szpuli o masie  \( 0{,}5\,\mathrm{kg} \) zwijany jest aluminiowy drut o długości \( 100\,\mathrm{m} \). Wybierz funkcję opisującą zależność masy szpuli z drutem \( m \)  (w kilogramach) od średnicy drutu \( d \) (w milimetrach). Gęstość drutu wynosi \( 2\,700\frac{kg}{m^3} \).
Wskazówka: Gęstość obiektu definiowana jest jako stosunek masy do objętości obiektu.}
{\( m=\frac{27\pi}{400} d^2+0{,}5 \)}{\( m= 67 500\pi d^2+0{,}5 \)}{\( m=\frac{27\pi}{400} d^2-0{,}5 \)}{\( m=\frac{27\pi}{200} d^2+0{,}5 \)}{}{}}
\LANGcs{
\Question{
Hliníkový drát o délce \( 100\,\mathrm{m} \) je navinutý na cívce o hmotnosti \( 0{,}5\,\mathrm{kg} \). Z následujících možností vyberte funkci, která vyjadřuje závislost hmotnosti cívky s drátem \( m \) (v kilogramech) na průměru drátu \( d \) (v milimetrech). Hustota drátu (hliníku) je \( 2\,700\frac{kg}{m^3} \).
\[ \]
Nápověda: Hustotu látky lze vypočítat jako poměr hmotnosti a objemu stejnorodého tělesa z této látky.}
{\( m=\frac{27\pi}{400} d^2+0{,}5 \)}{\( m= 67 500\pi d^2+0{,}5 \)}{\( m=\frac{27\pi}{400} d^2-0{,}5 \)}{\( m=\frac{27\pi}{200} d^2+0{,}5 \)}{}{}}
\LANGsk{
\Question{
Hliníkový drôt s dĺžkou \( 100\,\mathrm{m} \) je navinutý na cievke s hmotnosťou \( 0{,}5\,\mathrm{kg} \). Z následujúcich možností vyberte funkciu, ktorá vyjadruje závislosť hmotnosti cievky s drôtom \( m \) (v kilogramoch) na priemere drôtu \( d \) (v milimetroch). Hustota drôtu (hliníku) je \( 2\,700\frac{kg}{m^3} \).
\[ \]
Pomôcka: Hustotu počítame ako podiel hmotnosti a objemu telesa.}
{\( m=\frac{27\pi}{400} d^2+0{,}5 \)}{\( m= 67 500\pi d^2+0{,}5 \)}{\( m=\frac{27\pi}{400} d^2-0{,}5 \)}{\( m=\frac{27\pi}{200} d^2+0{,}5 \)}{}{}}
\LANGes{
\Question{
Una bobina de \( 0.5\,\mathrm{kg} \) de masa está  rodeada por un alambre de aluminio de longitud de \( 100\,\mathrm{m} \). Elige la función que describe la dependencia de la masa de la bobina con el alambre  \( m \) (en kilos) respecto al diámetro del alambre \( d \) (en milímetros). La densidad del alambre es \( 2\,700\frac{kg}{m^3} \).
\[ \] Pista: La densidad de un objeto se define como la proporción entre la masa y el volumen del objeto,}
{\( m=\frac{27\pi}{400} d^2+0.5 \)}{\( m= 67 500\pi d^2+0.5 \)}{\( m=\frac{27\pi}{400} d^2-0.5 \)}{\( m=\frac{27\pi}{200} d^2+0.5 \)}{}{}}
\End

\Data\ID{1003124806}
\Updated{2021-04-07 08:04:01}
\Level{C}
\SubArea{Quadratic Functions}
\LANGen{
\Question{
We should fence the land in a shape of an equilateral triangle. Choose the function that describes the dependence of the fenced land area \( S \) (in square meters) on the length \( d \) (in meters) of the fence used.}
{\( S=\frac{\sqrt3}{36} d^2 \)}{\( S=\frac{\sqrt3}{18} d^2 \)}{\( S=\frac{\sqrt3}4 d^2 \)}{\( S=\frac1{36} d^2 \)}{}{}}
\LANGpl{
\Question{
Należy ogrodzić pole w kształcie trójkąta równobocznego. Wybierz funkcję, która przedstawia zależność ogrodzonej ziemi \( S \) (w metrach kwadratowych) od długości \( d \) (w metrach) użytego ogrodzenia.}
{\( S=\frac{\sqrt3}{36} d^2 \)}{\( S=\frac{\sqrt3}{18} d^2 \)}{\( S=\frac{\sqrt3}4 d^2 \)}{\( S=\frac1{36} d^2 \)}{}{}}
\LANGcs{
\Question{
Plotem o délce \( d \) (v metrech) máme ohradit pozemek tvaru rovnostranného trojúhelníka. Z následujících možností vyberte funkci, která vyjadřuje závislost výměry ohrazeného pozemku \( S \) (v metrech čtverečních) na délce použitého plotu.}
{\( S=\frac{\sqrt3}{36} d^2 \)}{\( S=\frac{\sqrt3}{18} d^2 \)}{\( S=\frac{\sqrt3}4 d^2 \)}{\( S=\frac1{36} d^2 \)}{}{}}
\LANGsk{
\Question{
Pozemok tvaru rovnostranného trojuholníka ohradíme plotom s dĺžkou \( d \) (v metroch). Z následujúcich možností vyberte funkciu, ktorá vyjadruje závislosť výmery ohradeného pozemku \( S \) (v metroch štvorcových) na dĺžke použitého plotu.}
{\( S=\frac{\sqrt3}{36} d^2 \)}{\( S=\frac{\sqrt3}{18} d^2 \)}{\( S=\frac{\sqrt3}4 d^2 \)}{\( S=\frac1{36} d^2 \)}{}{}}
\LANGes{
\Question{
Necesitamos cercar un trozo de tierra en forma de triángulo equilátero. Elige la función que determina la dependencia del área de tierra cercado \( S \) (en metros cuadrados) respecto a la longitud \( d \) (en metros) de la cerca usada.}
{\( S=\frac{\sqrt3}{36} d^2 \)}{\( S=\frac{\sqrt3}{18} d^2 \)}{\( S=\frac{\sqrt3}4 d^2 \)}{\( S=\frac1{36} d^2 \)}{}{}}
\End

\Data\ID{1003148601}
\Updated{2021-04-07 08:04:58}
\Level{C}
\SubArea{Quadratic Functions}
\LANGen{
\Question{
Consider an object thrown upwards from the ground with the initial velocity of \( 30\frac{\mathrm{m}}{\mathrm{s}} \). The object moves upwards with decreasing vertical velocity until it stops. Then it starts moving vertically downwards. Find the greatest height above the ground the object does reach.  
\[ \]
Note: The vertical distance \( y \) of a thrown object is described by the equation \( y=v_0t-\frac12gt^2 \), where \( v_0 \) is the initial velocity of the thrown object, \( g \) is gravitational acceleration (count with the  rounded value \( 10\frac{\mathrm{m}}{\mathrm{s}^2}\)), and \( t \) is the time period of the object motion in seconds.}
{\( 45\,\mathrm{m} \)}{\( 135\,\mathrm{m} \)}{\( 360\,\mathrm{m} \)}{\( 40\,\mathrm{m} \)}{}{}}
\LANGpl{
\Question{
Rozważmy przedmiot podrzucony w górę z ziemi z początkową prędkością \( 30\frac{\mathrm{m}}{\mathrm{s}} \). Przedmiot porusza się do góry z malejącą pionową  prędkością do momentu zatrzymania. Wyznacz największą wysokość nad ziemią jaką osiągnie przedmiot. 
\[ \]
Uwaga: Odległość pionowa  \( y \) podrzuconego przedmiotu określona została wzorem \( y=v_0t-\frac12gt^2 \), gdzie \( v_0 \) jest początkową prędkością przedmiotu, a \( g \) jest przyspieszeniem grawitacyjnym  (zaokrąglij wartości  \( 10\frac{\mathrm{m}}{\mathrm{s}^2}\)), i \( t \) jest czasem poruszania się przedmiotu w sekundach).}
{\( 45\,\mathrm{m} \)}{\( 135\,\mathrm{m} \)}{\( 360\,\mathrm{m} \)}{\( 40\,\mathrm{m} \)}{}{}}
\LANGcs{
\Question{
Jestliže vrhneme tělesem svisle vzhůru, je jeho pohyb ve svislém směru (osa y) popsán rovnicí  \( y=v_0t-\frac12gt^2 \), kde \( v_0 \) je počáteční rychlost, kterou těleso vrhneme, \( g \) je tíhové zrychlení (počítejme se zaokrouhlenou hodnotou \( 10\frac{\mathrm{m}}{\mathrm{s}^2}\)) a  \( t \) vyjadřuje dobu vrhu v sekundách. Určete, do jaké maximální výšky těleso vystoupá, je-li vrženo počáteční rychlostí \( 30\frac{\mathrm{m}}{\mathrm{s}} \).}
{\( 45\,\mathrm{m} \)}{\( 135\,\mathrm{m} \)}{\( 360\,\mathrm{m} \)}{\( 40\,\mathrm{m} \)}{}{}}
\LANGsk{
\Question{
Ak vrhneme teleso zvisle nahor, je jeho pohyb v zvislom smere (os y) popísaný rovnicou  \( y=v_0t-\frac12gt^2 \), kde \( v_0 \) je začiatočná rýchlosť, ktorou teleso vrhneme, \( g \) je tiažové zrýchlenie (počítajme so zaokrúhlenou hodnotou \( 10\frac{\mathrm{m}}{\mathrm{s}^2}\)) a  \( t \) vyjadruje čas vrhu v sekundách. Určte, akú maximálnu výšku teleso dosiahne, ak je vrhnuté začiatočnou rýchlosťou \( 30\frac{\mathrm{m}}{\mathrm{s}} \).}
{\( 45\,\mathrm{m} \)}{\( 135\,\mathrm{m} \)}{\( 360\,\mathrm{m} \)}{\( 40\,\mathrm{m} \)}{}{}}
\LANGes{
\Question{
Considera un objeto lanzado hacia arriba desde el suelo con una velocidad inicial de \( 30\frac{\mathrm{m}}{\mathrm{s}} \). El objeto se mueve hacia arriba disminuyendo su velocidad hasta que para. Luego empieza a moverse hacia abajo. Encuentra la mayor altura que alcanza.
\[ \]
Nota: La distancia vertical \( y \) de un objeto lanzado se puede describir mediante la ecuación \( y=v_0t-\frac12gt^2 \), dónde \( v_0 \) es la velocidad inicial del objeto lanzado,  \( g \) es la aceleración de la gravedad (cuenta con el valor aproximado \( 10\frac{\mathrm{m}}{\mathrm{s}^2}\)), y \( t \) es el tiempo de movimiento del objeto en segundos.}
{\( 45\,\mathrm{m} \)}{\( 135\,\mathrm{m} \)}{\( 360\,\mathrm{m} \)}{\( 40\,\mathrm{m} \)}{}{}}
\End

\Data\ID{1003148602}
\Updated{2021-04-07 08:04:40}
\Level{C}
\SubArea{Quadratic Functions}
\LANGen{
\Question{
Consider an object thrown at an angle of \( 30^{\circ} \) above the horizontal with the initial velocity of \( 40\frac{\mathrm{m}}{\mathrm{s}} \). How long does it take for the object to reach its maximum height? 
\[ \]
Note: The height \( y \) of an object thrown is described by the formula \( y=v_0t\sin\alpha-\frac12gt^2 \), where \( v_0 \) is the initial velocity, \( g \) is gravitational acceleration (count with the rounded value \( 10\frac{\mathrm{m}}{\mathrm{s}^2}\)), \( t \) is the time period of the object motion in seconds, and \( \alpha \) is the angle to the horizontal at which the object is thrown.}
{\( 2\,\mathrm{s} \)}{\( 4\,\mathrm{s} \)}{\( 8\,\mathrm{s} \)}{\( 1\,\mathrm{s} \)}{}{}}
\LANGpl{
\Question{
Rozważ obiekt rzucany pod kątem \( 30^{\circ} \) powyżej poziomu z początkową prędkością \( 40\frac{\mathrm{m}}{\mathrm{s}} \). Ile czasu zajmuje przedmiotowi osiągnięcie maksymalnej wysokości?
Note:  wysokość \( y \)   rzutowanego obiektu jest opisana za pomocą wzoru  \( y=v_0t\sin\alpha-\frac12gt^2 \), gdzie \( v_0 \) to prędkość początkowa, \( g \)  to przyspieszenie grawitacyjne (liczba z zaokrągloną wartością \( 10\frac{\mathrm{m}}{\mathrm{s}^2}\)), \( t \) to czas okres ruchu obiektu w sekundach, a  \( \alpha \) jest kątem do poziomu, w którym obiekt jest rzucany.}
{\( 2\,\mathrm{s} \)}{\( 4\,\mathrm{s} \)}{\( 8\,\mathrm{s} \)}{\( 1\,\mathrm{s} \)}{}{}}
\LANGcs{
\Question{
Jestliže vrhneme tělesem šikmo vzhůru, je jeho pohyb ve svislém směru (osa y) popsán rovnicí \( y=v_0t\sin\alpha-\frac12gt^2 \), kde \( v_0 \) je počáteční rychlost tělesa, \( \alpha \) (tzv. elevační úhel) je úhel mezi vodorovným směrem a směrem \( v_0 \), \( g \) je tíhové zrychlení (počítejte se zaokrouhlenou hodnotou \( 10\frac{\mathrm{m}}{\mathrm{s}^2}\)) a \( t \) vyjadřuje dobu vrhu v sekundách. Určete, jak dlouho bude těleso stoupat do maximální výšky, je-li  \( \alpha=30^{\circ} \) a \( v_0=40\frac{\mathrm{m}}{\mathrm{s}} \).}
{\( 2\,\mathrm{s} \)}{\( 4\,\mathrm{s} \)}{\( 8\,\mathrm{s} \)}{\( 1\,\mathrm{s} \)}{}{}}
\LANGsk{
\Question{
Ak vrhneme teleso šikmo nahor, je jeho pohyb v zvislom smere (os y) popísaný rovnicou \( y=v_0t\sin\alpha-\frac12gt^2 \), kde \( v_0 \) je začiatočná rýchlosť telesa v smere, ktorý zviera s vodorovnou rovinou uhol \( \alpha \) (tzv. elevačný uhol), \( g \) je tiažové zrýchlenie (počítajte so zaokrúhlenou hodnotou \( 10\frac{\mathrm{m}}{\mathrm{s}^2}\)) a \( t \) vyjadruje čas vrhu v sekundách. Ako dlho bude trvať, kým teleso dosiahne maximálnu výšku, ak platí \( \alpha=30^{\circ} \) a \( v_0=40\frac{\mathrm{m}}{\mathrm{s}} \)?}
{\( 2\,\mathrm{s} \)}{\( 4\,\mathrm{s} \)}{\( 8\,\mathrm{s} \)}{\( 1\,\mathrm{s} \)}{}{}}
\LANGes{
\Question{
Considera un objeto lanzado con un ángulo de \( 30^{\circ} \) sobre la horizontal a una velocidad inicial de \( 40\frac{\mathrm{m}}{\mathrm{s}} \). ¿Cuánto tiempo tardará el objeto en alcanzar su altura máxima?
\[ \]
Nota: La altura \( y \) de un objeto lanzado se puede describir mediante la ecuación \( y=v_0t\sin\alpha-\frac12gt^2 \), dónde \( v_0 \) es la velocidad inicial,  \( g \) es la aceleración de la gravedad (considera un valor aproximado de \( 10\frac{\mathrm{m}}{\mathrm{s}^2}\)), y \( t \) es el tiempo de movimiento del objeto en segundos y \( \alpha \) es el ángulo horizontal con el cuál lanzamos el objeto.}
{\( 2\,\mathrm{s} \)}{\( 4\,\mathrm{s} \)}{\( 8\,\mathrm{s} \)}{\( 1\,\mathrm{s} \)}{}{}}
\End

\Data\ID{1003158901}
\Updated{2021-04-07 09:04:46}
\Level{C}
\SubArea{Quadratic Functions}
\LANGen{
\Question{
An object is moving with a constant deceleration in a straight line. Displacement \( s \) (in metres) in time \( t \) (in seconds) is modelled by \( s=24t-3t^2 \). Find the displacement of the object from the moment it starts to decelerate until it stops.}
{\( 48\,\mathrm{m} \)}{\( 144\,\mathrm{m} \)}{\( 16\,\mathrm{m} \)}{\( 96\,\mathrm{m} \)}{}{}}
\LANGpl{
\Question{
Przedmiot porusza się w linii prostej ze stałym zwalnianiem. Przemieszczanie się \( s \) (w metrach) w czasie  \( t \) (w sekundach) określa wzór  \( s=24t-3t^2 \).  Wyznacz wartość przemieszczenia się przedmiotu od moment kiedy zaczyna zwalniać do momentu zatrzymania.}
{\( 48\,\mathrm{m} \)}{\( 144\,\mathrm{m} \)}{\( 16\,\mathrm{m} \)}{\( 96\,\mathrm{m} \)}{}{}}
\LANGcs{
\Question{
Těleso se pohybuje pohybem rovnoměrně zpomaleným. Uražená vzdálenost (dráha \( s \)) je funkcí času a je určena předpisem \( s=24t-3t^2 \). Do jaké vzdálenosti se těleso dostane, než zastaví? Dráhu \( s \) udáváme v metrech, čas \( t \) v sekundách.}
{\( 48\,\mathrm{m} \)}{\( 144\,\mathrm{m} \)}{\( 16\,\mathrm{m} \)}{\( 96\,\mathrm{m} \)}{}{}}
\LANGsk{
\Question{
Teleso sa pohybuje rovnomerne spomaleným pohybom. Prejdená vzdialenosť  (dráha \( s \)) je funkciou času a je určená predpisom \( s=24t-3t^2 \). Do akej vzdialenosti sa teleso dostane kým zastaví? Dráhu \( s \) udávame v metroch a čas \( t \) v sekundách.}
{\( 48\,\mathrm{m} \)}{\( 144\,\mathrm{m} \)}{\( 16\,\mathrm{m} \)}{\( 96\,\mathrm{m} \)}{}{}}
\LANGes{
\Question{
Un objeto se mueve con una desaceleración constante en línea recta. Su desplazamiento \( s \) (en metros) se puede describir por \( s=24t-3t^2 \) siendo \( t \) el tiempo (en segundos) . Encuentra el desplazamiento del objeto desde el momento en el cual empieza a desacelerar hasta que para.}
{\( 48\,\mathrm{m} \)}{\( 144\,\mathrm{m} \)}{\( 16\,\mathrm{m} \)}{\( 96\,\mathrm{m} \)}{}{}}
\End

\Data\ID{1003158902}
\Updated{2021-04-07 09:04:48}
\Level{C}
\SubArea{Quadratic Functions}
\IMG{1103158902.pdf}
\LANGen{
\Question{
The length of a rectangle is \( 4\,\mathrm{cm} \) and width is \( x\,\mathrm{cm} \). The rectangle is divided by vertical crossing line into two parts so that one part is a square with the side of \( x\,\mathrm{cm} \) (see the picture). What is the maximum area of the remaining part of the rectangle?}
{\( 4\,\mathrm{cm}^2 \)}{\( 2\,\mathrm{cm}^2 \)}{\( 16\,\mathrm{cm}^2 \)}{\( 1\,\mathrm{cm}^2 \)}{}{}}
\LANGpl{
\Question{
Długość prostokąta wynosi \( 4\,\mathrm{cm} \), a szerokość \( x\,\mathrm{cm} \). Prostokąt jest podzielony pionową linią na dwie części, tak aby jedna część była kwadratem o boku \( x\,\mathrm{cm} \) (patrz obrazek). Jaki jest maksymalny obszar pozostałej części prostokąta?}
{\( 4\,\mathrm{cm}^2 \)}{\( 2\,\mathrm{cm}^2 \)}{\( 16\,\mathrm{cm}^2 \)}{\( 1\,\mathrm{cm}^2 \)}{}{}}
\LANGcs{
\Question{
Obdélník se stranami o velikostech \( 4\,\mathrm{cm} \) a \( x\,\mathrm{cm} \) rozdělíme příčkou tak, že vznikne čtverec o straně \( x\,\mathrm{cm} \) (viz obrázek). Jaký bude maximální možný obsah zbývající části obdélníka?}
{\( 4\,\mathrm{cm}^2 \)}{\( 2\,\mathrm{cm}^2 \)}{\( 16\,\mathrm{cm}^2 \)}{\( 1\,\mathrm{cm}^2 \)}{}{}}
\LANGsk{
\Question{
Obdĺžnik s dĺžkou \( 4\,\mathrm{cm} \) a šírkou \( x\,\mathrm{cm} \) rozdelíme priečkou na dve časti tak, aby vznikol štvorec so stranou \( x\,\mathrm{cm} \) (pozri obrázok). Aký bude maximálny možný obsah zostávajúcej časti?}
{\( 4\,\mathrm{cm}^2 \)}{\( 2\,\mathrm{cm}^2 \)}{\( 16\,\mathrm{cm}^2 \)}{\( 1\,\mathrm{cm}^2 \)}{}{}}
\LANGes{
\Question{
La base de un rectángulo es \( 4\,\mathrm{cm} \) y su altura es \( x\,\mathrm{cm} \). El rectángulo está dividido en dos partes por un segmento vertical para que una parte sea un cuadrado de lado \( x\,\mathrm{cm} \) (observa en el dibujo). ¿Cuál es el área máxima de la parte que queda del rectángulo?}
{\( 4\,\mathrm{cm}^2 \)}{\( 2\,\mathrm{cm}^2 \)}{\( 16\,\mathrm{cm}^2 \)}{\( 1\,\mathrm{cm}^2 \)}{}{}}
\End

\Data\ID{1003162301}
\Updated{2022-04-23 05:04:54}
\Level{C}
\SubArea{Quadratic Functions}
\LANGen{
\Question{
Find all the values of the real parameter \( a \) such that \( f(x)=ax^2-2 \) is decreasing on \( (0;\infty) \).}
{\( a\in(-\infty;0) \)}{\( a\in(0;\infty) \)}{\( a\in[2;+\infty) \)}{\( a\in(-\infty;2] \)}{}{}}
\LANGpl{
\Question{
Wyznacz wszystkie wartości rzeczywistego parametru \( a \) takie, że \( f(x)=ax^2-2 \) jest malejąca w  \( (0;\infty) \).}
{\( a\in(-\infty;0) \)}{\( a\in(0;\infty) \)}{\( a\in\langle2;+\infty) \)}{\( a\in(-\infty;2\rangle \)}{}{}}
\LANGcs{
\Question{
Pro které hodnoty parametru \( a \) bude funkce \( f(x)=ax^2-2 \) klesající na intervalu \( (0;\infty) \)?}
{\( a\in(-\infty;0) \)}{\( a\in(0;\infty) \)}{\( a\in\langle2;+\infty) \)}{\( a\in(-\infty;2\rangle \)}{}{}}
\LANGsk{
\Question{
Pre ktoré hodnoty reálneho parametra \( a \), bude funkcia \( f(x)=ax^2-2 \) klesajúca na intervale \( (0;\infty) \)?}
{\( a\in(-\infty;0) \)}{\( a\in(0;\infty) \)}{\( a\in\langle2;+\infty) \)}{\( a\in(-\infty;2\rangle \)}{}{}}
\LANGes{
\Question{
Encuentra todos los valores del parametro real \( a \) para los que la función \( f(x)=ax^2-2 \) sea decreciente en \( (0;\infty) \).}
{\( a\in(-\infty;0) \)}{\( a\in(0;\infty) \)}{\( a\in[2;+\infty) \)}{\( a\in(-\infty;2] \)}{}{}}
\End

\Data\ID{1003162302}
\Updated{2022-04-23 05:04:54}
\Level{C}
\SubArea{Quadratic Functions}
\LANGen{
\Question{
Find all the values of the real parameter \( m \) such that \( f(x)=-2(x-m)^2+3 \) is an even function.}
{\( m=0 \)}{\( m=3 \)}{\( m=-3 \)}{\( m\in(-\infty;\infty) \)}{}{}}
\LANGpl{
\Question{
Wyznacz wszystkie wartości rzeczywistego parametru \( m \) takie, dla których \( f(x)=-2(x-m)^2+3 \) jest funkcją parzystą.}
{\( m=0 \)}{\( m=3 \)}{\( m=-3 \)}{\( m\in(-\infty;\infty) \)}{}{}}
\LANGcs{
\Question{
Najděte všechny hodnoty reálného  parametru \( m \), pro které je funkce \( f(x)=-2(x-m)^2+3 \) sudá.}
{\( m=0 \)}{\( m=3 \)}{\( m=-3 \)}{\( m\in(-\infty;\infty) \)}{}{}}
\LANGsk{
\Question{
Pre ktoré hodnoty reálneho parametra \( m \), bude funkcia \( f(x)=-2(x-m)^2+3 \) párna?}
{\( m=0 \)}{\( m=3 \)}{\( m=-3 \)}{\( m\in(-\infty;\infty) \)}{}{}}
\LANGes{
\Question{
¿ Para qué valores del parámetro real \( m \)  \( f(x)=-2(x-m)^2+3 \) es una función par?}
{\( m=0 \)}{\( m=3 \)}{\( m=-3 \)}{\( m\in(-\infty;\infty) \)}{}{}}
\End

\Data\ID{1003162303}
\Updated{2022-04-23 05:04:54}
\Level{C}
\SubArea{Quadratic Functions}
\LANGen{
\Question{
Find all the values of the real parameter \( m \) such that \( f(x)=3(x+m)^2-2 \) is increasing on \( (0;\infty) \).}
{\( m\in[0;\infty) \)}{\( m\in(-\infty;0) \)}{\( m\in(-\infty;0] \)}{\( m\in(-\infty;2] \)}{}{}}
\LANGpl{
\Question{
Wyznacz wszystkie wartości rzeczywistego parametru \( m \) takie, dla których \( f(x)=3(x+m)^2-2 \) jest rosnąca w \( (0;\infty) \).}
{\( m\in\langle0;\infty) \)}{\( m\in(-\infty;0) \)}{\( m\in(-\infty;0\rangle \)}{\( m\in(-\infty;2\rangle \)}{}{}}
\LANGcs{
\Question{
Najděte všechny hodnoty reálného parametru \( m \), pro které je funkce \( f(x)=3(x+m)^2-2 \) rostoucí na intervalu \( (0;\infty) \).}
{\( m\in\langle0;\infty) \)}{\( m\in(-\infty;0) \)}{\( m\in(-\infty;0\rangle \)}{\( m\in(-\infty;2\rangle \)}{}{}}
\LANGsk{
\Question{
Pre ktoré hodnoty reálneho parametra \( m \), bude funkcia \( f(x)=3(x+m)^2-2 \) rastúca na intervale \( (0;\infty) \)?}
{\( m\in\langle0;\infty) \)}{\( m\in(-\infty;0) \)}{\( m\in(-\infty;0\rangle \)}{\( m\in(-\infty;2\rangle \)}{}{}}
\LANGes{
\Question{
¿Para qué valores del parámetro real  \( m \) \( f(x)=3(x+m)^2-2 \) es una función creciente en \( (0;\infty) \)?}
{\( m\in[0;\infty) \)}{\( m\in(-\infty;0) \)}{\( m\in(-\infty;0] \)}{\( m\in(-\infty;2] \)}{}{}}
\End

\Data\ID{1003162304}
\Updated{2021-02-15 07:02:59}
\Level{C}
\SubArea{Quadratic Functions}
\LANGen{
\Question{
Find all the values of the real parameter \( m \) such that \( f(x)=-x^2+2xm-m^2+2 \) is increasing on \( (-\infty;0) \).}
{\( m\in[0;\infty) \)}{\( m\in(-\infty;0) \)}{\( m\in(-\infty;0] \)}{\( m\in(-\infty;2] \)}{}{}}
\LANGpl{
\Question{
Wyznacz wszystkie wartości parametru rzeczywistego  \( m \) takie, dla których \( f(x)=-x^2+2xm-m^2+2 \) jest rosnąca w \( (-\infty;0) \).}
{\( m\in\langle0;\infty) \)}{\( m\in(-\infty;0) \)}{\( m\in(-\infty;0\rangle \)}{\( m\in(-\infty;2\rangle \)}{}{}}
\LANGcs{
\Question{
Pro které hodnoty parametru \( m \) bude funkce \( f(x)=-x^2+2xm-m^2+2 \) rostoucí na intervalu \( (-\infty;0) \)?}
{\( m\in\langle0;\infty) \)}{\( m\in(-\infty;0) \)}{\( m\in(-\infty;0\rangle \)}{\( m\in(-\infty;2\rangle \)}{}{}}
\LANGsk{
\Question{
Pre ktoré hodnoty reálneho parametra \( m \), bude funkcia \( f(x)=-x^2+2xm-m^2+2 \)  rastúca na intervale \( (-\infty;0) \)?}
{\( m\in\langle0;\infty) \)}{\( m\in(-\infty;0) \)}{\( m\in(-\infty;0\rangle \)}{\( m\in(-\infty;2\rangle \)}{}{}}
\LANGes{
\Question{
Encuentra todos los valores del parámetro real \( m \) para los cuales la función\( f(x)=-x^2+2xm-m^2+2 \) es creciente en \( (-\infty;0) \).}
{\( m\in[0;\infty) \)}{\( m\in(-\infty;0) \)}{\( m\in(-\infty;0] \)}{\( m\in(-\infty;2] \)}{}{}}
\End

\Data\ID{1003162305}
\Updated{2022-04-23 05:04:54}
\Level{C}
\SubArea{Quadratic Functions}
\LANGen{
\Question{
Find all the values of the real parameter \( p \) such that \( f(x)=3(x-2)^2+p \) is nonnegative for all \( x\in\mathbb{R} \).}
{\( p\in[0;\infty) \)}{\( p\in(-\infty;0) \)}{\( p=0 \)}{\( p\in(0;\infty) \)}{}{}}
\LANGpl{
\Question{
Wyznacz wszystkie wartości parametru rzeczywistego \( p \), takie dla których \( f(x)=3(x-2)^2+p \) jest nieujemna dla wszystkich \( x\in\mathbb{R} \).}
{\( p\in\langle0;\infty) \)}{\( p\in(-\infty;0) \)}{\( p=0 \)}{\( p\in(0;\infty) \)}{}{}}
\LANGcs{
\Question{
Určete obor parametru \( p \) tak, aby funkce \( f(x)=3(x-2)^2+p \) nabývala na celém svém definičním oboru nezáporných hodnot.}
{\( p\in\langle0;\infty) \)}{\( p\in(-\infty;0) \)}{\( p=0 \)}{\( p\in(0;\infty) \)}{}{}}
\LANGsk{
\Question{
Nájdite všetky hodnoty reálneho parametra  \( p \), pre ktoré bude funkcia \( f(x)=3(x-2)^2+p \) nezáporná pre všetky \( x\in\mathbb{R} \).}
{\( p\in\langle0;\infty) \)}{\( p\in(-\infty;0) \)}{\( p=0 \)}{\( p\in(0;\infty) \)}{}{}}
\LANGes{
\Question{
Encuentra todos los valores del parámetro real \( p \) para que \( f(x)=3(x-2)^2+p \) sea no negativa para todo \( x\in\mathbb{R} \).}
{\( p\in[0;\infty) \)}{\( p\in(-\infty;0) \)}{\( p=0 \)}{\( p\in(0;\infty) \)}{}{}}
\End

\Data\ID{1003162306}
\Updated{2020-09-30 04:09:04}
\Level{C}
\SubArea{Quadratic Functions}
\LANGen{
\Question{
Find all the values of the real parameter \( p \) such that \( f(x)=2x^2+px+p \) is positive for all \( x\in\mathbb{R} \).}
{\( p\in(0;8) \)}{\( p\in(-\infty;0)\cup(8;+\infty) \)}{\( p\in(-\infty;0) \)}{\( p\in(0;\infty) \)}{}{}}
\LANGpl{
\Question{
Wyznacz wszystkie wartości rzeczywistego parametru \( p \), takie, dla których \( f(x)=2x^2+px+p \) jest dodatnia dla wszystkich \( x\in\mathbb{R} \).}
{\( p\in(0;8) \)}{\( p\in(-\infty;0)\cup(8;+\infty) \)}{\( p\in(-\infty;0) \)}{\( p\in(0;\infty) \)}{}{}}
\LANGcs{
\Question{
Určete obor parametru \( p \) tak, aby funkce \( f(x)=2x^2+px+p \) nabývala na celém svém definičním oboru kladných hodnot.}
{\( p\in(0;8) \)}{\( p\in(-\infty;0)\cup(8;+\infty) \)}{\( p\in(-\infty;0) \)}{\( p\in(0;\infty) \)}{}{}}
\LANGsk{
\Question{
Nájdite všetky hodnoty reálneho parametra \( p \), pre ktoré bude funkcia \( f(x)=2x^2+px+p \) kladná pre všetky \( x\in\mathbb{R} \).}
{\( p\in(0;8) \)}{\( p\in(-\infty;0)\cup(8;+\infty) \)}{\( p\in(-\infty;0) \)}{\( p\in(0;\infty) \)}{}{}}
\LANGes{
\Question{
Encuentra todos los valores del parámetro real \( p \) para que \( f(x)=2x^2+px+p \) sea positiva para todo \( x\in\mathbb{R} \).}
{\( p\in(0;8) \)}{\( p\in(-\infty;0)\cup(8;+\infty) \)}{\( p\in(-\infty;0) \)}{\( p\in(0;\infty) \)}{}{}}
\End

\Data\ID{1003162307}
\Updated{2020-09-30 04:09:17}
\Level{C}
\SubArea{Quadratic Functions}
\LANGen{
\Question{
Find all the values of the real parameter \( p \) such that \( f(x)=2x^2+3px+2 \) has a minimum.}
{\( p\in(-\infty;\infty) \)}{\( p\in(-\infty;0)\cup(0;+\infty) \)}{\( p=0 \)}{\( p\in[0;\infty) \)}{}{}}
\LANGpl{
\Question{
Wyznacz wszystkie wartości parametru rzeczywistego \( p \) takie, dla których \( f(x)=2x^2+3px+2 \) ma minimum.}
{\( p\in(-\infty;\infty) \)}{\( p\in(-\infty;0)\cup(0;+\infty) \)}{\( p=0 \)}{\( p\in\langle0;\infty) \)}{}{}}
\LANGcs{
\Question{
Určete obor parametru  \( p \) tak, aby funkce \( f(x)=2x^2+3px+2 \) měla minimum.}
{\( p\in(-\infty;\infty) \)}{\( p\in(-\infty;0)\cup(0;+\infty) \)}{\( p=0 \)}{\( p\in\langle0;\infty) \)}{}{}}
\LANGsk{
\Question{
Nájdite všetky hodnoty reálneho parametra \( p \) tak, aby funkcia \( f(x)=2x^2+3px+2 \) mala minimum.}
{\( p\in(-\infty;\infty) \)}{\( p\in(-\infty;0)\cup(0;+\infty) \)}{\( p=0 \)}{\( p\in\langle0;\infty) \)}{}{}}
\LANGes{
\Question{
Encuentra todos los valores del parámetro real \( p \) para que \( f(x)=2x^2+3px+2 \) tenga mínimo.}
{\( p\in(-\infty;\infty) \)}{\( p\in(-\infty;0)\cup(0;+\infty) \)}{\( p=0 \)}{\( p\in[0;\infty) \)}{}{}}
\End

\Data\ID{1003162308}
\Updated{2020-09-30 04:09:57}
\Level{C}
\SubArea{Quadratic Functions}
\LANGen{
\Question{
Find all the values of the real parameter \( p \) such that \( f(x)=(p-2) x^2+px+2 \) has a maximum.}
{\( p\in(-\infty;2) \)}{\( p\in(-\infty;-2) \)}{\( p\in(2;+\infty) \)}{\( p\in(-\infty;0) \)}{}{}}
\LANGpl{
\Question{
Wyznacz wszystkie wartości parametru rzeczywistego \( p \) takie, dla których  \( f(x)=(p-2) x^2+px+2 \) ma maksimum.}
{\( p\in(-\infty;2) \)}{\( p\in(-\infty;-2) \)}{\( p\in(2;+\infty) \)}{\( p\in(-\infty;0) \)}{}{}}
\LANGcs{
\Question{
Určete obor parametru \( p \) tak, aby funkce \( f(x)=(p-2) x^2+px+2 \) měla maximum.}
{\( p\in(-\infty;2) \)}{\( p\in(-\infty;-2) \)}{\( p\in(2;+\infty) \)}{\( p\in(-\infty;0) \)}{}{}}
\LANGsk{
\Question{
Nájdite všetky hodnoty reálneho parametra \( p \) tak, aby funkcia \( f(x)=(p-2) x^2+px+2 \) mala maximum.}
{\( p\in(-\infty;2) \)}{\( p\in(-\infty;-2) \)}{\( p\in(2;+\infty) \)}{\( p\in(-\infty;0) \)}{}{}}
\LANGes{
\Question{
Encuentra todos los valores del parámetro real \( p \) para que \( f(x)=(p-2) x^2+px+2 \) tenga máximo.}
{\( p\in(-\infty;2) \)}{\( p\in(-\infty;-2) \)}{\( p\in(2;+\infty) \)}{\( p\in(-\infty;0) \)}{}{}}
\End

\Data\ID{1003162309}
\Updated{2020-09-30 04:09:52}
\Level{C}
\SubArea{Quadratic Functions}
\LANGen{
\Question{
Find all the values of the real parameter \( p \) such that \( f(x)=px^2-4px+4p-3 \) is a negative quadratic function for all \( x\in\mathbb{R} \).}
{\( p\in(-\infty;0) \)}{\( p=0 \)}{\( p\in(0;\infty) \)}{\( p\in(-2;2) \)}{}{}}
\LANGpl{
\Question{
Wyznacz wszystkie wartości parametru rzeczywistego \( p \), takiego, dla którego \( f(x)=px^2-4px+4p-3 \) ujemną funkcją kwadratową dla wszystkich \( x\in\mathbb{R} \).}
{\( p\in(-\infty;0) \)}{\( p=0 \)}{\( p\in(0;\infty) \)}{\( p\in(-2;2) \)}{}{}}
\LANGcs{
\Question{
Určete obor parametru \( p \) tak, aby funkce \( f(x)=px^2-4px+4p-3 \) nabývala na celém svém definičním oboru záporných hodnot.}
{\( p\in(-\infty;0) \)}{\( p=0 \)}{\( p\in(0;\infty) \)}{\( p\in(-2;2) \)}{}{}}
\LANGsk{
\Question{
Nájdite všetky hodnoty reálneho parametra \( p \), pre ktoré bude funkcia \( f(x)=px^2-4px+4p-3 \) záporná pre všetky \( x\in\mathbb{R} \).}
{\( p\in(-\infty;0) \)}{\( p=0 \)}{\( p\in(0;\infty) \)}{\( p\in(-2;2) \)}{}{}}
\LANGes{
\Question{
Encuentra todos los valores del parámetro real \( p \) para que \( f(x)=px^2-4px+4p-3 \) sea una función cuadrática negativa para todo  \( x\in\mathbb{R} \).}
{\( p\in(-\infty;0) \)}{\( p=0 \)}{\( p\in(0;\infty) \)}{\( p\in(-2;2) \)}{}{}}
\End

\Data\ID{1003206002}
\Updated{2021-02-15 07:02:34}
\Level{C}
\SubArea{Quadratic Functions}
\LANGen{
\Question{
We are given three quadratic functions:
\[ \begin{aligned}
f_1(x)&=ax^2+2ax+a-3, \\
f_2(x)&=a(x-1)^2+2, \\
f_3(x)&=ax^2,
\end{aligned} \]
where \( a\in(-\infty;0) \). If possible, determine which of the given functions has the highest output value for \( x = 0.5 \).}
{\( f_2 \)}{\( f_3 \)}{\( f_1 \)}{Given information is insufficient to decide.}{}{}}
\LANGpl{
\Question{
Dane są trzy funkcje kwadratowe:
\[ \begin{aligned}
f_1(x)&=ax^2+2ax+a-3, \\
f_2(x)&=a(x-1)^2+2, \\
f_3(x)&=ax^2,
\end{aligned} \]
gdzie \( a\in(-\infty;0) \). Jeśli to możliwe określ, która z podanych funkcji ma najwyższą wartość wyjściową dla \( x = 0{,}5 \).}
{\( f_2 \)}{\( f_3 \)}{\( f_1 \)}{Podana informacja nie jest wystarczająca.}{}{}}
\LANGcs{
\Question{
Jsou dány tři kvadratické funkce:
\[ \begin{aligned}
f_1(x)&=ax^2+2ax+a-3, \\
f_2(x)&=a(x-1)^2+2, \\
f_3(x)&=ax^2,
\end{aligned} \]
kde \( a\in(-\infty;0) \). Jestliže to je možné, rozhodněte, která z uvedených funkcí má pro \( x = 0{,}5 \) největší hodnotu.}
{\( f_2 \)}{\( f_3 \)}{\( f_1 \)}{Z daných informací to není možné jednoznačně určit.}{}{}}
\LANGsk{
\Question{
Dané sú tri kvadratické funkcie:
\[ \begin{aligned}
f_1(x)&=ax^2+2ax+a-3, \\
f_2(x)&=a(x-1)^2+2, \\
f_3(x)&=ax^2,
\end{aligned} \]
kde \( a\in(-\infty;0) \). Ak je to možné, určte ktorá funkcia nadobúda pre \( x = 0{,}5 \) najväčšiu hodnotu.}
{\( f_2 \)}{\( f_3 \)}{\( f_1 \)}{Z daných informácií to nie je možné určiť.}{}{}}
\LANGes{
\Question{
Tenemos tres funciones cuadráticas:
\[ \begin{aligned}
f_1(x)&=ax^2+2ax+a-3, \\
f_2(x)&=a(x-1)^2+2, \\
f_3(x)&=ax^2,
\end{aligned} \]
dónde \( a\in(-\infty;0) \). Si es posible, determina cuál de las funciones dadas tiene el mayor valor funcional para \( x = 0.5 \).}
{\( f_2 \)}{\( f_3 \)}{\( f_1 \)}{La información dada no es suficente.}{}{}}
\End

\Data\ID{1103067801}
\Updated{2021-04-07 06:04:17}
\Level{C}
\SubArea{Quadratic Functions}
\IMG{1103067801_0.pdf}
\LANGen{
\Question{
Given the graph of the function \( f(x)=x^2 + 2x - 4 \), find the solution set of the following equation. 
\[ \left|x^2 + 2x - 4\right|= 4 \]}
{\( \{ -4; -2; 0 ; 2 \} \)}{\( \{ -4; 2 \} \)}{\( \{ -4; 0; 2 \} \)}{\( \left\{ -1-\sqrt5;-1+ \sqrt5 \right\} \)}{}{}}
\LANGpl{
\Question{
Dany jest wykres funkcji \( f(x)=x^2 + 2x - 4 \), wyznacz zbiór rozwiązań następującego równania. 
\[ \left|x^2 + 2x - 4\right|= 4 \]}
{\( \{ -4; -2; 0 ; 2 \} \)}{\( \{ -4; 2 \} \)}{\( \{ -4; 0; 2 \} \)}{\( \left\{ -1-\sqrt5;-1+ \sqrt5 \right\} \)}{}{}}
\LANGcs{
\Question{
Pomocí grafu funkce \( f(x)=x^2 + 2x - 4 \) určete množinu řešení dané rovnice.
\[ \left|x^2 + 2x - 4\right|= 4 \]}
{\( \{ -4; -2; 0 ; 2 \} \)}{\( \{ -4; 2 \} \)}{\( \{ -4; 0; 2 \} \)}{\( \left\{ -1-\sqrt5;-1+ \sqrt5 \right\} \)}{}{}}
\LANGsk{
\Question{
Pomocou grafu funkcie \( f(x)=x^2 + 2x - 4 \), určte množinu riešení danej rovnice. 
\[ \left|x^2 + 2x - 4\right|= 4 \]}
{\( \{ -4; -2; 0 ; 2 \} \)}{\( \{ -4; 2 \} \)}{\( \{ -4; 0; 2 \} \)}{\( \left\{ -1-\sqrt5;-1+ \sqrt5 \right\} \)}{}{}}
\LANGes{
\Question{
Dada la gráfica de la función \( f(x)=x^2 + 2x - 4 \), halla el cunjunto solución de la siguiente ecuación.
\[ \left|x^2 + 2x - 4\right|= 4 \]}
{\( \{ -4; -2; 0 ; 2 \} \)}{\( \{ -4; 2 \} \)}{\( \{ -4; 0; 2 \} \)}{\( \left\{ -1-\sqrt5;-1+ \sqrt5 \right\} \)}{}{}}
\End

\Data\ID{1103067809}
\Updated{2021-04-07 06:04:34}
\Level{C}
\SubArea{Quadratic Functions}
\IMG{1103067809.pdf}
\LANGen{
\Question{
Given the graphs of the functions  \( f(x)=\frac12x^2-3 \) and \( g(x)=\frac12x \), find the solution set of the following equation.
\[ \left|\frac12 x^2-3\right|=\left|\frac12 x\right| \]}
{\( \{ -3; -2; 2; 3 \} \)}{\( \{ -2; 3 \} \)}{\( \{ 2; 3 \} \)}{\( \left\{ -\sqrt6; -2; \sqrt6; 3 \right\} \)}{}{}}
\LANGpl{
\Question{
Dany jest wykres funkcji  \( f(x)=\frac12x^2-3 \) i \( g(x)=\frac12x \), wyznacz zbiór rozwiązań następującego równania. 
\[ \left|\frac12 x^2-3\right|=\left|\frac12 x\right| \]}
{\( \{ -3; -2; 2; 3 \} \)}{\( \{ -2; 3 \} \)}{\( \{ 2; 3 \} \)}{\( \left\{ -\sqrt6; -2; \sqrt6; 3 \right\} \)}{}{}}
\LANGcs{
\Question{
Pomocí grafů funkcí  \( f(x)=\frac12x^2-3 \) a \( g(x)=\frac12x \) určete množinu řešení dané rovnice.
\[ \left|\frac12 x^2-3\right|=\left|\frac12 x\right| \]}
{\( \{ -3; -2; 2; 3 \} \)}{\( \{ -2; 3 \} \)}{\( \{ 2; 3 \} \)}{\( \left\{ -\sqrt6; -2; \sqrt6; 3 \right\} \)}{}{}}
\LANGsk{
\Question{
Pomocou grafov funkcií  \( f(x)=\frac12x^2-3 \) a \( g(x)=\frac12x \),  určte množinu riešení danej rovnice.
\[ \left|\frac12 x^2-3\right|=\left|\frac12 x\right| \]}
{\( \{ -3; -2; 2; 3 \} \)}{\( \{ -2; 3 \} \)}{\( \{ 2; 3 \} \)}{\( \left\{ -\sqrt6; -2; \sqrt6; 3 \right\} \)}{}{}}
\LANGes{
\Question{
Dadas las gráficas de las funciones \( f(x)=\frac12x^2-3 \) y \( g(x)=\frac12x \), encuentra el conjunto solución de la siguiente ecuación,
\[ \left|\frac12 x^2-3\right|=\left|\frac12 x\right| \]}
{\( \{ -3; -2; 2; 3 \} \)}{\( \{ -2; 3 \} \)}{\( \{ 2; 3 \} \)}{\( \left\{ -\sqrt6; -2; \sqrt6; 3 \right\} \)}{}{}}
\End

\Data\ID{1103082702}
\Updated{2022-04-23 05:04:48}
\Level{C}
\SubArea{Quadratic Functions}
\IMG{1103082702.pdf}
\LANGen{
\Question{
Function \( f \) is given completely by the graph. Identify which of the following statements is true.}
{\( f(x)=\left|x^2-1\right|;\ x\in[-2;2] \)}{\( f(x)=\left|x^2\right|-1;\ x\in[-2;2] \)}{\( f(x)=-\left|x^2+1\right|;\ x\in[-2;2] \)}{\( f(x)=\left|-x^2\right|+1;\ x\in[-2;2] \)}{}{}}
\LANGpl{
\Question{
Funkcję \( f \) przedstawiono na wykresie poniżej. Które z poniższych stwierdzeń jest prawdziwe?}
{\( f(x)=\left|x^2-1\right|;\ x\in\langle-2;2\rangle \)}{\( f(x)=\left|x^2\right|-1;\ x\in\langle-2;2\rangle \)}{\( f(x)=-\left|x^2+1\right|;\ x\in\langle-2;2\rangle \)}{\( f(x)=\left|-x^2\right|+1;\ x\in\langle-2;2\rangle \)}{}{}}
\LANGcs{
\Question{
Funkce \( f \) je dána grafem. Určete, které z následujících tvrzení je pravdivé.}
{\( f(x)=\left|x^2-1\right|;\ x\in\langle-2;2\rangle \)}{\( f(x)=\left|x^2\right|-1;\ x\in\langle-2;2\rangle \)}{\( f(x)=-\left|x^2+1\right|;\ x\in\langle-2;2\rangle \)}{\( f(x)=\left|-x^2\right|+1;\ x\in\langle-2;2\rangle \)}{}{}}
\LANGsk{
\Question{
Funkcia \( f \) je daná grafom. Vyberte pravdivé tvrdenie o funkcii \( f \).}
{\( f(x)=\left|x^2-1\right|;\ x\in\langle-2;2\rangle \)}{\( f(x)=\left|x^2\right|-1;\ x\in\langle-2;2\rangle \)}{\( f(x)=-\left|x^2+1\right|;\ x\in\langle-2;2\rangle \)}{\( f(x)=\left|-x^2\right|+1;\ x\in\langle-2;2\rangle \)}{}{}}
\LANGes{
\Question{
La función \( f \) está definida por la gráfica. Identifica cuál de las declaraciones es correcta.}
{\( f(x)=\left|x^2-1\right|;\ x\in[-2;2] \)}{\( f(x)=\left|x^2\right|-1;\ x\in[-2;2] \)}{\( f(x)=-\left|x^2+1\right|;\ x\in[-2;2] \)}{\( f(x)=\left|-x^2\right|+1;\ x\in[-2;2] \)}{}{}}
\End

\Data\ID{1103120009}
\Updated{2021-04-07 07:04:34}
\Level{C}
\SubArea{Quadratic Functions}
\IMG{1103120009.pdf}
\LANGen{
\Question{
In the picture there are two parabolas. One parabola can be mapped onto the other by shifting. These parabolas are graphs of the quadratic functions 
\[ f(x)=-(x-a)^2+b\ \text{ and }\ g(x)=-(x-c)^2+d, \]
where \( a \), \( b \), \( c \), \( d\in\mathbb{R} \).
Following statements describe the relations  between the pairs of the coefficients \( a \), \( b \), \( c \) and \( d \). Choose the true statement.}
{\( a=c-1\wedge b=d+4 \)}{\( a=c+1\wedge b=d-4 \)}{\( a=c-4\wedge b=d+1 \)}{\( a=c+4\wedge b=d-1 \)}{}{}}
\LANGpl{
\Question{
Na rysunku przedstawiono dwie parabole. Jedna parabola może być odwzorowana na drugiej poprzez przesunięcie.  Podane parabole są wykresami funkcji kwadratowych  
\[ f(x)=-(x-a)^2+b\ \text{ i }\ g(x)=-(x-c)^2+d, \]
gdzie \( a \), \( b \), \( c \), \( d\in\mathbb{R} \).
Poniższe stwierdzenia opisują relacje pomiędzy parami współczynników \( a \), \( b \), \( c \) i \( d \). Wybierz prawdziwe stwierdzenie.}
{\( a=c-1\wedge b=d+4 \)}{\( a=c+1\wedge b=d-4 \)}{\( a=c-4\wedge b=d+1 \)}{\( a=c+4\wedge b=d-1 \)}{}{}}
\LANGcs{
\Question{
Na obrázku jsou dvě paraboly (posunutím je možné zobrazit jednu na druhou). Tyto paraboly představují grafy kvadratických funkcí 
\[ f(x)=-(x-a)^2+b,\qquad g(x)=-(x-c)^2+d, \]
kde \( a \), \( b \), \( c \), \( d\in\mathbb{R} \).
Vyberte možnost, která správně vyjadřuje vztah mezi dvojicemi koeficientů \( a \), \( c \) a   \( b \), \( d \).}
{\( a=c-1\wedge b=d+4 \)}{\( a=c+1\wedge b=d-4 \)}{\( a=c-4\wedge b=d+1 \)}{\( a=c+4\wedge b=d-1 \)}{}{}}
\LANGsk{
\Question{
Na obrázku sú dve paraboly. Jedna parabola môže byť posunutím zobrazená do druhej.  Tieto paraboly sú grafmi kvadratických funkcií \[ f(x)=-(x-a)^2+b\ \text{ a }\ g(x)=-(x-c)^2+d, \]
kde \( a \), \( b \), \( c \), \( d\in\mathbb{R} \).
Ktoré z nasledujúcich tvrdení vyjadruje vzťah medzi dvojicami koeficientov \( a \), \( b \), \( c \) a \( d \)?}
{\( a=c-1\wedge b=d+4 \)}{\( a=c+1\wedge b=d-4 \)}{\( a=c-4\wedge b=d+1 \)}{\( a=c+4\wedge b=d-1 \)}{}{}}
\LANGes{
\Question{
En el dibujo hay dos parábolas. Una parabola puede identificarse con la otra mediante traslación. Estas parábolas son gráficas de las funciones cuadráticas
\[ f(x)=-(x-a)^2+b\ \text{ y }\ g(x)=-(x-c)^2+d, \]
dónde \( a \), \( b \), \( c \), \( d\in\mathbb{R} \).
Las declaraciones siguientes describen las relaciones entre las parejas de coeficientes \( a \), \( b \), \( c \) y \( d \). Elige la declaración correcta.}
{\( a=c-1\wedge b=d+4 \)}{\( a=c+1\wedge b=d-4 \)}{\( a=c-4\wedge b=d+1 \)}{\( a=c+4\wedge b=d-1 \)}{}{}}
\End

\Data\ID{1103120010}
\Updated{2021-04-07 07:04:19}
\Level{C}
\SubArea{Quadratic Functions}
\LANGen{
\Question{
Let \( f(x)=(x+a)^2+b \) and \( g(x)=(x+a-2)^2+b+3 \), where \( a \), \( b\in\mathbb{R} \). In which of the next four given pictures are the graphs of both functions \( f \) and \( g \)?}
{}{}{}{}{}{}}
\LANGpl{
\Question{
Niech \( f(x)=(x+a)^2+b \) i \( g(x)=(x+a-2)^2+b+3 \), gdzie \( a \), \( b\in\mathbb{R} \). Na którym z poniższych rysunków przedstawiono wykresy obydwu funkcji \( f \) i \( g \)?}
{}{}{}{}{}{}}
\LANGcs{
\Question{
Mějme funkce \( f(x)=(x+a)^2+b \), \( g(x)=(x+a-2)^2+b+3 \), kde \( a \), \( b\in\mathbb{R} \). Vyberte obrázek, na kterém jsou grafy funkcí  \( f \) a \( g \).}
{}{}{}{}{}{}}
\LANGsk{
\Question{
Dané sú funkcie \( f(x)=(x+a)^2+b \) a \( g(x)=(x+a-2)^2+b+3 \), kde \( a \), \( b\in\mathbb{R} \). Na ktorom s nasledujúcich štyroch obrázkov sa nachádzajú grafy obidvoch funkcií \( f \) a \( g \)?}
{}{}{}{}{}{}}
\LANGes{
\Question{
Sean \( f(x)=(x+a)^2+b \) y \( g(x)=(x+a-2)^2+b+3 \), dónde \( a \), \( b\in\mathbb{R} \). ¿En cuál de los cuatro dibujos aparecen gráficas de ambas funciones \( f \) y \( g \)?}
{}{}{}{}{}{}}

\IMGanswer{1103120010a_0.pdf}
\IMGanswer{1103120010b_0.pdf}
\IMGanswer{1103120010c_0.pdf}
\IMGanswer{1103120010d_0.pdf}\End

\Data\ID{1103148603}
\Updated{2022-01-22 07:01:58}
\Level{C}
\SubArea{Quadratic Functions}
\IMG{1103148603_0.pdf}
\LANGen{
\Question{
Consider a simple circuit in which a battery of electromotive force \( U_e \) and internal resistance \( R_i \) drives a current \( I \) through an external resistor of resistance \( R \) (see figure). The external resistor could be for example an electric light, an electric heating element, or, maybe, an electric motor. The basic purpose of the circuit is to transfer energy from the battery to the external resistor, where it actually does something useful for us (e.g. lighting a light bulb, or lifting a weight).
\[ \]
The power \( P \) transferred to the external resistor is described by the formula \( P=U_eI-R_i I^2 \). What maximum power can be transferred to the external resistor if we have the source with \( R_i=0.25\,\Omega \) and \( U_e=20\,\mathrm{V} \)?}
{\( 400\,\mathrm{W} \)}{\( 80\,\mathrm{W} \)}{\( 40\,\mathrm{W} \)}{\( 790\,\mathrm{W} \)}{}{}}
\LANGpl{
\Question{
Rozważmy prosty obwód, w którym bateria siły elektromotorycznej  \( U_e \)  i opór wewnętrzny \( R_i \)  prowadzą prąd  \( I \) przez zewnętrzny rezystor o oporze \( R \)  (patrz rysunek). Zewnętrznym rezystorem może być na przykład światło elektryczne, elektryczny element grzejny lub, być może, silnik elektryczny. Podstawowym celem obwodu jest przesyłanie energii z akumulatora do zewnętrznego rezystora, gdzie faktycznie robi on coś użytecznego dla nas (np. zapalanie żarówki lub podnoszenie ciężaru).
\[ \]

Moc \( P \) przekazywana do zewnętrznego rezystora opisana jest wzorem \( P=U_eI-R_i I^2 \). Jaką maksymalną moc można przenieść na zewnętrzny rezystor, jeśli mamy źródło z $R_i=0{,}25\,\Omega$ i $U_e=20\,V$?}
{\( 400\,\mathrm{W} \)}{\( 80\,\mathrm{W} \)}{\( 40\,\mathrm{W} \)}{\( 790\,\mathrm{W} \)}{}{}}
\LANGcs{
\Question{
Výkon elektrického proudu ve spotřebiči je určen vztahem \( P=U_eI-R_i I^2 \), kde \( U_e \) a \( R_i \) charakterizují zdroj (\( U_e \) -elektromotorické napětí zdroje a \( R_i \) -vnitřní odpor zdroje). Jakého maximálního výkonu může dosáhnout proud ve spotřebiči, jestliže máme v obvodu zdroj o parametrech \( R_i=0{,}25\,\Omega \) a \( U_e=20\,\mathrm{V} \)?}
{\( 400\,\mathrm{W} \)}{\( 80\,\mathrm{W} \)}{\( 40\,\mathrm{W} \)}{\( 790\,\mathrm{W} \)}{}{}}
\LANGsk{
\Question{
Výkon elektrického prúdu v spotrebiči je určený vzťahom \( P=U_eI-R_i I^2 \), kde \( U_e \) a \( R_i \) charakterizujú zdroj (\( U_e \) -elektromotorické napätie zdroja a \( R_i \) -vnútorný odpor zdroja). Aký maximálny výkon môže dosiahnuť prúd v spotrebiči, ak máme v obvode zdroj s parametrami \( R_i=0{,}25\,\Omega \) a \( U_e=20\,\mathrm{V} \)?}
{\( 400\,\mathrm{W} \)}{\( 80\,\mathrm{W} \)}{\( 40\,\mathrm{W} \)}{\( 790\,\mathrm{W} \)}{}{}}
\LANGes{
\Question{
Considera un circuito eléctrico con una batería  \( U_e \) , una resitencia interna \( R_i \) y que conduce una corriente \( I \) hacia un receptor \( R \) (ve dibujo ). El receptor podría ser por ejemplo una luz eléctrica, un elemento de calefacción eléctrica, o posiblemente, un motor eléctrico. El objeto elemental del circuito es la transferencia de energía de la batería al receptor dónde se usa. (por ejemplo enciende una luz)
\[ \]
La fuerza \( P \) transferida al receptor se describe mediante la fórmula \( P=U_eI-R_i I^2 \). ¿Cuál es la fuerza máxima que se puede transferir al receptor si tenemos una fuente con  \( R_i=0.25\,\Omega \) y \( U_e=20\,\mathrm{V} \)?}
{\( 400\,\mathrm{W} \)}{\( 80\,\mathrm{W} \)}{\( 40\,\mathrm{W} \)}{\( 790\,\mathrm{W} \)}{}{}}
\End

\Data\ID{1103148605}
\Updated{2021-04-07 07:04:37}
\Level{C}
\SubArea{Quadratic Functions}
\IMG{1103148605.pdf}
\LANGen{
\Question{
Suppose, an object, that is in rest, starts to accelerate with the constant acceleration \( a \). The distance \( s \) travelled by the object in time \( t \) is given by the formula \( s=\frac12at^2 \). You can see the graph of the distance \( s \) on the time \( t \) dependency in the picture. Find the acceleration \( a \) of the object.}
{\( 8\frac{\mathrm{m}}{\mathrm{s}^2} \)}{\( 16\frac{\mathrm{m}}{\mathrm{s}^2} \)}{\( 4\frac{\mathrm{m}}{\mathrm{s}^2} \)}{\( 2\frac{\mathrm{m}}{\mathrm{s}^2} \)}{}{}}
\LANGpl{
\Question{
Zakładając, że przedmiot w stanie spoczynku zaczyna poruszać się ze stałym przyspieszeniem \( a \). Odległość \( s \) przebyta przez przedmiot w czasie \( t \) wyrażona jest wzorem \( s=\frac12at^2 \). Możesz zobaczyć wykres zależności odległości \( s \) od czasu \( t \)  na rysunku poniżej. Wyznacz przyspieszenie  \( a \) tego przedmiotu.}
{\( 8\frac{\mathrm{m}}{\mathrm{s}^2} \)}{\( 16\frac{\mathrm{m}}{\mathrm{s}^2} \)}{\( 4\frac{\mathrm{m}}{\mathrm{s}^2} \)}{\( 2\frac{\mathrm{m}}{\mathrm{s}^2} \)}{}{}}
\LANGcs{
\Question{
Jestliže těleso z klidu rovnoměrně zrychluje, je jeho dráha \( s \) funkcí času \( t \) s předpisem \( s=\frac12at^2 \), kde \( a \) je zrychlení tělesa.  Určete zrychlení tělesa, jehož graf dráhy (závislost dráhy na čase) je na obrázku.}
{\( 8\frac{\mathrm{m}}{\mathrm{s}^2} \)}{\( 16\frac{\mathrm{m}}{\mathrm{s}^2} \)}{\( 4\frac{\mathrm{m}}{\mathrm{s}^2} \)}{\( 2\frac{\mathrm{m}}{\mathrm{s}^2} \)}{}{}}
\LANGsk{
\Question{
Ak teleso z pokoja rovnomerne zrýchľuje, tak je jeho dráha \( s \) funkciou času \( t \) s predpisom \( s=\frac12at^2 \), kde \( a \) je zrýchlenie telesa.  Určte zrýchlenie telesa, ktorého graf dráhy (závislosti dráhy na čase) je na obrázku.}
{\( 8\frac{\mathrm{m}}{\mathrm{s}^2} \)}{\( 16\frac{\mathrm{m}}{\mathrm{s}^2} \)}{\( 4\frac{\mathrm{m}}{\mathrm{s}^2} \)}{\( 2\frac{\mathrm{m}}{\mathrm{s}^2} \)}{}{}}
\LANGes{
\Question{
Supongamos que un objeto en reposo empieza a acelerar con una aceleración constante \( a \). La distancia \( s \) recorrida por el objeto en tiempo \( t \) viene dada dada por la fórmula \( s=\frac12at^2 \). En el dibujo se puede ver la gráfica de la dependencia de la distancia  \( s \) respecto al tiempo \( t \). Halla la aceleración \( a \) del objeto.}
{\( 8\frac{\mathrm{m}}{\mathrm{s}^2} \)}{\( 16\frac{\mathrm{m}}{\mathrm{s}^2} \)}{\( 4\frac{\mathrm{m}}{\mathrm{s}^2} \)}{\( 2\frac{\mathrm{m}}{\mathrm{s}^2} \)}{}{}}
\End

\Data\ID{1103148606}
\Updated{2021-04-07 08:04:00}
\Level{C}
\SubArea{Quadratic Functions}
\LANGen{
\Question{
If an object moving with an initial velocity \( v_0 \) is slowing down at a constant deceleration \( a \), then the distance \( s \) travelled while decelerating is described by the formula \( s=v_0t-\frac12at^2 \), where \( t \) is time of decelerating. Choose the graph, which could represent the dependency of the distance \( s \) on the time \( t \).}
{}{}{}{}{}{}}
\LANGpl{
\Question{
Jeżeli przedmiot poruszający się z początkową prędkością  \( v_0 \) zmniejsza prędkość ze stałym zwalnianiem \( a \), wtedy odległość \( s \) przebyta przy zwalnianiu określona jest wzorem \( s=v_0t-\frac12at^2 \), gdzie \( t \) jest czasem zwalniania. Wybierz wykres, który mógłby przedstawiać zależność odległości  \( s \) od czasu \( t \).}
{}{}{}{}{}{}}
\LANGcs{
\Question{
Vyberte obrázek, který může vyjadřovat závislost dráhy na čase pro rovnoměrně zpomalený pohyb. Danou závislost popisuje pro tento typ pohybu rovnice \( s=v_0t-\frac12at^2 \), kde \( v_0 \) je počáteční rychlost a \( a \) je hodnota zpomalení pohybu.}
{}{}{}{}{}{}}
\LANGsk{
\Question{
Vyberte graf, ktorý môže vyjadrovať závislosť dráhy od času pre rovnomerne spomalený pohyb. Danú závislosť pre tento typ pohybu môžeme vyjadriť rovnicou \( s=v_0t-\frac12at^2 \), kde \( a \) je konštantné spomalenie pohybu a \( v_0 \) je počiatočná rýchlosť.}
{}{}{}{}{}{}}
\LANGes{
\Question{
Si un objeto moviéndose con velocidad inicial \( v_0 \) está decelerando con una deceleración constante \( a \), la distancia \( s \) recorrida mientras decelera se describe mediante la fórmula \( s=v_0t-\frac12at^2 \), dónde \( t \) es el tiempo que está decelerando. Elige la gráfica que representa la dependencia de la distancia \( s \) respecto al tiempo \( t \).}
{}{}{}{}{}{}}

\IMGanswer{1103148606a.pdf}
\IMGanswer{1103148606b.pdf}
\IMGanswer{1103148606c.pdf}
\IMGanswer{1103148606d.pdf}\End

\Data\ID{1103206101}
\Updated{2021-04-07 07:04:23}
\Level{C}
\SubArea{Quadratic Functions}
\IMG{1103206101_0.pdf}
\LANGen{
\Question{
There are graphs of three quadratic functions in the picture. Choose the formula which corresponds to all three functions graphed in the picture. Suppose that \( a\in\mathbb{R}^+ \).}
{\( y=a(x+2)^2-1 \)}{\( y=a(x-2)^2-1 \)}{\( y=a(x-1)^2+2 \)}{\( y=2(x-a)^2+1 \)}{}{}}
\LANGpl{
\Question{
Na rysunku przedstawiono wykresy trzech funkcji kwadratowych. Wybierz wzór, który odnosi się do wszystkich trzech wykresów. Załóżmy, że \( a\in\mathbb{R}^+ \).}
{\( y=a(x+2)^2-1 \)}{\( y=a(x-2)^2-1 \)}{\( y=a(x-1)^2+2 \)}{\( y=2(x-a)^2+1 \)}{}{}}
\LANGcs{
\Question{
Na obrázku jsou grafy tří kvadratických funkcí. Z následujících možností vyberte předpis, který odpovídá všem funkcím na obrázku. Předpokládejte, že  \( a\in\mathbb{R}^+ \).}
{\( y=a(x+2)^2-1 \)}{\( y=a(x-2)^2-1 \)}{\( y=a(x-1)^2+2 \)}{\( y=2(x-a)^2+1 \)}{}{}}
\LANGsk{
\Question{
Na obrázku sú grafy troch kvadratických funkcií. Vyberte predpis, ktorý zodpovedá všetkým trom funkciám, ktoré sú nakreslené na obrázku. Predpokladajme, že platí \( a\in\mathbb{R}^+ \).}
{\( y=a(x+2)^2-1 \)}{\( y=a(x-2)^2-1 \)}{\( y=a(x-1)^2+2 \)}{\( y=2(x-a)^2+1 \)}{}{}}
\LANGes{
\Question{
En el dibujo aparecen gráficas de tres funciones cuadráticas. Elije la fórmula que corresponde a las tres funciones. Supongamos que  \( a\in\mathbb{R}^+ \).}
{\( y=a(x+2)^2-1 \)}{\( y=a(x-2)^2-1 \)}{\( y=a(x-1)^2+2 \)}{\( y=2(x-a)^2+1 \)}{}{}}
\End

\Data\ID{1103206102}
\Updated{2021-04-07 07:04:45}
\Level{C}
\SubArea{Quadratic Functions}
\IMG{1103206102_0.pdf}
\LANGen{
\Question{
There are graphs of three quadratic functions in the picture. Choose the formula which corresponds to all three functions graphed in the picture.}
{\( y=-(x+a)^2+3 \), \( a\in(-\infty; 0] \)}{\( y=-(x+a)^2+3 \), \( a\in\mathbb{R}^+ \)}{\( y=-(x+3)^2+a \), \( a\in\mathbb{R}^+ \)}{\( y=-(x-3)^2+a \), \( a\in\mathbb{R}^+ \)}{}{}}
\LANGpl{
\Question{
Na rysunku przedstawione  są wykresy trzech funkcji kwadratowych. Wybierz wzór, który odnosi się do wszystkich trzech funkcji przedstawionych na rysunku.}
{\( y=-(x+a)^2+3 \), \( a\in(-\infty; 0\rangle \)}{\( y=-(x+a)^2+3 \), \( a\in\mathbb{R}^+ \)}{\( y=-(x+3)^2+a \), \( a\in\mathbb{R}^+ \)}{\( y=-(x-3)^2+a \), \( a\in\mathbb{R}^+ \)}{}{}}
\LANGcs{
\Question{
Na obrázku jsou grafy tří kvadratických funkcí. Z následujících možností vyberte předpis, který odpovídá všem funkcím na obrázku.}
{\( y=-(x+a)^2+3 \), \( a\in(-\infty; 0\rangle \)}{\( y=-(x+a)^2+3 \), \( a\in\mathbb{R}^+ \)}{\( y=-(x+3)^2+a \), \( a\in\mathbb{R}^+ \)}{\( y=-(x-3)^2+a \), \( a\in\mathbb{R}^+ \)}{}{}}
\LANGsk{
\Question{
Na obrázku sú grafy troch kvadratických funkcií. Vyberte predpis, ktorý zodpovedá všetkým trom funkciám, ktoré sú nakreslené na obrázku.}
{\( y=-(x+a)^2+3 \), \( a\in(-\infty; 0\rangle \)}{\( y=-(x+a)^2+3 \), \( a\in\mathbb{R}^+ \)}{\( y=-(x+3)^2+a \), \( a\in\mathbb{R}^+ \)}{\( y=-(x-3)^2+a \), \( a\in\mathbb{R}^+ \)}{}{}}
\LANGes{
\Question{
En el dibujo aparecen las gráficas de tres funciones cuadráticas. Elije la fórmula que corresponde a las tres funciones.}
{\( y=-(x+a)^2+3 \), \( a\in(-\infty; 0] \)}{\( y=-(x+a)^2+3 \), \( a\in\mathbb{R}^+ \)}{\( y=-(x+3)^2+a \), \( a\in\mathbb{R}^+ \)}{\( y=-(x-3)^2+a \), \( a\in\mathbb{R}^+ \)}{}{}}
\End

\Data\ID{2000004307}
\Updated{2022-02-20 08:02:01}
\Level{C}
\SubArea{Quadratic Functions}
\IMG{2000004401-figure7.pdf}
\LANGen{
\Question{
The function \(f\) is graphed in the picture. Which of the following is its equation?}
{\(f(x) = \left|x^2-4\right|\)}{\(f(x) = \left|x^2+4\right|\)}{\(f(x) = \left|x^2-2\right|\)}{\(f(x) = \left|x^2\right|+4\)}{}{}}
\LANGpl{
\Question{
Na rysunku przedstawiona jest funkcja \(f\). Wskaż wzór tej funkcji.}
{\(f(x) = \left|x^2-4\right|\)}{\(f(x) = \left|x^2+4\right|\)}{\(f(x) = \left|x^2-2\right|\)}{\(f(x) = \left|x^2\right|+4\)}{}{}}
\LANGcs{
\Question{
Funkce \(f\) je dána grafem. Z nabídnutých možností vyberte její předpis.}
{\(f(x) = \left|x^2-4\right|\)}{\(f(x) = \left|x^2+4\right|\)}{\(f(x) = \left|x^2-2\right|\)}{\(f(x) = \left|x^2\right|+4\)}{}{}}
\LANGsk{
\Question{
Funkcia \(f\) je daná grafom. Z ponúknutých možností vyberte jej prepis.}
{\(f(x) = \left|x^2-4\right|\)}{\(f(x) = \left|x^2+4\right|\)}{\(f(x) = \left|x^2-2\right|\)}{\(f(x) = \left|x^2\right|+4\)}{}{}}
\LANGes{
\Question{
La función \(f\) viene dada por la gráfica. Elige su ecuación.}
{\(f(x) = \left|x^2-4\right|\)}{\(f(x) = \left|x^2+4\right|\)}{\(f(x) = \left|x^2-2\right|\)}{\(f(x) = \left|x^2\right|+4\)}{}{}}
\End

\Data\ID{2000004308}
\Updated{2022-02-20 08:02:25}
\Level{C}
\SubArea{Quadratic Functions}
\IMG{2000004401-figure8.pdf}
\LANGen{
\Question{
The function \(f\) is graphed in the picture. Which of the following is its equation?}
{\( f(x) = \left|x^2-4x\right|\)}{\( f(x) = \left|x^2+4x\right|\)}{\( f(x) = \left|x^2\right|-4x\)}{\( f(x) = \left|x^2\right|+4x\)}{}{}}
\LANGpl{
\Question{
Dany jest wykres funkcji \(f\). Wskaż wzór tej funkcji.}
{\( f(x) = \left|x^2-4x\right|\)}{\( f(x) = \left|x^2+4x\right|\)}{\( f(x) = \left|x^2\right|-4x\)}{\( f(x) = \left|x^2\right|+4x\)}{}{}}
\LANGcs{
\Question{
Funkce \(f\) je dána grafem. Z nabídnutých možností vyberte její předpis.}
{\( f(x) = \left|x^2-4x\right|\)}{\( f(x) = \left|x^2+4x\right|\)}{\( f(x) = \left|x^2\right|-4x\)}{\( f(x) = \left|x^2\right|+4x\)}{}{}}
\LANGsk{
\Question{
Funkcia \(f\) je daná grafom. Z ponúknutých možností vyberte jej prepis.}
{\( f(x) = \left|x^2-4x\right|\)}{\( f(x) = \left|x^2+4x\right|\)}{\( f(x) = \left|x^2\right|-4x\)}{\( f(x) = \left|x^2\right|+4x\)}{}{}}
\LANGes{
\Question{
La función \(f\) viene dada por la gráfica. Elige su ecuación.}
{\( f(x) = \left|x^2-4x\right|\)}{\( f(x) = \left|x^2+4x\right|\)}{\( f(x) = \left|x^2\right|-4x\)}{\( f(x) = \left|x^2\right|+4x\)}{}{}}
\End

\Data\ID{2000004309}
\Updated{2022-02-20 08:02:52}
\Level{C}
\SubArea{Quadratic Functions}
\IMG{2000004401-figure9.pdf}
\LANGen{
\Question{
The function \(f\) is graphed in the picture. Which of the following is its equation?}
{\( f(x) =- \left|x^2-4\right|\)}{\( f(x) = -\left|x^2+4\right|\)}{\( f(x) = -\left|x^2\right|-4\)}{\( f(x) = -\left|x^2\right|+4\)}{}{}}
\LANGpl{
\Question{
Wykres przedstawia funkcję \(f\). Wskaż wzór tej funkcji.}
{\( f(x) =- \left|x^2-4\right|\)}{\( f(x) = -\left|x^2+4\right|\)}{\( f(x) = -\left|x^2\right|-4\)}{\( f(x) = -\left|x^2\right|+4\)}{}{}}
\LANGcs{
\Question{
Funkce \(f\) je dána grafem. Z nabídnutých možností vyberte její předpis.}
{\( f(x) =- \left|x^2-4\right|\)}{\( f(x) = -\left|x^2+4\right|\)}{\( f(x) = -\left|x^2\right|-4\)}{\( f(x) = -\left|x^2\right|+4\)}{}{}}
\LANGsk{
\Question{
Funkcia \(f\) je daná grafom. Z ponúknutých možností vyberte jej prepis.}
{\( f(x) =- \left|x^2-4\right|\)}{\( f(x) = -\left|x^2+4\right|\)}{\( f(x) = -\left|x^2\right|-4\)}{\( f(x) = -\left|x^2\right|+4\)}{}{}}
\LANGes{
\Question{
La función \(f\) viene dada por la gráfica. Elige su ecuación.}
{\( f(x) =- \left|x^2-4\right|\)}{\( f(x) = -\left|x^2+4\right|\)}{\( f(x) = -\left|x^2\right|-4\)}{\( f(x) = -\left|x^2\right|+4\)}{}{}}
\End

\Data\ID{2000004310}
\Updated{2022-02-20 08:02:14}
\Level{C}
\SubArea{Quadratic Functions}
\IMG{2000004401-figure10.pdf}
\LANGen{
\Question{
The function \(f\) is graphed in the picture. Which of the following is its equation?}
{\( f(x) = x^2-4|x|\)}{\( f(x) = \left|x^2-4x\right|\)}{\( f(x) = \left|x^2\right|-4x\)}{\( f(x) = \left|x^2\right|+4x\)}{}{}}
\LANGpl{
\Question{
Wykres przedstawia funkcję \(f\). Wskaż wzór tej funkcji.}
{\( f(x) = x^2-4|x|\)}{\( f(x) = \left|x^2-4x\right|\)}{\( f(x) = \left|x^2\right|-4x\)}{\( f(x) = \left|x^2\right|+4x\)}{}{}}
\LANGcs{
\Question{
Funkce \(f\) je dána grafem. Z nabídnutých možností vyberte její předpis.}
{\( f(x) = x^2-4|x|\)}{\( f(x) = \left|x^2-4x\right|\)}{\( f(x) = \left|x^2\right|-4x\)}{\( f(x) = \left|x^2\right|+4x\)}{}{}}
\LANGsk{
\Question{
Funkcia \(f\) je daná grafom. Z ponúknutých možností vyberte jej prepis.}
{\( f(x) = x^2-4|x|\)}{\( f(x) = \left|x^2-4x\right|\)}{\( f(x) = \left|x^2\right|-4x\)}{\( f(x) = \left|x^2\right|+4x\)}{}{}}
\LANGes{
\Question{
La función \(f\) viene dada por la gráfica. Elige su ecuación.}
{\( f(x) = x^2-4|x|\)}{\( f(x) = \left|x^2-4x\right|\)}{\( f(x) = \left|x^2\right|-4x\)}{\( f(x) = \left|x^2\right|+4x\)}{}{}}
\End

\Data\ID{2000004311}
\Updated{2022-02-15 07:02:41}
\Level{C}
\SubArea{Quadratic Functions}
\IMG{2000004401-figure11.pdf}
\LANGen{
\Question{
The function \(f\) is graphed in the picture. Which of the following is its equation?}
{\( f(x) = \left|x^2-1\right|-1\)}{\( f(x) = \left|x^2+1\right|-1\)}{\( f(x) = \left|x^2+1\right|+1\)}{\( f(x) = \left|x^2-1\right|+1\)}{}{}}
\LANGpl{
\Question{
Wykres przedstawia funkcję \(f\). Wskaż wzór tej funkcji.}
{\( f(x) = \left|x^2-1\right|-1\)}{\( f(x) = \left|x^2+1\right|-1\)}{\( f(x) = \left|x^2+1\right|+1\)}{\( f(x) = \left|x^2-1\right|+1\)}{}{}}
\LANGcs{
\Question{
Funkce \(f\) je dána grafem. Z nabídnutých možností vyberte její předpis.}
{\( f(x) = \left|x^2-1\right|-1\)}{\( f(x) = \left|x^2+1\right|-1\)}{\( f(x) = \left|x^2+1\right|+1\)}{\( f(x) = \left|x^2-1\right|+1\)}{}{}}
\LANGsk{
\Question{
Funkcia \(f\) je daná grafom. Z ponúknutých možností vyberte jej prepis.}
{\( f(x) = \left|x^2-1\right|-1\)}{\( f(x) = \left|x^2+1\right|-1\)}{\( f(x) = \left|x^2+1\right|+1\)}{\( f(x) = \left|x^2-1\right|+1\)}{}{}}
\LANGes{
\Question{
Tenemos dada la función \(f\)  por la gráfica. Elige su ecuación.}
{\( f(x) = \left|x^2-1\right|-1\)}{\( f(x) = \left|x^2+1\right|-1\)}{\( f(x) = \left|x^2+1\right|+1\)}{\( f(x) = \left|x^2-1\right|+1\)}{}{}}
\End

\Data\ID{2010012201}
\Updated{2022-11-02 09:11:23}
\Level{C}
\SubArea{Quadratic Functions}
\IMG{2010012201-figure0_0.pdf}
\LANGen{
\Question{
Function \( f \) is given completely by the graph. Identify which of the following statements is true.}
{\( f(x)=-\left|x^2-4\right|;\ x\in[ -3;3] \)}{\( f(x)=-\left|x^2+4\right|;\ x\in[ -3;3]\)}{\( f(x)=-\left|x^2\right|-4;\ x\in[ -3;3]\)}{\( f(x)=\left|-x^2\right|-4;\ x\in[ -3;3] \)}{}{}}
\LANGpl{
\Question{
Na wykresie przedstawiono funkcję \( f \). Wskaż prawdziwe stwierdzenie.}
{\( f(x)=-\left|x^2-4\right|;\ x\in\langle -3;3\rangle \)}{\( f(x)=-\left|x^2+4\right|;\ x\in\langle -3;3\rangle\)}{\( f(x)=-\left|x^2\right|-4;\ x\in\langle -3;3\rangle\)}{\( f(x)=\left|-x^2\right|-4;\ x\in\langle -3;3\rangle \)}{}{}}
\LANGcs{
\Question{
Funkce \( f \)  je dána grafem. Určete, který z následujících tvrzení je pravdivé.}
{\( f(x)=-\left|x^2-4\right|;\ x\in\langle -3;3\rangle \)}{\( f(x)=-\left|x^2+4\right|;\ x\in\langle -3;3\rangle\)}{\( f(x)=-\left|x^2\right|-4;\ x\in\langle -3;3\rangle\)}{\( f(x)=\left|-x^2\right|-4;\ x\in\langle -3;3\rangle \)}{}{}}
\LANGsk{
\Question{
Funkcia \( f \)  je daná grafom. Určte, ktoré z nasledujúcich tvrdení je pravdivé.}
{\( f(x)=-\left|x^2-4\right|;\ x\in\langle -3;3\rangle \)}{\( f(x)=-\left|x^2+4\right|;\ x\in\langle -3;3\rangle\)}{\( f(x)=-\left|x^2\right|-4;\ x\in\langle -3;3\rangle\)}{\( f(x)=\left|-x^2\right|-4;\ x\in\langle -3;3\rangle \)}{}{}}
\LANGes{
\Question{
La función \( f \) es dada por la siguiente gráfica. Identifica cuál de las siguientes afirmaciones es verdadera.}
{\( f(x)=-\left|x^2-4\right|;\ x\in[ -3;3] \)}{\( f(x)=-\left|x^2+4\right|;\ x\in[ -3;3]\)}{\( f(x)=-\left|x^2\right|-4;\ x\in[ -3;3]\)}{\( f(x)=\left|-x^2\right|-4;\ x\in[ -3;3] \)}{}{}}
\End

\Data\ID{2010012203}
\Updated{2022-11-02 09:11:23}
\Level{C}
\SubArea{Quadratic Functions}
\LANGen{
\Question{
Find all the values of the real parameter \( m \) such that \( f(x)=3(x+m)^2-2\) is an even function.}
{\( m=0 \)}{\( m>0 \)}{\( m=-2 \)}{\( m=2 \)}{}{}}
\LANGpl{
\Question{
Znajdź wszystkie wartości rzeczywistego parametru \( m \), dla których funkcja \( f(x)=3(x+m)^2-2\) jest parzysta.}
{\( m=0 \)}{\( m>0 \)}{\( m=-2 \)}{\( m=2 \)}{}{}}
\LANGcs{
\Question{
Najděte všechny hodnoty reálného parametru  \( m \), pro které je funkce  \( f(x)=3(x+m)^2-2\) sudá.}
{\( m=0 \)}{\( m>0 \)}{\( m=-2 \)}{\( m=2 \)}{}{}}
\LANGsk{
\Question{
Nájdite všetky hodnoty reálneho parametra  \( m \), pre ktoré je funkcia  \( f(x)=3(x+m)^2-2\) párna.}
{\( m=0 \)}{\( m>0 \)}{\( m=-2 \)}{\( m=2 \)}{}{}}
\LANGes{
\Question{
Halla todos los valores del parámetro real \( m \) tales que \( f(x)=3(x+m)^2-2\) sea una función par.}
{\( m=0 \)}{\( m>0 \)}{\( m=-2 \)}{\( m=2 \)}{}{}}
\End

\Data\ID{2010012204}
\Updated{2022-11-02 09:11:23}
\Level{C}
\SubArea{Quadratic Functions}
\LANGen{
\Question{
Find all the values of the real parameter \( m \) such that \( f(x)=-2(x-m)^2-5\) is decreasing on \( (0;\infty) \).}
{\( m\in(-\infty;0] \)}{\( m\in [0;\infty) \)}{\( m\in(-\infty;-5] \)}{\( m\in(-\infty;5] \)}{}{}}
\LANGpl{
\Question{
Znajdź wszystkie wartości rzeczywistego parametru \( m \), dla których funkcja\( f(x)=-2(x-m)^2-5\) jest malejąca w przedziale \( (0;\infty) \).}
{\( m\in(-\infty;0\rangle \)}{\( m\in \langle0;\infty) \)}{\( m\in(-\infty;-5\rangle \)}{\( m\in(-\infty;5\rangle \)}{}{}}
\LANGcs{
\Question{
Najděte všechny hodnoty reálného parametru  \( m \), pro které je funkce \( f(x)=-2(x-m)^2-5\) klesající na intervalu \( (0;\infty) \).}
{\( m\in(-\infty;0\rangle \)}{\( m\in \langle0;\infty) \)}{\( m\in(-\infty;-5\rangle \)}{\( m\in(-\infty;5\rangle \)}{}{}}
\LANGsk{
\Question{
Nájdite všetky hodnoty reálneho parametra  \( m \), pre ktorý je funkcia \( f(x)=-2(x-m)^2-5\) klesajúca na intervale \( (0;\infty) \).}
{\( m\in(-\infty;0\rangle \)}{\( m\in \langle0;\infty) \)}{\( m\in(-\infty;-5\rangle \)}{\( m\in(-\infty;5\rangle \)}{}{}}
\LANGes{
\Question{
Halla todos los valores del parámetro real \( m \) tales que \( f(x)=-2(x-m)^2-5\) sea decreciente en \( (0;\infty) \).}
{\( m\in(-\infty;0] \)}{\( m\in [0;\infty) \)}{\( m\in(-\infty;-5] \)}{\( m\in(-\infty;5] \)}{}{}}
\End

\Data\ID{2010012205}
\Updated{2022-11-02 09:11:23}
\Level{C}
\SubArea{Quadratic Functions}
\LANGen{
\Question{
Find all the values of the real parameter \( p \) such that \( f(x)=-2(x+3)^2+p \) is nonpositive for all \( x\in\mathbb{R} \).}
{\( p\in(-\infty;0] \)}{\( p\in [ 0;\infty) \)}{\( p\in(-\infty;-3] \)}{\( p\in [ -3;\infty) \)}{}{}}
\LANGpl{
\Question{
Znajdź wszystkie wartości rzeczywistego parametru \( p \), dla których funkcja\( f(x)=-2(x+3)^2+p \) przybiera wartości niedodatnie w całej swojej dziedzinie.}
{\( p\in(-\infty;0\rangle \)}{\( p\in \langle 0;\infty) \)}{\( p\in(-\infty;-3\rangle \)}{\( p\in \langle -3;\infty) \)}{}{}}
\LANGcs{
\Question{
Najděte všechny hodnoty reálného parametru  \( p \), pro které funkce  \( f(x)=-2(x+3)^2+p \) nabývá na celém svém definičním oboru nekladných hodnot.}
{\( p\in(-\infty;0\rangle \)}{\( p\in \langle 0;\infty) \)}{\( p\in(-\infty;-3\rangle \)}{\( p\in \langle -3;\infty) \)}{}{}}
\LANGsk{
\Question{
Nájdite všetky hodnoty reálneho parametra  \( p \), pre ktorý funkcia  \( f(x)=-2(x+3)^2+p \) nadobúda na celom svojom definičnom obore nekladné hodnoty.}
{\( p\in(-\infty;0\rangle \)}{\( p\in \langle 0;\infty) \)}{\( p\in(-\infty;-3\rangle \)}{\( p\in \langle -3;\infty) \)}{}{}}
\LANGes{
\Question{
Halla todos los valores del parámetro real \( p \) tales que \( f(x)=-2(x+3)^2+p \) sea no positiva para todo \( x\in\mathbb{R} \).}
{\( p\in(-\infty;0] \)}{\( p\in [ 0;\infty) \)}{\( p\in(-\infty;-3] \)}{\( p\in [ -3;\infty) \)}{}{}}
\End

\Data\ID{9000007101}
\Updated{2021-04-07 08:04:46}
\Level{C}
\SubArea{Quadratic Functions}
\IMG{000071_01-figure0.pdf}
\LANGen{
\Question{
Consider a family \(M\) 
of quadratic functions, as shown in the picture. Any quadratic function in this family
is given by the analytic expression 
\[ 
 y = ax^{2} + bx + c
\]
where \(a\),
\(b\) and
\(c\) are real constants and
\(a\not = 0\). For each function
the set \(K\) denotes the
set of \(x\)-intercepts. Complete the statement. „The analytic expressions for functions in the family
\(M\) differ
only in ....”}
{the coefficient \(a\)}{the coefficient \(b\)}{the coefficient \(c\)}{the set \(K\)}{}{}}
\LANGpl{
\Question{
Rozważ rodzinę  funkcji kwadratowych \(M\) pokazaną na rysunku poniżej. Każda funkcja kwadratowa w tej rodzinie jest wyrażona wzorem  
\[ 
 y = ax^{2} + bx + c
\]
gdzie \(a\),
\(b\) i 
\(c\) są rzeczywistymi stałymi i 
\(a\not = 0\) zaś
\(K\) określa zbiór rozwiązań równania
\(ax^{2} + bx + c=0\). Dokończ zdanie. „Wzory funkcji w rodzinie 
\(M\) różnią się jedynie ....”}
{współczynnikiem \(a\)}{współczynnikiem  \(b\)}{współczynnikiem  \(c\)}{zbiorem \(K\)}{}{}}
\LANGcs{
\Question{
Předpokládejte kvadratické funkce, které jsou dány předpisem ve tvaru
\(y = ax^{2} + bx + c\), kde
\(a,\ b,\ c\) jsou reálné
koeficienty, přičemž \(a\not = 0\) 
a \(K\) je množina
kořenů rovnice \(ax^{2} + bx + c = 0\). Doplňte tvrzení: „Předpis
funkcí, které jsou znázorněny grafem, se liší pouze ....”}
{hodnotou koeficientu \(a\)}{hodnotou koeficientu \(b\)}{hodnotou koeficientu \(c\)}{v množině kořenů \(K\)}{}{}}
\LANGsk{
\Question{
Predpokladajme množinu \(M\) kvadratických funkcií, ktoré sú znázornené na obrázku. Každá kvadratická funkcia tejto množiny má predpis v tvare
\(y = ax^{2} + bx + c\), kde \(a,\ b,\ c\) sú reálne koeficienty, pričom \(a\not = 0\) 
a \(K\) je množina koreňov rovnice \(ax^{2} + bx + c = 0\).
Doplňte tvrdenie: 
„Predpisy funkcií množiny \(M\) sa líšia len ....”}
{hodnotou koeficientu \(a\)}{hodnotou koeficientu \(b\)}{hodnotou koeficientu \(c\)}{množinou koreňov \(K\)}{}{}}
\LANGes{
\Question{
Considera una familia \(M\) 
de funciones cuadráticas, como aparecen en el dibujo. Cualquier función cuadrática en esta familia se puede expresar analíticamente mediante
\[ 
 y = ax^{2} + bx + c
\]
dónde \(a\),
\(b\) y
\(c\) son constantes reales y
\(a\not = 0\). Para cada función el conjunto \(K\) denota el conjunto de \(x\)-intesecantes.. Completa la frase: Las expresiónes analíticas para las funciones de la familia 
\(M\) difieren solamente en ...}
{el coeficiente \(a\)}{el coeficiente \(b\)}{el coeficiente \(c\)}{el conjunto \(K\)}{}{}}
\End

\Data\ID{9000007102}
\Updated{2021-04-07 08:04:09}
\Level{C}
\SubArea{Quadratic Functions}
\IMG{000071_02-figure0.pdf}
\LANGen{
\Question{
Consider a family \(M\) 
of quadratic functions, as shown in the picture. Any quadratic function in this family
is given by the analytic expression 
\[ 
 y = ax^{2} + bx + c
\]
where \(a\),
\(b\) and
\(c\) are real constants and
\(a\not = 0\). For each function
the set \(K\) denotes the
set of \(x\)-intercepts. Complete the statement. „The analytic expressions for functions in the family
\(M\) differ
only in ....”}
{the coefficient \(c\)}{the coefficient \(a\)}{the coefficient \(b\)}{the solution set \(K\)}{}{}}
\LANGpl{
\Question{
Rozważ rodzinę  \(M\) 
funkcji kwadratowych, jak przedstawiono na rysunku. Dowolna funkcja kwadratowa w tej rodzinie jest wyrażona wzorem 
\[ 
 y = ax^{2} + bx + c
\]
gdzie \(a\),
\(b\) i
\(c\) rzeczywistymi stałymi a
\(a\not = 0\). Dla każdej funkcji
zbiór \(K\) oznacza zbiór punktów przecięcia z osia współrzędnych \(x\).  Dokończ zdanie: „Wzory dla funkcji w rodzinie 
\(M\) różnią się tylko ....”}
{współczynnikiem  \(c\)}{współczynnikiem \(a\)}{współczynnikiem \(b\)}{zbiorem rozwiązań \(K\)}{}{}}
\LANGcs{
\Question{
Předpokládejte kvadratické funkce, které jsou dány předpisem ve tvaru
\(y = ax^{2} + bx + c\), kde
\(a,\ b,\ c\) jsou reálné
koeficienty, přičemž \(a\not = 0\) 
a \(K\) je množina
kořenů rovnice \(ax^{2} + bx + c = 0\). Doplňte tvrzení:
„Předpis
funkcí, které jsou znázorněny grafem, se liší pouze ....”}
{hodnotou koeficientu \(c\)}{hodnotou koeficientu \(a\)}{hodnotou koeficientu \(b\)}{v množině kořenů \(K\)}{}{}}
\LANGsk{
\Question{
Predpokladajme množinu \(M\) kvadratických funkcií, ktoré sú znázornené na obrázku. Každá kvadratická funkcia tejto množiny má predpis v tvare
\(y = ax^{2} + bx + c\), kde \(a,\ b,\ c\) sú reálne koeficienty, pričom \(a\not = 0\) 
a \(K\) je množina koreňov rovnice \(ax^{2} + bx + c = 0\).
Doplňte tvrdenie: 
„Predpisy funkcií množiny \(M\) sa líšia len ....”}
{hodnotou koeficientu \(c\)}{hodnotou koeficientu \(a\)}{hodnotou koeficientu \(b\)}{množinou koreňov \(K\)}{}{}}
\LANGes{
\Question{
Considera una familia \(M\) 
de funciones cuadráticas, como aparecen en el dibujo. Cualquier función cuadrática en esta familia se puede expresar analíticamente mediante
\[ 
 y = ax^{2} + bx + c
\]
dónde \(a\),
\(b\) y
\(c\) son constantes reales y
\(a\not = 0\). Para cada función el conjunto \(K\) denota el conjunto de \(x\)-intesecantes.. Completa la frase: Las expresiónes analíticas para las funciones de la familia 
\(M\) difieren solamente en ...}
{el coeficiente \(c\)}{el coeficiente \(a\)}{el coeficiente \(b\)}{el conjunto de soluciones \(K\)}{}{}}
\End

\Data\ID{9000007103}
\Updated{2021-04-07 08:04:25}
\Level{C}
\SubArea{Quadratic Functions}
\IMG{000071_03-figure0.pdf}
\LANGen{
\Question{
Consider a family \(M\) 
of quadratic functions, as shown in the picture. Any quadratic function in this family
is given by the analytic expression 
\[ 
 y = ax^{2} + bx + c
\]
where \(a\),
\(b\) and
\(c\) are real constants and
\(a\not = 0\). For each function
the set \(K\) denotes the
set of \(x\)-intercepts. Complete the statement. „The analytic expressions for functions in the family
\(M\) share
only ....”}
{the value of the coefficient \(a\)}{the value of the coefficient \(b\)}{the value of the coefficient \(c\)}{the solution set \(K\)}{}{}}
\LANGpl{
\Question{
Rozważ rodzinę \(M\) 
funkcji kwadratowych, jak przedstawiono na rysunku. Dowolna funkcja kwadratowa w tej rodzinie wyrażona jest wzorem 
\[ 
 y = ax^{2} + bx + c
\]
gdzie \(a\),
\(b\) i
\(c\) są rzeczywistymi stałymi, a 
\(a\not = 0\). Dla każdej funkcji  \(K\) oznacza zbiór punktów przecięcia z osia współrzędnych  \(x\). Dokończ zdanie: „Wzory dla funkcji w rodzinie 
\(M\) dzielą 
jedynie ....”}
{wartość współczynnika \(a\)}{wartość współczynnika \(b\)}{wartość współczynnika \(c\)}{zbiór rozwiązań \(K\)}{}{}}
\LANGcs{
\Question{
Předpokládejte kvadratické funkce, které jsou dány předpisem ve tvaru
\(y = ax^{2} + bx + c\), kde
\(a,\ b,\ c\) jsou reálné
koeficienty, přičemž \(a\not = 0\) 
a \(K\) je množina
kořenů rovnice \(ax^{2} + bx + c = 0\).
Doplňte tvrzení: „Předpis
všech funkcí, které jsou znázorněny grafem, se shoduje pouze ....”}
{hodnotou koeficientu \(a\)}{hodnotou koeficientu \(b\)}{hodnotou koeficientu \(c\)}{v množině kořenů \(K\)}{}{}}
\LANGsk{
\Question{
Predpokladajme množinu \(M\) kvadratických funkcií, ktoré sú znázornené na obrázku. Každá kvadratická funkcia tejto množiny má predpis v tvare
\(y = ax^{2} + bx + c\), kde \(a,\ b,\ c\) sú reálne koeficienty, pričom \(a\not = 0\) 
a \(K\) je množina koreňov rovnice \(ax^{2} + bx + c = 0\).
Doplňte tvrdenie: 
„Predpisy funkcií množiny \(M\) sa zhodujú len ....”}
{hodnotou koeficientu \(a\)}{hodnotou koeficientu \(b\)}{hodnotou koeficientu \(c\)}{množinou koreňov \(K\)}{}{}}
\LANGes{
\Question{
Considera una familia \(M\) 
de funciones cuadráticas, como aparecen en el dibujo. Cualquier función cuadrática de esta familia se puede expresar analíticamente mediante
\[ 
 y = ax^{2} + bx + c
\]
dónde \(a\),
\(b\) y
\(c\) son constantes reales y
\(a\not = 0\). Para cada función el conjunto \(K\) denota el conjunto de \(x\)-intesecantes. Completa la frase: Las expresiónes analíticas para las funciones de la familia 
\(M\) tienen en común solamente ...}
{el valor del coeficiente \(a\)}{el valor del coeficiente \(b\)}{el valor del coeficiente \(c\)}{el conjunto de soluciones \(K\)}{}{}}
\End

\Data\ID{9000007104}
\Updated{2021-04-07 08:04:04}
\Level{C}
\SubArea{Quadratic Functions}
\IMG{000071_04-figure0.pdf}
\LANGen{
\Question{
Consider a family \(M\) 
of quadratic functions, as shown in the picture. Any quadratic function in this family
is given by the analytic expression 
\[ 
 y = ax^{2} + bx + c
\]
where \(a\),
\(b\) and
\(c\) are real constants and
\(a\not = 0\). For each function
the set \(K\) denotes the
set of \(x\)-intercepts. Complete the statement. „The analytic expressions for functions in the family
\(M\) share
only ....”}
{the value of the coefficient \(b\)}{the value of the coefficient \(a\)}{the value of the coefficient \(c\)}{the solution set \(K\)}{}{}}
\LANGpl{
\Question{
Rozważ rodzinę \(M\) 
funkcji kwadratowych, jak przedstawiono na rysunku. Dowolna funkcja kwadratowa w tej rodzinie wyrażona jest wzorem 
\[ 
 y = ax^{2} + bx + c
\]
gdzie  \(a\),
\(b\) i
\(c\) są rzeczywistymi stałymi, a
\(a\not = 0\). Dla każdej funkcji zbiór \(K\) oznacza zbiór punktów przecięcia z osią współrzędnych \(x\).  Dokończ zdanie: „Wzory dla funkcji w tej rodzinie 
\(M\) dzielą
tylko....”}
{wartość współczynnika \(b\)}{wartość współczynnika \(a\)}{wartość współczynnika \(c\)}{zbiór rozwiązań \(K\)}{}{}}
\LANGcs{
\Question{
Předpokládejte kvadratické funkce, které jsou dány předpisem ve tvaru
\(y = ax^{2} + bx + c\), kde
\(a,\ b,\ c\) jsou reálné
koeficienty, přičemž \(a\not = 0\) 
a \(K\) je množina
kořenů rovnice \(ax^{2} + bx + c = 0\).
Doplňte tvrzení: „Předpis
všech funkcí, které jsou znázorněny grafem, se shoduje pouze ....”}
{hodnotou koeficientu \(b\)}{hodnotou koeficientu \(a\)}{hodnotou koeficientu \(c\)}{v množině kořenů \(K\)}{}{}}
\LANGsk{
\Question{
Predpokladajme množinu \(M\) kvadratických funkcií, ktoré sú znázornené na obrázku. Každá kvadratická funkcia tejto množiny má predpis v tvare
\(y = ax^{2} + bx + c\), kde \(a,\ b,\ c\) sú reálne koeficienty, pričom \(a\not = 0\) 
a \(K\) je množina koreňov rovnice \(ax^{2} + bx + c = 0\).
Doplňte tvrdenie: 
„Predpisy funkcií množiny \(M\) sa zhodujú len ....”}
{hodnotou koeficientu \(b\)}{hodnotou koeficientu \(a\)}{hodnotou koeficientu \(c\)}{množinou koreňov \(K\)}{}{}}
\LANGes{
\Question{
Considera una familia \(M\) 
de funciones cuadráticas, como aparecen en el dibujo. Cualquier función cuadrática de esta familia se puede expresar analíticamente mediante
\[ 
 y = ax^{2} + bx + c
\]
dónde \(a\),
\(b\) y
\(c\) son constantes reales y
\(a\not = 0\). Para cada función el conjunto \(K\) denota el conjunto de \(x\)-intesecantes. Completa la frase: Las expresiónes analíticas para las funciones de la familia 
\(M\) tienen en común solamente ...}
{el valor del coeficiente \(b\)}{el valor del coeficiente \(a\)}{el valor del coeficiente \(c\)}{el conjunto de soluciones \(K\)}{}{}}
\End

\Data\ID{9000007105}
\Updated{2021-04-07 08:04:22}
\Level{C}
\SubArea{Quadratic Functions}
\IMG{000071_05-figure0.pdf}
\LANGen{
\Question{
Consider a family \(M\) 
of quadratic functions, as shown in the picture. Any quadratic function in this family
is given by the analytic expression 
\[ 
 y = ax^{2} + bx + c
\]
where \(a\),
\(b\) and
\(c\) are real constants and
\(a\not = 0\). For each function
the set \(K\) denotes the
set of \(x\)-intercepts. Complete the statement. „The analytic expressions for functions in the family
\(M\) share
only ....”}
{the solution set \(K\)}{the value of the coefficient \(a\)}{the value of the coefficient \(b\)}{the value of the coefficient \(c\)}{}{}}
\LANGpl{
\Question{
Rozważ rodzinę \(M\) 
funkcji kwadratowych, jak przedstawiono na rysunku. Dowolna funkcja kwadratowa w tej rodzinie wyrażona jest wzorem 
\[ 
 y = ax^{2} + bx + c
\]
gdzie \(a\),
\(b\) i
\(c\) są rzeczywistymi stałymi a 
\(a\not = 0\). Dla każdej funkcji zbiór \(K\) oznacza zbiór punktów przecięcia z osią współrzędnych \(x\). Dokończ zdanie: „Wzory dla funkcji w tej rodzinie 
\(M\) dzielą 
jedynie....”}
{zbiór rozwiązań \(K\)}{wartość współczynnika \(a\)}{wartość współczynnika \(b\)}{wartość współczynnika \(c\)}{}{}}
\LANGcs{
\Question{
Předpokládejte kvadratické funkce, které jsou dány předpisem ve tvaru
\(y = ax^{2} + bx + c\), kde
\(a,\ b,\ c\) jsou reálné
koeficienty, přičemž \(a\not = 0\) 
a \(K\) je množina
kořenů rovnice \(ax^{2} + bx + c = 0\).
Doplňte tvrzení: „Předpis všech
funkcí, které jsou znázorněny grafem, se shoduje pouze ....”}
{v množině kořenů \(K\)}{hodnotou koeficientu \(a\)}{hodnotou koeficientu \(b\)}{hodnotou koeficientu \(c\)}{}{}}
\LANGsk{
\Question{
Predpokladajme množinu \(M\) kvadratických funkcií, ktoré sú znázornené na obrázku. Každá kvadratická funkcia tejto množiny má predpis v tvare
\(y = ax^{2} + bx + c\), kde \(a,\ b,\ c\) sú reálne koeficienty, pričom \(a\not = 0\) 
a \(K\) je množina koreňov rovnice \(ax^{2} + bx + c = 0\).
Doplňte tvrdenie: 
„Predpisy funkcií množiny \(M\) sa zhodujú .... ”}
{množinou koreňov \(K\)}{hodnotou koeficientu \(a\)}{hodnotou koeficientu \(b\)}{hodnotou koeficientu \(c\)}{}{}}
\LANGes{
\Question{
Considera una familia \(M\) 
de funciones cuadráticas, como aparecen en el dibujo. Cualquier función cuadrática de esta familia se puede expresar analíticamente mediante
\[ 
 y = ax^{2} + bx + c
\]
dónde \(a\),
\(b\) y
\(c\) son constantes reales y
\(a\not = 0\). Para cada función el conjunto \(K\) denota el conjunto de \(x\)-intesecantes. Completa la frase: Las expresiónes analíticas para las funciones de la familia 
\(M\) tienen en común solamente ...}
{el conjunto de soluciones \(K\)}{el valor del coeficiente \(a\)}{el valor del coeficiente \(b\)}{el valor del coeficiente \(c\)}{}{}}
\End

\Data\ID{9000033707}
\Updated{2020-11-20 12:11:18}
\Level{C}
\SubArea{Quadratic Functions}
\IMG{000337_07-figure0.pdf}
\LANGen{
\Question{
Identify the inequality with a solution graphed in the picture.}
{\(|x(3 - x)| > 3 - x\)}{\(|x(x - 3)| < x - 3\)}{\(|3x - x^{2}| > x - 3\)}{\(|x^{2} - 3x| < 3 - x\)}{\(x^{2} - 3|x| > 3 - x\)}{\(x^{2} - 3|x| < x - 3\)}}
\LANGpl{
\Question{
Zidentyfikuj nierówność z rozwiązaniem przedstawionym na rysunku.}
{\(|x(3 - x)| > 3 - x\)}{\(|x(x - 3)| < x - 3\)}{\(|3x - x^{2}| > x - 3\)}{\(|x^{2} - 3x| < 3 - x\)}{\(x^{2} - 3|x| > 3 - x\)}{\(x^{2} - 3|x| < x - 3\)}}
\LANGcs{
\Question{
Vyberte tu z nerovnic, jejíž řešení je graficky ilustrováno na obrázku.}
{\(|x(3 - x)| > 3 - x\)}{\(|x(x - 3)| < x - 3\)}{\(|3x - x^{2}| > x - 3\)}{\(|x^{2} - 3x| < 3 - x\)}{\(x^{2} - 3|x| > 3 - x\)}{\(x^{2} - 3|x| < x - 3\)}}
\LANGsk{
\Question{
Vyberte nerovnicu, ktorej riešenie je graficky znázornené na obrázku.}
{\(|x(3 - x)| > 3 - x\)}{\(|x(x - 3)| < x - 3\)}{\(|3x - x^{2}| > x - 3\)}{\(|x^{2} - 3x| < 3 - x\)}{\(x^{2} - 3|x| > 3 - x\)}{\(x^{2} - 3|x| < x - 3\)}}
\LANGes{
\Question{
En el dibujo está representada la solución de una de las inecuaciones de abajo, ¿Cuál es?}
{\(|x(3 - x)| > 3 - x\)}{\(|x(x - 3)| < x - 3\)}{\(|3x - x^{2}| > x - 3\)}{\(|x^{2} - 3x| < 3 - x\)}{\(x^{2} - 3|x| > 3 - x\)}{\(x^{2} - 3|x| < x - 3\)}}
\End
